
\documentclass[12pt,a4paper]{scrartcl}
\usepackage{lmodern}
\usepackage{amssymb,amsmath}

\usepackage[T1]{fontenc}
\usepackage[utf8]{inputenc}
\usepackage[magyar]{babel}

\usepackage{upquote}
\usepackage[]{microtype}
\usepackage[dvipsnames,svgnames*,x11names*]{xcolor}
\usepackage[unicode=true]{hyperref}
\hypersetup{
            pdftitle={Analízis III},
            pdfauthor={Izsák Ferenc; Tarcsay Zsigmond; Tüzes Dániel},
            colorlinks=true,
            linkcolor=Maroon,
            citecolor=Blue,
            urlcolor=Blue,
            breaklinks=true}

\usepackage{graphicx,grffile}
\makeatletter
\def\fps@figure{htbp}
\renewcommand\paragraph{\@startsection{paragraph}{4}{\z@}%
                                    {3.25ex \@plus1ex \@minus.2ex}%
                                    {0.5em} % meg kellett növelni, mert amúgy -0.2em volt itt, és kicsinek tűnt
                                    {\normalfont\normalsize\bfseries}}
\makeatother

\usepackage[normalem]{ulem}

\providecommand{\tightlist}{%
  \setlength{\itemsep}{0pt}\setlength{\parskip}{0pt}}
\setcounter{secnumdepth}{4}


\usepackage{environ} % azért kell, mert egyes környezeteket el akarok rejteni
\usepackage{amsthm}
\usepackage{newunicodechar} % unicode-ban egyszerűbb tárolni, viszont a latex alapból nem ismeri ezeket
\newunicodechar{■}{\qed}
\newunicodechar{∎}{\qed}
\newunicodechar{↣}{\rightarrowtail}

\renewcommand{\contentsname}{Tartalomjegyzék}


\definecolor{sotetvoros}{HTML}{900000} % a színekre így hivatkozok
\definecolor{sotetzold}{HTML}{009900} % a színekre így hivatkozok
\usepackage{tcolorbox} % a def, tétel, bizonyítás, állítás, megjegyzés körüli színes kerethez.
\tcbuselibrary{most} % a def, tétel, bizonyítás, állítás, megjegyzés körüli színes kerethez.

%
% Saját környezetek definiálása, hogy a definíciók, tételek, bizonyítások és megjegyzések külön,
% de egységesen formázhatóak legyenek. A NewEnviron azért kell, hogy egyes környezetek
% elrejthetőek legyenek. Pl. azon extra definíciók, tételek, bizonyítások, amelyek a régi
% jegyzetben benne voltak, de most már nem szükségesek. Elrejthető továbbá az "ajánló" rovat,
% amely nem a tananyaghoz kapcsolódó vicces, közéleti, de a természettudományokhoz, vagy az
% emberi kultúrkincshez valahogy másképp hozzájárul.
%

\newenvironment{definicio}{}{}
\tcolorboxenvironment{definicio}{blanker,breakable,left=5mm, right=5mm,
before skip=10pt,after skip=10pt,borderline west={1.5mm}{0pt}{black},borderline east={1.5mm}{0pt}{black}}

\NewEnviron{def_extra}{
\begin{definicio} % ennek és az alábbi sornak a kikommentelésel eltűnik az apró oldalsó színes vonal a régi, már nem szükséges tartalomról
%\BODY
\end{definicio} % ennek és az alábbi sornak a kikommentelésel eltűnik az apró oldalsó színes vonal a régi, már nem szükséges tartalomról
}


\newenvironment{tetel}{}{}
\tcolorboxenvironment{tetel}{blanker,breakable,left=5mm,right=5mm,
before skip=10pt,after skip=10pt,borderline west={1.5mm}{0pt}{sotetvoros},borderline east={1.5mm}{0pt}{sotetvoros}}

\NewEnviron{tetel_extra}{
\begin{tetel} % ennek és az alábbi sornak a kikommentelésel eltűnik az apró oldalsó színes vonal a régi, már nem szükséges tartalomról
%\BODY
\end{tetel} % ennek és az alábbi sornak a kikommentelésel eltűnik az apró oldalsó színes vonal a régi, már nem szükséges tartalomról
}


\newenvironment{bizonyitas}{}{}
\tcolorboxenvironment{bizonyitas}{blanker,breakable,left=5mm,right=5mm,
before skip=10pt,after skip=10pt,borderline west={1.5mm}{0pt}{sotetvoros},borderline east={1.5mm}{0pt}{sotetvoros}}

\NewEnviron{biz_extra}{
\begin{bizonyitas} % ennek és az alábbi sornak a kikommentelésel eltűnik az apró oldalsó színes vonal a régi, már nem szükséges tartalomról
%\BODY
\end{bizonyitas} % ennek és az alábbi sornak a kikommentelésel eltűnik az apró oldalsó színes vonal a régi, már nem szükséges tartalomról
}

\newenvironment{allitas}{}{}
\tcolorboxenvironment{allitas}{blanker,breakable,left=5mm,right=5mm,
before skip=10pt,after skip=10pt,borderline west={0.75mm}{0pt}{sotetvoros},borderline east={0.75mm}{0pt}{sotetvoros}}

\NewEnviron{all_extra}{
\begin{allitas} % ennek és az alábbi sornak a kikommentelésel eltűnik az apró oldalsó színes vonal a régi, már nem szükséges tartalomról
%\BODY
\end{allitas} % ennek és az alábbi sornak a kikommentelésel eltűnik az apró oldalsó színes vonal a régi, már nem szükséges tartalomról
}

\newenvironment{megjegyzes}{}{}
\tcolorboxenvironment{megjegyzes}{blanker,breakable,left=5mm,right=5mm,
before skip=10pt,after skip=10pt,borderline west={0.75mm}{0pt}{lightgray},borderline east={0.75mm}{0pt}{lightgray}}

\NewEnviron{megj_extra}{
\begin{megjegyzes} % ennek és az alábbi sornak a kikommentelésel eltűnik az apró oldalsó színes vonal a régi, már nem szükséges tartalomról
%\BODY
\end{megjegyzes} % ennek és az alábbi sornak a kikommentelésel eltűnik az apró oldalsó színes vonal a régi, már nem szükséges tartalomról
}


\newenvironment{pelda}{}{}
\tcolorboxenvironment{pelda}{blanker,breakable,left=5mm,right=5mm,
before skip=10pt,after skip=10pt,borderline west={1.5mm}{	0pt}{sotetzold},borderline east={1.5mm}{0pt}{sotetzold}}


\newenvironment{ajanlofig}{\begin{figure}\begin{center}}{
\end{center}\end{figure}}

\NewEnviron{ajanlo}{
\begin{samepage} % az alábbi sorok kikommentelésével eltűnnek az ajánlók
\footnotesize
\BODY
\normalsize
\end{samepage} % eddig kell kikommentelni, hogy az ajánlók eltűnjenek
}


\title{Analízis III}
\providecommand{\subtitle}[1]{}
\subtitle{Előadásjegyzet fizikusoknak matematikusoktól}
\author{Izsák Ferenc \and Tarcsay Zsigmond \and Tüzes Dániel}
\providecommand{\institute}[1]{}
\institute{ELTE}
\date{}

\begin{document}
\maketitle


\setcounter{tocdepth}{6}
\tableofcontents

~

Előadó e-mail címe: simonl a ludens.elte.hu-nál

Ez a jegyzet \textbf{nem} szakirodalom s nem garantált, hogy az órai
anyagot teljesen lefedi, az előadásokra bejárni ajánlott.

Ha a jegyzetben helyesírási, tartalmi vagy formai hibát találsz, kérlek
jelezd az előadónak!

\hypertarget{differencialegyenletek}{%
\section{Differenciálegyenletek}\label{differencialegyenletek}}

Mi a differenciálegyenlet?

Pl

\begin{enumerate}
\def\labelenumi{\arabic{enumi}.}
\tightlist
\item
  \(\ddot x\left( t \right) = - {\omega ^2}x\left( t \right)\)
\item
  \(\ddot x\left( t \right) = F\left( t \right)/m\)
\item
  \(\partial_{t}u = \Delta u\)
\item
  \(\overset{}{x}\left( t \right) = x\left( {t - 1} \right)\)
\end{enumerate}

Ezeket lehet rendszerezni: ODE (ordenary differential equation, azaz
közönséges differenciál-egyenlet, 1-es és 2-es), PDE (partial
differential equation, 3-as), FDE (functional differential equation,
4-es).

Most az ODE-val foglalkozunk. Mi a közönséges differenciál-egyenlet?

\begin{definicio}

Definíció:\\
Legyen \(\left. F:{\mathbb{R}}^{n + 2}\rightarrow{\mathbb{R}} \right.\),
\(n\)-edrendű közönséges differenciálegyenlet: \(\forall t\)-re
\(0 = F\left( {t,x\left( t \right),\dot x\left( t \right),\ddot x\left( t \right),...,{x^{\left( n \right)}}\left( t \right)} \right)\).

\end{definicio}

\begin{megjegyzes}

Megjegyzés:\\
Egy ilyen \(n\)-edrendű egyenlet átírató elsőrendű rendszerré. Pl:
\(\ddot x\left( t \right) = - {\omega ^2}x\left( t \right)\) egyenletet
átírjuk: \(y_{1}\left( t \right) = x\left( t \right)\),
\(y_{2}\left( t \right) = \overset{}{x}\left( t \right)\). Ekkor
\(y\)-ra az alábbi elsőrendű, kétismeretlenes rendszer áll fenn:
\[\begin{array}{l}
{y_{1}\left( t \right) = y_{2}\left( t \right)} \\
{y_{2}\left( t \right) = - \omega^{2} \cdot y_{1}{\left( t \right).}} \\
\end{array}\] Általánosan \(n\)-ed rendűnél: \(y_{1} = x\),
\(y_{2} = \overset{}{x}\), \ldots{} , \(y_{n} = x^{({n - 1})}\). Ekkor
\(\left( {y_{1},...,y_{n}} \right)\)-re elsőrendű rendszert kapunk.

\end{megjegyzes}

\begin{definicio}

Definíció:\\
Legyen
\(\left. f:{\mathbb{R}} \times {\mathbb{R}}^{n}\rightarrow{\mathbb{R}}^{n} \right.\),
\(\overset{}{x}\left( t \right) = f\left( {t,x\left( t \right)} \right)\)
elsőrendű (explicit) közönséges differenciálegyenlet-rendszer.
Ismeretlen az
\(\left. x:{\mathbb{R}}\rightarrow{\mathbb{R}}^{n} \right.\) függvény.
Koordinátánként kiírva:\\
\(\begin{array}{l} {{\overset{}{x}}_{1}\left( t \right) = f_{1}\left( {t,x_{1}\left( t \right),x_{2}\left( t \right),...,x_{n}\left( t \right)} \right)} \\ {{\overset{}{x}}_{2}\left( t \right) = f_{2}\left( {t,x_{1}\left( t \right),x_{2}\left( t \right),...,x_{n}\left( t \right)} \right)} \\  \vdots \\ {{\overset{}{x}}_{n}\left( t \right) = f_{n}\left( {t,x_{1}\left( t \right),x_{2}\left( t \right),...,x_{n}\left( t \right)} \right)} \\ \end{array}\)

\end{definicio}

Mivel foglalkozik a közönséges differenciálelmélet?

\begin{enumerate}
\def\labelenumi{\arabic{enumi}.}
\tightlist
\item
  Mi a megoldás? Azaz számítsuk ki a megoldást. (Ezt már tanultuk.)
  Vannak:

  \begin{enumerate}
  \def\labelenumii{\alph{enumii}.}
  \tightlist
  \item
    képlettel megoldhatók
  \item
    képlettel nem megoldhatók (de numerikusan közelíthetők)
  \end{enumerate}
\item
  Megoldás létezésének, egyértelműségének keresése, függése a
  paraméterektől
\item
  Milyen a megoldás? Pl periodikus-e, korlátos-e\ldots{} A megoldást
  szeretnénk jellemezni annak kiszámítása nélkül. Pl
  \(\overset{}{x} = x\) és \(x\left( 0 \right) > 0\). Ekkor egyből
  látjuk, hogy \(x\) szigmon nő, akkor is, amikor még nem tudtuk, hogy
  konkrétan mi a megoldás.
\end{enumerate}

\hypertarget{kozonseges-differencialegyenlet-megoldasanak-letezese-es-egyertelmusege}{%
\subsection{Közönséges differenciálegyenlet megoldásának létezése és
egyértelműsége}\label{kozonseges-differencialegyenlet-megoldasanak-letezese-es-egyertelmusege}}

\begin{pelda}

Példák:

\begin{itemize}
\tightlist
\item
  \(\overset{}{x}\left( t \right) = x\left( t \right)\) , ennek egy jó
  megoldása \(x\left( t \right) = c \cdot e^{t}\) ,
  \(c \in {\mathbb{R}}\) , azaz végtelen sok megoldás van. Legyen
  kezdeti feltétel: \(x\left( 0 \right) = a \in {\mathbb{R}}\) adott.
  Ekkor már csak 1 megoldás van az ilyen fajtákból:
  \(\left. c \cdot e^{0} = a\Rightarrow c = a \right.\) , vagyis a
  megoldás \(x\left( t \right) = a \cdot e^{t}\). De más fajtából
  lehetne még megoldás? Nem, ugyanis: \[\begin{array}{l}
  {\overset{}{x}\left( t \right) = x\left( t \right)} \\
  {\overset{}{x}\left( t \right) \cdot e^{- t} - x\left( t \right)e^{- t} = 0} \\
  \left. \left( {x\left( t \right) \cdot e^{- t}} \right)^{\prime} = 0\Rightarrow x\left( t \right) \cdot e^{- t} = c. \right. \\
  \end{array}\] Az implikáció csak akkor igaz, ha \(D\left( x \right)\)
  (azaz a differenciáloperátor) egy intervallumon van értelmezve. Tehát
  \(\left. \exists k \in {\mathbb{R}}:x\left( t \right)e^{- t} = k\Leftrightarrow x\left( t \right) = k \cdot e^{t} \right.\).
  A megoldás egyértelmű, mert bármilyen kezdőfeltételt adok meg, lesz
  pontosan 1 megoldás.
\item
  \(\overset{}{x}\left( t \right) = \sqrt{\left| {x\left( t \right)} \right|}\).
  Mi a megoldás \(x > 0\) -ra?
  \[\left. \frac{\overset{}{x}\left( t \right)}{\sqrt{x\left( t \right)}} = 1\Rightarrow 2\sqrt{x\left( t \right)} = t + c\Rightarrow x\left( t \right) = \left( \frac{t + c}{2} \right)^{2} \right.\]
  Hamis gyökök a parabolák "bal oldalai". \(x < 0\) esetén a megoldás
  "lefelé fordított parabolák bal oldalai", hamis megoldás a parabolák
  "jobb oldalai". \(x = 0\) esetén mindkét fajta megoldás jó. Így adott
  kezdeti feltétel mellett végtelen sok megoldás létezik. Ha
  \(x\left( t_{0} \right) = a\) a kezdeti feltétel, akkor \(a > 0\)
  esetén a megoldás csak lokálisan egyértelmű, de globálisan nem.
\end{itemize}

\end{pelda}

\hypertarget{mitol-lesz-a-megoldas-egyertelmu}{%
\subsubsection{Mitől lesz a megoldás
egyértelmű?}\label{mitol-lesz-a-megoldas-egyertelmu}}

\begin{tetel}

Tétel:\\
Ha
\(\overset{}{x}\left( t \right) = f\left( {t,x\left( t \right)} \right)\)
közönséges diffegyenletben az \(f\) függvény az \(x\) változóban
teljesíti a lokális Lipschitz feltételt, akkor a megoldás egyértelmű.
Vagyis ha minden pont egy alkalmas környezetéhez
\(\exists L \in {\mathbb{R}}^{+}:\left| {f\left( {t,p} \right) - f\left( {t,q} \right)} \right| \leq L \cdot \left| {p - q} \right|\),
akkor a megoldás egyértelmű.

\end{tetel}

\begin{pelda}

Példa:\\
\(g\left( x \right) = 5x\), vagy \(g\left( x \right) = x^{2}\)
teljesítik a lokális Lipschitz feltételt, de a
\(g\left( x \right) = \sqrt{\left| x \right|}\) már nem. Ez utóbbi 0-ban
nem lok. Lip, csak 1-ben pl.

\end{pelda}

\textbf{Észrevétel}: ha a derivált létezik, és korlátos minden pont
környezetében, akkor lok. Lip.

\begin{tetel}

Gronwall lemma (egyszerű eset):\\
(Az előző tétel bizonyítása ezen a lemmán alapszik.) Legyen
\(\left. u:\left\lbrack {a,b} \right\rbrack\rightarrow{\mathbb{R}} \right.\)
diffható, melyhez
\(\exists k \in {\mathbb{R}}^{+}:\overset{}{u}\left( t \right) \leq k \cdot u\left( t \right)\,\forall t \in \left\lbrack {a,b} \right\rbrack\).
Ekkor
\(u\left( t \right) \leq u\left( a \right) \cdot e^{k{({t - a})}}\,\forall t \in \left\lbrack {a,b} \right\rbrack\).

\end{tetel}

\begin{bizonyitas}

Bizonyítás (lemmáé):\\
Beszorzunk \(e^{- kt}\)-vel:\\
\(\begin{array}{l} {\overset{}{u}\left( t \right) \cdot e^{- kt} - k \cdot u\left( t \right) \cdot e^{- kt} \leq 0} \\ {\left( {u\left( t \right)e^{- kt}} \right)^{\prime} \leq 0} \\ {u\left( t \right)e^{- kt} \leq u\left( a \right)e^{- ka}} \\ {u\left( t \right) \leq u\left( a \right)e^{k{({t - a})}}} \\ \end{array}\)

\end{bizonyitas}

\begin{bizonyitas}

Tétel bizonyítása:\\
legyen \(x\) és \(y\) két megoldás, amelyekhez
\(\exists\tau \in {\mathbb{R}}:x\left( \tau \right) = y\left( \tau \right)\).
Belátjuk, hogy \(x\left( t \right) = y\left( t \right)\,\forall t\).
Bizonyítás \(n = 1\) esetre:
\(u\left( t \right) = \left( {x\left( t \right) - y\left( t \right)} \right)^{2}\),
\(\overset{}{u}\left( t \right) = 2\left( {x\left( t \right) - y\left( t \right)} \right) \cdot \left( {\overset{}{x}\left( t \right) - \overset{}{y}\left( t \right)} \right) = 2\left( {x\left( t \right) - y\left( t \right)} \right)\left( {f\left( {t,x\left( t \right)} \right) - f\left( {t,y\left( t \right)} \right)} \right)\).
\(\overset{}{u}\left( t \right) \leq \left| {\overset{}{u}\left( t \right)} \right| = 2\left| {x\left( t \right) - y\left( t \right)} \right| \cdot \left| {f\left( {t,x\left( t \right)} \right) - f\left( {t,y\left( t \right)} \right)} \right| \leq 2\left| {x\left( t \right) - y\left( t \right)} \right| \cdot L \cdot \left| {x\left( t \right) - y\left( t \right)} \right| = 2L \cdot u\left( t \right)\)
Gronwall alkalmazása:
\(u\left( t \right) \leq u\left( a \right) \cdot e^{2L{({t - a})}}\),
\(\left. u\left( \tau \right) = 0\Rightarrow u\left( t \right) = \left( {x\left( t \right) - y\left( t \right)} \right)^{2} \leq 0\Rightarrow x\left( t \right) = y\left( t \right)\,\forall t \geq \tau \right.\).
Hasonlóan igaz a \(t \leq \tau\)-ra is.

\end{bizonyitas}

\begin{ajanlo}

\begin{ajanlofig}

\href{https://xkcd.com}{\includegraphics[width=5.20833in,height=2.82292in]{wikipedian_protester.png}}

\end{ajanlofig}

Text

\end{ajanlo}

\hypertarget{a-hilbert-ter-geometriaja-fourier-sorfejtes}{%
\section{A Hilbert tér geometriája, Fourier
sorfejtés}\label{a-hilbert-ter-geometriaja-fourier-sorfejtes}}

Kiegésztés: fogalmaink használatához be kell vezetni a komplex
Euklideszi tér fogalmát.\\
Komplex vektortér: a definíció analóg a valós vektortér definíciójával,
kivéve: komplex számmal való szorzás is értelmezve van, a műveleti
tulajdonságok ugyanazok.

Komplex Euklideszi tér: komplex vektortér (az alaptest a komplex számok
halmaza, \(\mathbb{C}\)), plusz 2 elem skalárszorzata is értelmezve van,
értéke komplex szám. A műveleti tulajdonságok analógok, eltérés:
\(\left\langle {x,y} \right\rangle = \overline{\left\langle {y,x} \right\rangle}\)
(a felülhúzás a komplex konjugálás), ekkor amúgy
\(\left\langle {\lambda x,y} \right\rangle = \lambda\left\langle {x,y} \right\rangle\)
és
\(\left\langle {x,\lambda y} \right\rangle = \overline{\lambda}\left\langle {x,y} \right\rangle\).
(Vegyük észre, hogy a komplex vektortereken értelmezett skaláris szorzás
kétféleképp definiálható. Itt - és a matematikában általában - a
skaláris szorzás az első változójában lineáris és a másodikban konjugált
lineáris. Fizikában fordítva, azaz az első változójában lineáris, a
másodikban konjugált
lineáris:\(\left\langle {\lambda x,y} \right\rangle = \overline{\lambda}\left\langle {x,y} \right\rangle\),
illetve
\(\left\langle {x,\lambda y} \right\rangle = \lambda\left\langle {x,y} \right\rangle\).)

\begin{megjegyzes}

Példák (és megjegyzések) komplex euklideszi térre:

\begin{itemize}
\tightlist
\item
  \({\mathbb{C}}^{n}\) esetén
  \(x = \left( {x_{1},x_{2},...,x_{n}} \right)\),
  \(x_{j} \in {\mathbb{C}}\), akkor
  \(\lambda x = \left( {\lambda x_{1},\lambda x_{2},...,\lambda x_{n}} \right)\),
  \(\left\langle {x,y} \right\rangle = {\sum\limits_{j = 1}^{n}{x_{j}\overline{y_{j}}}}\)
\item
  \(L^{2}\left( M \right)\) tér (komplex esetben), ha
  \(M \subset {\mathbb{R}}^{n}\) mérhető halmaz: legyen
  \(\left. f:M\rightarrow{\mathbb{C}} \right.\),
  \(f = f_{1} + i \cdot f_{2}\). Legyen továbbá \(f_{1}\), \(f_{2}\)
  valós függvények. \(f\) mérhetősége azt jelenti, hogy \(f_{1},f_{2}\)
  mérhető
  \(\left. \Rightarrow{\int_{M}f}: = {\int_{M}f_{1}} + i{\int_{M}f_{2}} \right.\).
  \(\left. f:M\rightarrow{\mathbb{C}} \right.\) integrálható
  \(\left. \Leftrightarrow\left| f \right| \right.\) integrálható,
  \(\left. \left| f \right|:M\rightarrow{\mathbb{R}} \right.\) mérhető.
  Ekkor jelölje \(L^{2}\left( M \right)\) az olyan
  \(\left. f:M\rightarrow{\mathbb{C}} \right.\) mérhető függvények
  összességét, amelyekre \(\left| f \right|^{2}\) integrálható. Könnyen
  belátható, hogy \(L^{2}\left( M \right)\) komplex vektortér. Vezessük
  be ebben a következő skalárszorzatot:
  \(\left\langle {f,g} \right\rangle: = {\int_{M}{f\overline{g}}}\). Így
  egy Euklideszi teret kapunk. Sőt, a tér teljes, vagyis
  \(L^{2}\left( M \right)\) Hilbert tér.
\item
  Komplex \(l^{2}\) tér,
  \(x: = \left( {x_{1},x_{2},...,x_{j},...} \right)\),
  \(x_{j} \in {\mathbb{C}}\), \(l^{2}\) komplex euklideszi tér, ebben a
  skaláris szorzás
  \(\left\langle {x,y} \right\rangle = {\sum\limits_{j = 1}^{\infty}{x_{j}\overline{y_{j}}}}\).
  Bizonyítható, hogy teljes is.
\end{itemize}

\end{megjegyzes}

\hypertarget{ortogonalis-kiegeszito-alter}{%
\subsection{Ortogonális kiegészítő
altér}\label{ortogonalis-kiegeszito-alter}}

\begin{definicio}

Definíció:\\
Legyen \(X\) Hilbert tér (vagy akár Banach is). Egy \(Y \subset X\)
halmazt altérnek nevezzük, ha az összeadás és számmal való szorzás nem
vezet ki belőle és zárt részhalmaz (a konvergencia nem vezet ki).

\end{definicio}

\begin{definicio}

Definíció:\\
Legyen \(X\) Hilbert tér, s két eleme \(x\) és \(y\). Ezek merőlegesek,
vagyis \(x\bot y\), ha \(\left\langle {x,y} \right\rangle = 0\).

\end{definicio}

\begin{definicio}

Definíció:\\
Legyen \(X\) Hilbert tér, \(Y \subset X\) altér. Azt mondjuk, hogy az
\(x \in X\) elem \(Y\) ortogonális, ha \(\forall y \in Y\)-ra
\(\left\langle {x,y} \right\rangle = 0\).

\end{definicio}

\begin{definicio}

Definíció:\\
Legyen \(X\) Hilbert tér, \(Y \subset X\) altér. Az \(Y\) altér
ortogonális kiegészítő altérét, \(Y^{\bot}\)-t így értelmezzük:
\(Y^{\bot}: = \left\{ {x \in X:x\bot Y} \right\}\).

\end{definicio}

\begin{allitas}

Állítás:\\
\(Y^{\bot} \subset X\) is altér.

\end{allitas}

\begin{bizonyitas}

Bizonyítás:\\
Az összeadás és számmal való szorzás nem vezet ki belőle, ugyanis tfh
\(y_{1},y_{2} \in Y^{\bot}\), \(x \in Y\) tetszőleges. Ekkor
\(\left\langle {\lambda_{1}y_{1} + \lambda_{2}y_{2},x} \right\rangle = \lambda_{1}\left\langle {y_{1},x} \right\rangle + \lambda_{2}\left\langle {y_{2},x} \right\rangle = 0\).
\(Y^{\bot}\) zárt halmaz, ugyanis legyen \(y_{j} \in Y^{\bot}\),
\(\lim\left( y_{j} \right) = y \in X\). Tudjuk, hogy
\(\left\langle {y_{j},x} \right\rangle = 0\,\forall x \in Y\).
\(\left. y_{j}\rightarrow y\Rightarrow\left\langle {y_{j},x} \right\rangle\rightarrow\left\langle {y,x} \right\rangle \right.\)
minden rögzített \(x\)-re, ugyanis a skalárszorzat a tényezőktől
folytonosan függ, tehát \(\left\langle {y,x} \right\rangle = 0\),
\(\forall x \in X\)-re, vagyis \(y \in Y^{\bot}\).

\end{bizonyitas}

\begin{megjegyzes}

Megjegyzés:\\
Komplex Cauchy-Schwarz egyenlőtlenség, azaz
\(\left| \left\langle {x,y} \right\rangle \right| \leq \left\| x \right\| \cdot \left\| y \right\|\)
bizonyítása:\\
\(\begin{array}{l} {0 \leq \left\langle {x + \lambda y,x + \lambda y} \right\rangle = \left\langle {x,x} \right\rangle + \lambda\left\langle {y,x} \right\rangle + \overline{\lambda}\left\langle {x,y} \right\rangle + \lambda\overline{\lambda}\left\langle {y,y} \right\rangle} \\ {0 \leq \left\langle {x,x} \right\rangle + \lambda\left\langle {y,x} \right\rangle + \overline{\lambda}\left\lbrack {\left\langle {x,y} \right\rangle + \lambda\left\langle {y,y} \right\rangle} \right\rbrack} \\ \end{array}\)\\
A \(\lambda \in {\mathbb{C}}\) számot válasszuk meg úgy, hogy
\(\overline{\lambda}\) együtthatója 0 legyen. Ez teljesül, ha
\(\lambda = - \frac{\left\langle {x,y} \right\rangle}{\left\langle {y,y} \right\rangle}\)
(\(y = 0\) triviális eset, így feltesszük, hogy \(y \neq 0\)),
behelyettesítve:
\(\left. 0 \leq \left\langle {x,x} \right\rangle - \frac{\left\langle {x,y} \right\rangle}{\left\langle {y,y} \right\rangle}\left\langle {y,x} \right\rangle = \left\langle {x,x} \right\rangle - \frac{\left| \left\langle {x,y} \right\rangle \right|^{2}}{\left\langle {y,y} \right\rangle}\Rightarrow\left| \left\langle {x,y} \right\rangle \right|^{2} \leq \left\langle {x,x} \right\rangle\left\langle {y,y} \right\rangle \right.\).

\end{megjegyzes}

\begin{tetel}

Riesz-féle felbontási tétel:\\
Legyen \(X\) Hilbert tér, \(Y\) egy altere, \(Y^{\bot}\) az \(Y\)-nak
ortogonális kiegészítő altere! Ekkor \(\forall x \in X\) elemre
\(x = y + z\), ahol \(y \in Y\), \(z \in Y^{\bot}\) és a felbontás
egyértelmű.

\end{tetel}

\begin{allitas}

Lemma (paralelogramma egyenlőség):\\
Legyen \(X\) egy Hilbert tér. Ekkor \(\forall a,b \in X\) esetén
\(\left\| {a + b} \right\|^{2} + \left\| {a - b} \right\|^{2} = 2\left\| a \right\|^{2} + 2\left\| b \right\|^{2}\).

\end{allitas}

\begin{bizonyitas}

Bizonyítás (lemmáé):\\
\(\left\| {a + b} \right\|^{2} + \left\| {a - b} \right\|^{2} = \left\langle {a + b,a + b} \right\rangle + \left\langle {a - b,a - b} \right\rangle =\)
\(= \left\| a \right\|^{2} + \left\| b \right\|^{2} + \left\langle {a,b} \right\rangle + \left\langle {b,a} \right\rangle + \left\| a \right\|^{2} + \left\| b \right\|^{2} - \left\langle {a,b} \right\rangle - \left\langle {b,a} \right\rangle = 2\left\| a \right\|^{2} + 2\left\| b \right\|^{2}\).

\end{bizonyitas}

\begin{bizonyitas}

Bizonyítás (tételé):\\
Legyen
\(d: = \inf\left\{ {\left\| {x - y} \right\|:y \in Y} \right\} \geq 0\)
(\(d\) véges). Belátjuk, hogy
\[\exists y_{0} \in Y:\left\| {x - y_{0}} \right\| = d.\] Az infinimum
definíciója miatt
\[\exists y_{j} \in Y:d^{2} \leq \left\| {x - y_{j}} \right\|^{2} < d^{2} + 1/j\quad j \in {\mathbb{N}}.\]Tekintsük
az \(\left( y_{j} \right)\) sorozatot!\\
Állítás: \(\left( y_{j} \right)\) Cauchy sorozat. Ehhez felhasználjuk a
paralelogramma egyenlőséget \(a: = x - y_{j}\) , \(b: = x - y_{k}\)
választással.
\[\left\| {\left( {x - y_{j}} \right) + \left( {x - y_{k}} \right)} \right\|^{2} + \left\| {\left( {x - y_{j}} \right) + \left( {x - y_{k}} \right)} \right\|^{2} = 2\left\| {x - y_{j}} \right\|^{2} + 2\left\| {x - y_{k}} \right\|^{2}\]
\[\begin{aligned}
  {\left\| {{y_k} - {y_j}} \right\|^2} =  & 2{\left\| {x - {y_j}} \right\|^2} + 2{\left\| {x - {y_k}} \right\|^2} - \underbrace {{{\left\| {2x - \left( {{y_j} + {y_k}} \right)} \right\|}^2}}_{4{{\left\| {x - \frac{{{y_j} + {y_k}}}{2}} \right\|}^2}} \\ 
   \leqslant  & 2\left( {{d^2} + 1/j} \right) + 2\left( {{d^2} + 1/k} \right) - 4{d^2} \\ 
   =  & \frac{2}{j} + \frac{2}{k} < \varepsilon , \text{ ha } j,k \geqslant {j_0}. \\ 
\end{aligned} \] Mivel \(X\) tér teljes
\(\left. \Rightarrow\exists y_{0} \in X:\lim\limits_{j\rightarrow\infty}\left\| {y_{j} - y_{0}} \right\| = 0 \right.\).
Mivel \(Y\) altér zár halmaz
\(\left. \Rightarrow y_{0} = \lim\left( y_{j} \right) \in Y \right.\).

Másrészt (mint ahogy már írtuk)
\(d = \inf\left\{ {\left\| {x - y} \right\|:y \in Y} \right\}\) ,
\(d^{2} \leq \left\| {x - y_{j}} \right\|^{2} < d^{2} + \frac{1}{j}\),
továbbá
\[\lim \left( {{y_j}} \right) = {y_0}\; \Rightarrow \;\lim \left\| {x - {y_j}} \right\| = \left\| {x - {y_0}} \right\|\; \Rightarrow \;{\left\| {x - {y_0}} \right\|^2} = {d^2}.\]
Legyen \(z_{0} = x - y_{0}\). Be kellene még látni, hogy \(z_{0}\bot Y\)
, vagyis \(x = y_{0} + z_{0}\) , ahol
\(y_{0} \in Y,z_{0} \in Y^{\bot}\).

Legyen \(y \in Y\) ! Mivel \(d\) a fenti infinimum, ezért tetszőleges
\(\lambda \in {\mathbb{K}}\) esetén
\(d^{2} = \left\| {x - y_{0}} \right\|^{2} \leq \left\| {x - y_{0} - \lambda y} \right\|^{2} =\)
\(= \left\| {z_{0} - \lambda y} \right\|^{2} = \left\langle {z_{0} - \lambda y,z_{0} - \lambda y} \right\rangle = \left\| z_{0} \right\|^{2} - \lambda\left\langle {y,z_{0}} \right\rangle - \overline{\lambda}\left\lbrack {\left\langle {z_{0},y} \right\rangle - \lambda\left\| y \right\|^{2}} \right\rbrack\).
Most \(\lambda\) -t megint úgy választjuk, hogy \(\overline{\lambda}\)
együtthatója 0 legyen, vagyis legyen
\(\lambda = \frac{\left\langle {z_{0},y} \right\rangle}{\left\| y \right\|^{2}}\)
(megint feltehetjük, hogy \(y \neq 0\) ). Tehát
\(d^{2} \leq d^{2} - \lambda\left\langle {y,z_{0}} \right\rangle = d^{2} - \frac{\left\langle {z_{0},y} \right\rangle}{\left\| y \right\|^{2}}\left\langle {y,z_{0}} \right\rangle = d^{2} - \frac{\left| \left\langle {z_{0},y} \right\rangle \right|^{2}}{\left\| y \right\|^{2}}\)
,
\(\left. 0 \leq - \frac{\left| \left\langle {z_{0},y} \right\rangle \right|^{2}}{\left\| y \right\|^{2}}\Rightarrow\left| \left\langle {z_{0},y} \right\rangle \right|^{2} = 0\Rightarrow\left\langle {z_{0},y} \right\rangle = 0 \right.\).
Tehát \(z_{0} \in Y^{\bot}\), vagyis valóban lehetséges ilyen
felbontás.\\
Indirekt bizonyítjuk, hogy a felbontás egyértelmű. Tfh két alakban is
felírható \(x\) : \(x = y_{0} + z_{0} = y_{1} + z_{1}\), ahol
\(y_{1},y_{2} \in Y\) és \(z_{1},z_{2} \in Y^{\bot}\).
\(Y \ni \left( {y_{0} - y_{1}} \right): = a = \left( {z_{1} - z_{0}} \right) \in Y^{\bot}\).

\(\left. \left\langle {y_{0} - y_{1},z_{1} - z_{0}} \right\rangle = \left\| a \right\|^{2} = 0\Rightarrow y_{0} - y_{1} = z_{0} - z_{1} = 0\Rightarrow y_{0} = y_{1},z_{0} = z_{1} \right.\)

\end{bizonyitas}

\begin{ajanlo}

\begin{ajanlofig}

\href{https://xkcd.com}{\includegraphics[width=5.20833in,height=2.82292in]{wikipedian_protester.png}}

\end{ajanlofig}

Text

\end{ajanlo}

\hypertarget{ortogonalis-rendszerek}{%
\subsubsection{Ortogonális rendszerek}\label{ortogonalis-rendszerek}}

\begin{definicio}

Definíció:\\
Egy \(X\) vektortérben az \(M\) halmaz elemei lineárisan függetlenek, ha
bármely véges sok lineárisan független.

\end{definicio}

\begin{definicio}

Definíció:\\
Legyen \(X\) normált tér! \(X\) dimenziója az olyan lineárisan független
elemek maximális száma, amelyek véges lineárkombinációi mindenütt sűrűn
vannak \(X\)-ben (egy \(A \subset X\) sűrű \(X\)-ben, ha
\(\overline{A} = X\), ahol a halmaz felülvonása a lezárást jelenti, ez
amúgy ekvivalens azzal, hogy \(\forall x \in X\)-nek minden
környezetében van \(A\)-beli elem). Másképp fogalmazva: jelöljük
\(\mathcal{L}\left( {x_{1},x_{2},...} \right)\)-val azt a lineáris
teret, amely az \(x_{1},x_{2},...\) elemek véges lineárkombinációjaként
előáll. (Az előálló lineáris tér egyértelmű, de egy teret több ilyen
vektorrendszer is előállíthat.) Ekkor \(X\) tér dimenziója az olyan
lineárisan független elemek maximális száma, melyekre
\(\overline{\mathcal{L}\left( {x_{1},x_{2},...} \right)} = X\). A \(D\)
dimenziószám egyértelmű, \(0 \leq D \leq \infty\).

\end{definicio}

\begin{definicio}

Definíció:\\
Egy \(X\) normált teret szeparábilisnak nevezünk, ha benne megadható
megszámlálhatóan sok (azaz véges vagy megszámlálhatóan végtelen sok)
lineárisan független elem, amelyek véges lineárkombinációi sűrűn vannak
\(X\)-ben.

\end{definicio}

\begin{definicio}

Definíció:\\
Legyen \(X\) Hilbert-tér! Azt mondjuk, hogy az
\(x_{1},x_{2},...,x_{k},...\) elemek ortogonális rendszert alkotnak, ha
\(\forall x_{j},x_{k} \neq 0\) esetén
\(\left\langle {x_{j},x_{k}} \right\rangle = \begin{cases} 0 & {j \neq k} \\ {\text{nem}0} & {j = k} \\ \end{cases}\).
A rendszer ortonormált, ha \(\forall x \in X\) esetén
\(\left\| x \right\| = 1\).

\end{definicio}

\textbf{Kérdés}: ha az \(X\) Hilbert-térben
\(y_{1},y_{2},...,y_{k},...\) lineárisan függetlenek, akkor lehet-e
ezekből ortonormált rendszert konstruálni, és ha igen, hogyan? Válasz:
lehet, az ún. Schmidt-féle ortogonalizációs eljárással.

\begin{tetel}

Tétel:\\
Az \(y_{1},y_{2},...,y_{k},...\) lineárisan független elemekhez az
\(x_{1},x_{2},...,x_{k},...\) elemek megkonstruálhatók úgy, hogy az
utóbbiak ortonormált rendszert alkossanak, mégpedig úgy, hogy
\(\forall k\)-ra
\(\mathcal{L}\left( {x_{1},x_{2},...,x_{k}} \right) = \mathcal{L}\left( {y_{1},y_{2},...,y_{k}} \right)\).

\end{tetel}

\begin{bizonyitas}

Bizonyítás:

\begin{enumerate}
\def\labelenumi{\arabic{enumi}.}
\tightlist
\item
  Legyen \(x_{1} = \frac{y_{1}}{\left\| y_{1} \right\|}\), ekkor
  \(\left\| x_{1} \right\| = 1\). \(y_{1} \neq 0\) , mert
  \(y_{1},y_{2},...\) lineárisan függetlenek.
\item
  \(z_{2}: = y_{2} - \lambda_{1}x_{1}\), ahol
  \(\lambda_{1} \in {\mathbb{R}}\). Ezt hogy válasszuk meg, hogy
  \(z_{2}\bot x_{1}\) teljesüljön?
  \(\left. 0 = \left\langle {z_{2},x_{1}} \right\rangle = \left\langle {y_{2} - \lambda_{1}x_{1},x_{1}} \right\rangle = \left\langle {y_{2},x_{1}} \right\rangle - \lambda_{1}\underset{= 1}{\underset{\}\ }{\left\langle {x_{1},x_{1}} \right\rangle}}\Rightarrow\lambda_{1} = \left\langle {y_{2},x_{1}} \right\rangle \right.\).
  Ekkor \(z_{2} \neq 0\), mert \(y_{1},y_{2}\) lineárisan függetlenek.
  \(x_{2}: = \frac{z_{2}}{\left\| z_{2} \right\|}\), ekkor
  \(\left\| x_{2} \right\| = 1\) és
  \(\left\langle {x_{1},x_{2}} \right\rangle = 0\).
\item
  \(z_{3}: = y_{3} - \mu_{1}x_{1} - \mu_{2}x_{2}\), ahol
  \(\mu_{1},\mu_{2} \in {\mathbb{R}}\). Ezeket hogy válasszuk meg, hogy
  \(z_{3}\bot x_{1},x_{2}\) teljesüljenek? \[\begin{array}{l}
  \left. 0 = \left\langle {y_{3} - \mu_{1}x_{1} - \mu_{2}x_{2},x_{1}} \right\rangle = \left\langle {y_{3},x_{1}} \right\rangle - \mu_{1} - 0\Leftrightarrow\mu_{1} = \left\langle {y_{3},x_{1}} \right\rangle \right. \\
  \left. 0 = \left\langle {y_{3} - \mu_{1}x_{1} - \mu_{2}x_{2},x_{2}} \right\rangle = \left\langle {y_{3},x_{2}} \right\rangle - 0 - \mu_{2}\Leftrightarrow\mu_{2} = \left\langle {y_{3},x_{2}} \right\rangle \right. \\
  \end{array}\] \(z_{3} \neq 0\) \(y_{1},y_{2},y_{3}\) lineáris
  függetlensége miatt, ezért
  \(x_{3}: = \frac{z_{3}}{\left\| z_{3} \right\|}\) jó választás, így
  \(\left\| x_{3} \right\| = 1\) és \(x_{3}\bot x_{1},x_{2}\).
\end{enumerate}

Nem nehéz belátni, hogy az eljárás folytatható \(\forall k\)-ra és
\(\mathcal{L}\left( {y_{1},y_{2},...,y_{k}} \right) = \mathcal{L}\left( {x_{1},x_{2},...,x_{k}} \right)\).

\end{bizonyitas}

\hypertarget{ortogonalis-sorok-fourier-sorok}{%
\subsubsection{Ortogonális sorok,
Fourier-sorok}\label{ortogonalis-sorok-fourier-sorok}}

A továbbiakban legyen \(X\) szeparábilis Hilbert-tér, véges vagy
végtelen dimenziós! Tudjuk, hogy ekkor \(X\)-ben megadható
\(x_{1},x_{2},...,x_{k},...\) ortonormált rendszer. Egy
\(\sum\limits_{k}{c_{k}x_{k}}\) alakú sort (összeget) -- ahol
\(c_{k} \in {\mathbb{K}}\) -- ortogonális sornak nevezünk.

\begin{tetel}

Tételek:

\begin{enumerate}
\def\labelenumi{\arabic{enumi}.}
\tightlist
\item
  Egy \(\sum\limits_{k}{c_{k}x_{k}}\) sor konvergens
  \(\left. \Leftrightarrow{\sum\limits_{k}{\left| c_{k} \right|^{2} < \infty}} \right.\)
\item
  Ha \(x = {\sum\limits_{k}{c_{k}x_{k}}}\), akkor
  \(c_{l} = \left\langle {x,x_{l}} \right\rangle\)
\item
  \(\left\| x \right\|^{2} = {\sum\limits_{k}\left| c_{k} \right|^{2}}\)
  (végtelen dimenziós Pitagorasz tétel).
\end{enumerate}

\end{tetel}

\begin{bizonyitas}

Bizonyítás:

\begin{enumerate}
\def\labelenumi{\arabic{enumi}.}
\tightlist
\item
  Véges dimenzióban triviális, így tegyük fel, hogy végtelen sok elemű
  az ortonormált rendszer! Legyen
  \(s_{j}: = {\sum\limits_{k = 1}^{j}{c_{k}x_{k}}}\)! A sor
  konvergenciája azt jelenti, hogy \(\left( s_{j} \right)\) sorozat
  konvergens \(\left. \Leftrightarrow\left( s_{j} \right) \right.\)
  Cauchy sorozat. \[\begin{aligned}
    {\left\| {{s_j} - {s_l}} \right\|^2} &  = \left\langle {{s_j} - {s_l},{s_j} - {s_l}} \right\rangle  \\ 
     &  = \left\langle {\mathop \sum \limits_{k = l + 1}^j {c_k}{x_k},\mathop \sum \limits_{k = l + 1}^j {c_k}{x_k}} \right\rangle  \\ 
     &  = \mathop \sum \limits_{k = l + 1}^j {c_k}{{\bar c}_k}\underbrace {\left\langle {{x_k},{x_k}} \right\rangle }_{ = 1} = \mathop \sum \limits_{k = l + 1}^j {\left| {{c_k}} \right|^2} \\ 
  \end{aligned} \] Ez a
  \(\sum\limits_{k = 1}^{\infty}\left| c_{k} \right|^{2}\) sor egy
  „szelete``. Tehát \(\left( s_{j} \right)\) \(X\)-beli sorozatra
  teljesül a Cauchy-kritérium
  \(\left. \Leftrightarrow{\sum\limits_{k = 1}^{\infty}\left| c_{k} \right|^{2}} \right.\)
  sorra teljesül a Cauchy-kritérium
  \(\left. \Leftrightarrow\left( s_{j} \right) \right.\) \(X\)-beli
  sorozat konvergens
  \(\left. \Leftrightarrow{\sum\limits_{k = 1}^{\infty}\left| c_{k} \right|^{2}} \right.\)
  sor konvergens.
\item
  Tfh \(x = {\sum\limits_{k}{c_{k}x_{k}}}\), \(x_{l}\)-lel szorozzuk
  skalárisan (jobbról) az egyenlőséget (ezt megtehetjük, hisz nem nehéz
  belátni, hogy egy konvergens sor tagonként szorozható skalárisan),
  \(\left\langle {x,x_{l}} \right\rangle = \left\langle {\sum\limits_{k}{c_{k}x_{k},x_{l}}} \right\rangle = {\sum\limits_{k}{c_{k}\left\langle {x_{k},x_{l}} \right\rangle = c_{l}}}\)
\item
  \({\left\| x \right\|^2} = \left\langle {x,x} \right\rangle = \left\langle {\mathop \sum \limits_k {c_k}{x_k},x} \right\rangle = \mathop \sum \limits_k {c_k}\underbrace {\left\langle {{x_k},x} \right\rangle }_{{{\bar c}_k}} = \mathop \sum \limits_k {\left| {{c_k}} \right|^2}\)
\end{enumerate}

\end{bizonyitas}

\begin{definicio}

Definíció:\\
Legyen \(x_{1},x_{2},...,x_{k}\) ortonormált rendszer, \(x \in X\) adott
elem! Értelmezzük az \(x\) elem \(k\)-adik Fourier-együtthatóját:
\(c_{k}: = \left\langle {x,x_{k}} \right\rangle\). Az így adódó
\(\sum\limits_{k}{c_{k}x_{k}}\) „sort" az \(x\) elem Fourier-sorának
nevezzük.

\end{definicio}

\textbf{Kérdés}: egy \(x\) elem Fourier-sora konvergens-e? Ha igen, mi
az összege?

\begin{tetel}

Tétel:\\
Egy \(x \in X\) elem Fourier sora mindig konvergens, ugyanis teljesül az
ún. Bessel-egyenlőtlenség:
\({\sum\limits_{k}\left| c_{k} \right|^{2}} \leq \left\| x \right\|^{2}\).
A sor összege pontosan akkor \(x\), ha teljesül az ún Parseval
egyenlőség, azaz
\({\sum\limits_{k}\left| c_{k} \right|^{2}} = \left\| x \right\|^{2}\).

\end{tetel}

\begin{bizonyitas}

Bizonyítás:\\
\(s_{j}: = {\sum\limits_{k = 1}^{j}{c_{k}x_{k}}}\), ekkor
\[\begin{array}{l}
{0 \leq \left\| {x - s_{j}} \right\|^{2} = \left\langle {x - s_{j},x - s_{j}} \right\rangle} \\
{= \left\| x \right\|^{2} - \left\langle {s_{j},x} \right\rangle - \left\langle {x,s_{j}} \right\rangle + \left\| s_{j} \right\|^{2}} \\
{= \left\| x \right\|^{2} - \left\langle {\sum\limits_{k = 1}^{j}c_{k}x_{k},x} \right\rangle - \left\langle {x,\sum\limits_{k = 1}^{j}c_{k}x_{k}} \right\rangle + \left\langle {\sum\limits_{k = 1}^{j}c_{k}x_{k},\sum\limits_{k = 1}^{j}c_{k}x_{k}} \right\rangle} \\
{= \left\| x \right\|^{2} - \sum\limits_{k = 1}^{j}c_{k}{\overline{c}}_{k} - \sum\limits_{k = 1}^{j}{\overline{c}}_{k}c_{k} + \sum\limits_{k = 1}^{j}c_{k}{\overline{c}}_{k}} \\
\left. {}{} = \left\| x \right\|^{2} - \sum\limits_{k = 1}^{j}\left| c_{k} \right|^{2}\Leftrightarrow\sum\limits_{k = 1}^{j}\left| c_{k} \right|^{2} \leq \left\| x \right\|^{2}\Rightarrow\sum\limits_{k = 1}^{\infty}\left| c_{k} \right|^{2} \leq \left\| x \right\|^{2}, \right. \\
\end{array}\]másrészt a fentiek szerint
\(\left\| {x - s_{j}} \right\|^{2} = \left\| x \right\|^{2} - {\sum\limits_{k = 1}^{j}\left| c_{k} \right|^{2}}\).
Ebből láthatjuk, hogy
\(\left. s_{j}\rightarrow x\Leftrightarrow\left\| x \right\|^{2} - {\sum\limits_{k}\left| c_{k} \right|^{2}} = 0 \right.\),
vagyis a sor összege pontosan akkor \(x\), ha
\(\left\| x \right\|^{2} - {\sum\limits_{k}\left| c_{k} \right|^{2}} = 0\).

\end{bizonyitas}

\begin{tetel}

Tétel:\\
Legyen \(x_{1},x_{2},...,x_{k},...\) ortonormált rendszer. Ekkor egy
\(x \in X\) elem Fourier-sorának összege az \(x\) elemnek az
\(X_{0}: = \overline{\mathcal{L}\left( {x_{1},x_{2},...,x_{k},...} \right)} \subset X\)
alterén vett merőleges vetülete.

\end{tetel}

\begin{bizonyitas}

Bizonyítás:\\
Jelölje \(x^{*}: = {\sum\limits_{k}{c_{k}x_{k}}}\), ahol
\(c_{k}: = \left\langle {x,x_{k}} \right\rangle\). Azt kellene belátni,
hogy \(x^{*} \in X_{0}\) és \(\left( {x - x^{*}} \right)\bot X_{0}\).
\(x^{*} \in X_{0}\), ugyanis
\({\sum\limits_{k = 1}^{j}{c_{k}x_{k}}} \in \mathcal{L}\left( {x_{1},x_{2},...,x_{j}} \right)\),
így \({\sum\limits_{k}{c_{k}x_{k}}} \in X_{0}\).
\(\left( {x - x^{*}} \right)\bot X_{0}\) ugyanis először legyen
\(y \in \mathcal{L}\left( {x_{1},x_{2},...,x_{l}} \right)\) tetszőleges!
Belátjuk, hogy \(\left\langle {x - x^{*},y} \right\rangle = 0\).
\(y = {\sum\limits_{j = 1}^{l}{d_{j}x_{j}}}\), \[\begin{aligned}
  \left\langle {x - {x^*},y} \right\rangle  &  = \left\langle {x,y} \right\rangle  - \left\langle {{x^ * },y} \right\rangle  \\ 
   &  = \left\langle {x,\mathop \sum \limits_{j = 1}^l {d_j}{x_j}} \right\rangle  - \left\langle {\mathop \sum \limits_k {c_k}{x_k},\mathop \sum \limits_{j = 1}^l {d_j}{x_j}} \right\rangle  \\ 
   &  = \mathop \sum \limits_{j = 1}^l {{\bar d}_j}\underbrace {\left\langle {x,{x_j}} \right\rangle }_{{c_j}} - \mathop \sum \limits_{j = 1}^l {{\bar d}_j}\underbrace {\left\langle {\mathop \sum \limits_k {c_k}{x_k},{x_j}} \right\rangle }_{{c_j}} = 0 .\\ 
\end{aligned} \] Most legyen
\(y \in X_{0} = \overline{\mathcal{L}\left( {x_{1},x_{2},...} \right)}\),
szeretnénk, ha ekkor \(\left\langle {x - x^{*},y} \right\rangle = 0\) is
igaz lenne. Ehhez vegyünk egy \(\left( y_{\nu} \right)\),
\(\mathcal{L}\left( {x_{1},x_{2},...} \right)\)-beli konvergens
sorozatot, melyre \(\left. y_{\nu}\rightarrow y \right.\). Ekkor
\(\left\langle {x - x^{*},y_{\nu}} \right\rangle = 0\). Így, mivel
\(\left. y_{\nu}\rightarrow y \right.\),
\(\left\langle {x - x^{*},y} \right\rangle = 0\), ugyanis
\[\begin{aligned}
  \left| {\left\langle {x - {x^ * },y} \right\rangle } \right| &  = \left| {\left\langle {x - {x^ * },y} \right\rangle  - \left\langle {x - {x^ * },{y_\nu }} \right\rangle } \right| \\ 
   &  = \left| {\left\langle {x - {x^ * },y - {y_\nu }} \right\rangle } \right| \\ 
   &  \leqslant \left\| {x - {x^ * }} \right\| \cdot \underbrace {\left\| {y - {y_\nu }} \right\|}_{ \to 0} \to 0. \\ 
\end{aligned} \]

\end{bizonyitas}

\begin{definicio}

Definíció:\\
Az \(x_{1},x_{2},...\) ortonormált rendszert zártnak nevezzük, ha
\(\overline{\mathcal{L}\left( {x_{1},x_{2},...} \right)} = X\).

\end{definicio}

Következmény: ha az \(x_{1},x_{2},...\) ortonormált rendszer zárt, akkor
\(\forall x \in X\) elem Fourier-sorának összege \(x\).

\begin{definicio}

Definíció:\\
Egy \(x_{1},x_{2},...\) ortonormált rendszert teljesnek nevezzük, ha
\(\left. x\bot x_{k}\,\forall k\Rightarrow x = 0 \right.\).

\end{definicio}

\begin{tetel}

Tétel (bizonyítás nélkül):\\
Egy \(x_{1},x_{2},...\) ortonormált rendszer teljes \(\Leftrightarrow\)
zárt.

\end{tetel}

\begin{ajanlo}

\begin{ajanlofig}

\href{https://xkcd.com}{\includegraphics[width=5.20833in,height=2.82292in]{wikipedian_protester.png}}

\end{ajanlofig}

Text

\end{ajanlo}

\begin{pelda}

Példák zárt (teljes) ortonormált rendszerekre.

Észrevétel: ha egy \(y_{1},y_{2},...,y_{k},...\) lineárisan független
rendszer olyan, hogy
\(\overline{\mathcal{L}\left( {y_{1},y_{2},...} \right)} = X\)
(\(ha\ pl.\ X\) Hilbert-tér, akkor már a lineárisan független rendszer
is zárt, ha kifeszíti a teljes teret, mert a Hilbert-tér zárt), akkor
ebből a Schmidt ortogonalizálási eljárással zárt (teljes) ortonormált
rendszert kapunk.

\begin{enumerate}
\def\labelenumi{\arabic{enumi}.}
\tightlist
\item
  \(X: = L^{2}\left( {a,b} \right)\) , ahol \(\left( {a,b} \right)\)
  véges intervallum.

  \begin{tetel}

  Tétel:\\
  Ebben az
  \(\left. t\mapsto 1,t\mapsto t,t\mapsto t^{2},...,t\mapsto t^{k},... \right.\)
  lineárisan független függvények zárt rendszert alkotnak.

  \end{tetel}

  \begin{bizonyitas}

  Bizonyítás (vázlat):\\
  Egyrészt a függvényrendszer lineárisan független:
  \[\left. {\sum\limits_{j = 0}^{k}{a_{j}t^{j}}} = 0\Leftrightarrow a_{j} = 0. \right.\](Egy
  valós \(k\) -ad fokú polinomnak legfeljebb \(k\) db gyöke lehet
  \(k \geq 1\).) Az, hogy a rendszer zárt, következik a Weierstrass
  approximációs tételéből. Eszerint tetszőleges
  \(\left. f:\left\lbrack {a,b} \right\rbrack\rightarrow{\mathbb{R}} \right.\)
  folytonos függvényhez \(\exists P_{k}\) polinom sorozat, amely
  egyenletesen tart \(f\) -hez. Legyen
  \(\left. g:\left( {a,b} \right)\rightarrow{\mathbb{R}},g \in L^{2}\left( {a,b} \right) \right.\).
  A Lebesgue integrál felépítéséből kiolvasható, hogy
  \(\left. g:\left\lbrack {a,b} \right\rbrack\rightarrow{\mathbb{R}} \right.\)
  folytonos függvények sűrűn vannak \(L^{2}\left( {a,b} \right)\) -n. A
  \(g\) folytonos függvényt Weierstrass approximációs tétele szerint
  tetszőleges előírt pontossággal meg lehet közelíteni polinomokkal, a
  szuprémum normában \(\Rightarrow\) ezek közelítik \(g\) -t \(L^{2}\)
  normában is.

  \end{bizonyitas}
\item
  Komplex trigonometrikus rendszer \(X: = L^{2}\left( {0,2\pi} \right)\)
  ,
  \(\phi_{k}\left( t \right): = e^{ikt},t \in \left( {0,2\pi} \right),k \in {\mathbb{Z}}\).

  \begin{tetel}

  Tétel:\\
  A fenti függvények egy zárt ortogonális rendszert alkotnak (biz.
  nélkül). Belátjuk, hogy
  \(\left( \phi_{k} \right)_{k \in {\mathbb{Z}}}\) ortogonális.
  \({\int\limits_{0}^{2\pi}{\phi_{k}\left( t \right)\overline{\phi_{l}\left( t \right)}dt}} = {\int\limits_{0}^{2\pi}{e^{ikt}e^{- ilt}dt}} = {\int\limits_{0}^{2\pi}e^{i{({k - l})}t}} = \left\lbrack \frac{e^{i{({k - l})}t}}{i\left( {k - l} \right)} \right\rbrack_{t = 0}^{2\pi} = 0\)
  ha \(k \neq l\). \(\psi_{k}: = \frac{1}{\sqrt{2\pi}}\phi_{k}\) már
  ortonormált rendszer.

  \end{tetel}
\item
  Valós trigonometrikus rendszerek.\\
  Legyen az \(X\) alaphalmaz a valós \(L^{2}\left( {0,2\pi} \right)\).
  \(e^{ikt} = \cos\left( {kt} \right) + i\sin\left( {kt} \right)\) ,
  \(\cos\left( {kt} \right) = \frac{e^{ikt} + e^{- ikt}}{2}\) ,
  \(\sin\left( {kt} \right) = \frac{e^{ikt} - e^{- ikt}}{2i}\). Egyszerű
  számolással adódik, hogy
  \(1,\cos t,\sin t,\cos\left( {2t} \right),\sin\left( {2t} \right),...,\cos\left( {kt} \right),\sin\left( {kt} \right),...\)
  függvények páronként merőlegesek. Tehát ezek ortogonális rendszert
  alkotnak a valós \(L^{2}\left( {0,2\pi} \right)\) -ben. Abból, hogy a
  komplex trigonometrikus rendszer zárt \(\Rightarrow\) a fenti rendszer
  valós ortogonális zárt rendszer. A fentiekből következik, hogy egy
  tetszőleges \(f \in L^{2}\left( {0,2\pi} \right)\) függvénynek akár a
  komplex, akár a valós trigonometrikus rendszer szerint Fourier sora
  előállítja a függvényt \(L^{2}\) normában.
\item
  Az
  \(1,\cos t,\cos\left( {2t} \right),...,\cos\left( {kt} \right),...\)
  függvényrendszer zárt és ortogonális a \(L^{2}\left( {0,\pi} \right)\)
  -ben. A szinuszos ugyanígy.
\end{enumerate}

\end{pelda}

\hypertarget{linearis-es-korlatos-operatorok}{%
\subsection{Lineáris és korlátos
operátorok}\label{linearis-es-korlatos-operatorok}}

\begin{allitas}

Állítás:\\
Legyen \(X\) , \(Y\) normált terek! Korábban bizonyítottuk, hogy
\(\left. A:X\rightarrow Y \right.\) lineráis operátor folytonos
\(\Leftrightarrow\) \(A\) korlátos.

\end{allitas}

\begin{definicio}

Definíció:\\
Egy \(\left. A:X\rightarrow Y \right.\) lineáris operátort korlátosnak
nevezzük, ha
\(\exists c \geq 0:\left\| {Ax} \right\|_{Y} \leq c\left\| x \right\|_{X}\,\forall x \in X\).

\end{definicio}

\begin{tetel}

Tétel:\\
Legyen \(X\) normált tér, \(Y\) teljes normált tér (Banach tér),
\(\left. A:M\rightarrow Y \right.\) korlátos lineáris operátor, ahol
\(M \subset X\) lineáris altér, de nem kell zártnak lennie. Ekkor az
\(A\) -nak egyértelműen létezik korlátos lineáris kiterjesztése az
\(\overline{M}\) -ra (\(M\) lezárására). Más szóval:
\(\left. \exists!\widetilde{A}:\overline{M}\rightarrow Y \right.\)
korlátos lineáris operátor, amelyre
\(\widetilde{A}x = Ax\quad\forall x \in M\). Spec eset, mikor
\(\overline{M} = X\).

\end{tetel}

\begin{bizonyitas}

Bizonyítás (vázlatos):\\
Legyen \(x \in \overline{M}\). Ehhez
\(\exists x_{k} \in M:\lim\left( x_{k} \right) = x\). Tekintsük az
\(\left( {Ax_{k}} \right)_{k \in {\mathbb{N}}}\) sorozatot \(Y\) -ban!
Belátjuk, hogy ez Cauchy sorozat.
\(\left\| {Ax_{k} - Ax_{l}} \right\|_{Y} = \left\| {A\left( {x_{k} - x_{l}} \right)} \right\|_{Y} \leq c \cdot \left\| {x_{k} - x_{l}} \right\|{}_{X}\).
Legyen \(\varepsilon > 0\) , \(\exists k_{0}:\forall k,l > k_{0}\)
esetén
\(\left. \left\| {x_{k} - x_{l}} \right\| < \varepsilon\Rightarrow\left\| {Ax_{k} - Ax_{l}} \right\| \leq c \cdot \varepsilon \right.\).
\(Y\) teljes
\(\left. \Rightarrow\exists y \in Y:\lim\left( {Ax_{k}} \right) = y \right.\).
\(y\) csak \(x\) -től függ, nem függ \(\left( x_{k} \right)\) -tól és
egyértelmű. \(\widetilde{A}\left( x \right): = y\) , \(\widetilde{A}\)
lineáris, korlátos (és folytonos).

\end{bizonyitas}

\begin{tetel}

Tétel (Hahn-Banach-tétel):\\
Legyen \(X\) Banach tér, \(X_{0} \subset X\) valódi (zárt lineáris)
altér, \(\left. f:X_{0}\rightarrow{\mathbb{K}} \right.\) korlátos
lineáris funkcionál (azaz számértékű operátor). Ekkor
\(\left. \exists\widetilde{f}:X\rightarrow{\mathbb{K}} \right.\)
korlátos lineáris kiterjesztés, és
\(\left\| \widetilde{f} \right\| = \left\| f \right\|\).

\end{tetel}

\hypertarget{korlatos-linearis-funkcionalok-dualis-ter-hilbert-ter-eseten}{%
\subsubsection{Korlátos lineáris funkcionálok, duális tér (Hilbert tér
esetén)}\label{korlatos-linearis-funkcionalok-dualis-ter-hilbert-ter-eseten}}

\textbf{Észrevétel}: legyen \(X\) Hilbert tér, \(y \in X\) tetszőleges
rögzített elem. Értelmezzük az
\(\left. f:X\rightarrow{\mathbb{K}},f\left( x \right): = \left\langle {x,y} \right\rangle \right.\)
funkcionált.

\begin{allitas}

Állítás:\\
Ekkor \(f\) korlátos lineáris funkcionál. \(f\) linearitása triviális,
és korlátos is, ugyanis
\(\left| {f\left( x \right)} \right| = \left| \left\langle {x,y} \right\rangle \right| \leq \left\| x \right\| \cdot \left\| y \right\|\).

\end{allitas}

\begin{tetel}

Tétel (Riesz-tétel):\\
Legyen \(X\) Hilbert tér (valós vagy komplex), \(f\) egy korlátos
lineáris funkcionál \(X\) -en. Ekkor létezik egyetlen \(y \in X\) , hogy
\(f\left( x \right) = \left\langle {x,y} \right\rangle\,\forall x \in X\).

\end{tetel}

\begin{bizonyitas}

Bizonyítás:\\
Jelölje \(X_{0}: = \left\{ {x \in X:f\left( x \right) = 0} \right\}\)
-vel \(f\) magterét. \(X_{0}\) altér \(X\) -ben, azaz az algebrai
műveletek nem vezetnek ki \(X_{0}\) -ból, és zárt részhalmaz \(X\) -ben.
Utóbbi azért igaz, mivel \(f\) folytonos, azaz ha \(x_{k} \in X_{0}\) ,
\(\left. \left( x_{k} \right)\rightarrow x\Rightarrow x \in X_{0} \right.\).
\(\left. f\left( x_{k} \right)\rightarrow f\left( x \right)\Rightarrow f\left( x \right) = 0 \right.\),
mivel jelen esetben \(f\left( x_{k} \right) = 0\).

\begin{enumerate}
\def\labelenumi{\arabic{enumi}.}
\tightlist
\item
  Ha \(X_{0} = X\) , \(f\left( x \right) = 0\,\forall x \in X\) ,
  triviális eset. Ekkor legyen \(y = 0\).
\item
  \(X_{0}\) valódi altér \(\Rightarrow\) (Riesz-féle felbontási tétel
  szerint) \(\exists x_{1} \neq 0:x_{1} \in X_{0}^{\bot}\). Legyen
  \(x \in X\) tetszőleges, és tekintsük az
  \(X \ni y_{1}: = f\left( x \right)x_{1} - f\left( x_{1} \right)x\)
  elemet. Ekkor
  \[\left. f\left( y_{1} \right) = f\left( x \right)f\left( x_{1} \right) - f\left( x_{1} \right)f\left( x \right) = 0\Rightarrow y_{1} \in X_{0}\Rightarrow\left\langle {y_{1},x_{1}} \right\rangle = 0. \right.\]
  Más szóval
  \[0 = \left\langle {y_{1},x_{1}} \right\rangle = \left\langle {f\left( x \right)x_{1} - f\left( x_{1} \right)x,x_{1}} \right\rangle = f\left( x \right)\left\| x_{1} \right\|^{2} - f\left( x_{1} \right)\left\langle {x,x_{1}} \right\rangle.\]
  Átrendezve kapjuk, hogy
  \(\left. f\left( x \right) = \frac{f\left( x_{1} \right)\left\langle {x,x_{1}} \right\rangle}{\left\| x_{1} \right\|^{2}} = \left\langle {x,\frac{\overline{f\left( x_{1} \right)}x_{1}}{\left\| x_{1} \right\|^{2}}} \right\rangle\Rightarrow\exists y \right.\)
  , nevezetesen
  \(y = \frac{\overline{f\left( x_{1} \right)}}{\left\| x_{1} \right\|^{2}}x_{1}\).
\item
  \(y\) egyértelmű. Tfh
  \[\left. \left\langle {x,y} \right\rangle = \left\langle {x,y^{*}} \right\rangle\,\forall x \in X\Rightarrow\left\langle {x,y - y^{*}} \right\rangle = 0\forall x \in X\Rightarrow y - y^{*} = 0\Rightarrow y = y^{*}. \right.\]
\end{enumerate}

\end{bizonyitas}

\begin{ajanlo}

\begin{ajanlofig}

\href{https://xkcd.com}{\includegraphics[width=5.20833in,height=2.82292in]{wikipedian_protester.png}}

\end{ajanlofig}

Text

\end{ajanlo}

\hypertarget{korlatos-linearis-funkcionalok}{%
\paragraph{Korlátos lineáris
funkcionálok}\label{korlatos-linearis-funkcionalok}}

Legyen \(X\) Hilbert tér \(y \in X\) egy rögzített eleme,
\(f\left( x \right): = \left\langle {x,y} \right\rangle\). Ekkor a
CS-ből következik: \(\left\| f \right\| \leq \left\| y \right\|\).

\begin{megjegyzes}

Megjegyzés:\\
\(\left\| f \right\| = \left\| y \right\|\) , ugyanis egyrészt
\(\left. \left| {f\left( x \right)} \right| = \left| \left\langle {x,y} \right\rangle \right| \leq \left\| x \right\| \cdot \left\| y \right\|\Rightarrow\left\| f \right\| \leq \left\| y \right\| \right.\).
Másrészt
\(\left\| f \right\| = \sup\left\{ {\left| {f\left( x \right)} \right|:\left\| x \right\| = 1} \right\}\).
Válasszuk \(x: = \frac{y}{\left\| y \right\|}\) (\(y \neq 0\) ,
máskülönben triviális), ekkor \(\left\| x \right\| = 1\) ,
\(\left| {f\left( x \right)} \right| = \left| \left\langle {\frac{y}{\left\| y \right\|},y} \right\rangle \right| = \left\| y \right\|\).
Tehát \(\left\| f \right\| = \left\| y \right\|\).

\end{megjegyzes}

Spec eset: \(X: = L^{2}\left( M \right)\) ,
\(M \subset {\mathbb{R}}^{n}\) mérhető halmaz. Ekkor egy tetszőleges
\(f\) korlátos lineáris funkcionál ilyen alakú:
\(f\left( \phi \right): = \left\langle {\phi,\psi} \right\rangle = {\int_{M}{\phi\overline{\psi}}}\)
, ahol \(\psi \in L^{2}\left( M \right)\) rögzített.
\(\psi_{0}: = \overline{\psi} \in L^{2}\left( M \right)\) jelöléssel
\(f\left( \phi \right) = {\int_{M}{\phi\psi_{0}}},\forall\phi \in L^{2}\left( M \right)\).

\textbf{Korlátos lineáris funkcionálok} \(L^{p}\left( M \right)\)
\textbf{-en, ahol} \(1 < p < \infty\) (azaz
\(L^{\infty}\left( M \right)\) teret nem tárgyaljuk)

Legyen \(\psi \in L^{q}\left( M \right)\) tetszőleges rögzített,
\(\frac{1}{p} + \frac{1}{q} = 1\)! Értelmezzük az \(f\) funkcionált:
\(f\left( \phi \right): = {\int_{M}{\phi\psi}}\), ahol
\(\phi \in L^{p}\left( M \right)\).

\begin{allitas}

Állítás:\\
\(f\) korlátos lineáris funkcionál \(L^{p}\left( M \right)\)-en.

\end{allitas}

\begin{bizonyitas}

Bizonyítás:\\
Tudjuk, hogy
\(\left. \phi \in L^{p}\left( M \right),\psi \in L^{q}\left( M \right)\Rightarrow\phi\psi \in L^{1}\left( M \right) \right.\),
tehát a funkcionál értelmezve van az egész \(L^{p}\left( M \right)\)-n,
nyilván lineáris. A Hölder egyenlőtlenség szerint
\(\left. \left| {\int_{M}{\phi\psi}} \right| \leq \left\| \phi \right\|_{L^{p}{(M)}} \cdot \left\| \psi \right\|_{L^{q}{(M)}}\Rightarrow\left\| f \right\| \leq \left\| \psi \right\|_{L^{q}{(M)}} \right.\),
vagyis korlátos is és normája
\(\leq \left\| \psi \right\|_{L^{q}{(M)}}\)

\end{bizonyitas}

\begin{tetel}

Tétel:\\
\(\left\| f \right\| = \left\| \psi \right\|_{L^{q}{(M)}}\).

\end{tetel}

\begin{tetel}

Tétel:\\
Legyen \(1 < p < \infty\). Ekkor tetszőleges
\(\left. f:L^{p}\left( M \right)\rightarrow{\mathbb{K}} \right.\)
korlátos lineáris funkcionálhoz
\(\exists!\psi \in L^{q}\left( M \right):f\left( \phi \right) = {\int_{M}{\psi\phi}}\).

\end{tetel}

\hypertarget{dualis-konjugalt-ter}{%
\paragraph{Duális (konjugált) tér}\label{dualis-konjugalt-ter}}

\begin{definicio}

Definíció:\\
Legyen \(X\) normált tér! Az \(X\)-en értelmezett korlátos lineáris
funkcionálok terét \(X\) duálisának nevezzük és \(X'\)-vel jelöljük
(van, ahol *-gal jelölik).

\end{definicio}

\begin{megjegyzes}

Megjegyzés:\\
\(X' = L\left( {X,{\mathbb{K}}} \right)\). Tudjuk, hogy
\(X' = L\left( {X,{\mathbb{K}}} \right)\) normált tér (norma az operátor
normája), \(X'\) tér teljes, mivel \(\mathbb{K}\) alaptest teljes, így
\(X'\) Banach tér.

\end{megjegyzes}

Értelmezzük az előbbieket ezen fogalom rögzítésével!

\(X\) Hilbert tér. Tudjuk, hogy
\(\forall f \in X'\exists y \in X:f\left( x \right) = \left\langle {x,y} \right\rangle\),
\(\left\| f \right\| = \left\| y \right\|\). Fordítva, \(y \in X\)
esetén
\(f\left( x \right): = \left\langle {x,y} \right\rangle,x \in X\)! Tehát
ha \(X\) Hilbert tér, bijekció létesíthető \(X'\) és \(X\) között.
Jelöljük: \(\Phi\left( y \right): = f\),
\(f\left( x \right): = \left\langle {x,y} \right\rangle\).
\(\left. \Phi:X\rightarrow X' \right.\) bijekció. Ennek tulajdonságai:

\begin{itemize}
\tightlist
\item
  \(\Phi\left( {y_{1} + y_{2}} \right) = \Phi\left( y_{1} \right) + \Phi\left( y_{2} \right)\).
  \(f_{1}\left( x \right) = \left\langle {x,y_{1}} \right\rangle\),
  \(f_{2}\left( x \right) = \left\langle {x,y_{2}} \right\rangle\).
  \(\left( {f_{1} + f_{2}} \right)\left( x \right) = f_{1}\left( x \right) + f_{2}\left( x \right) = \left\langle {x,y_{1}} \right\rangle + \left\langle {x,y_{2}} \right\rangle = \left\langle {x,y_{1} + y_{2}} \right\rangle\),
  vagyis \(\left. f_{1} + f_{2}\leftrightarrow y_{1} + y_{2} \right.\).
\item
  \(\lambda \in {\mathbb{K}}\) esetén
  \(\Phi\left( {\lambda y} \right) = \overline{\lambda}\Phi\left( y \right)\).
  \[\left. f\left( x \right) = \left\langle {x,y} \right\rangle\Rightarrow\left\langle {x,\lambda y} \right\rangle = \overline{\lambda}\left\langle {x,y} \right\rangle = \overline{\lambda}f\left( x \right) = \left( {\overline{\lambda}f} \right)x, \right.\]vagyis
  \(\left. \lambda y\leftrightarrow\overline{\lambda}f \right.\), tehát
  \(\Phi\) konjugált lineáris.
\end{itemize}

\(X = L^{p}\left( M \right)\) esete, mikor \(1 \leq p \leq \infty\) és
\(\frac{1}{p} + \frac{1}{q} = 1\).

Tudjuk, hogy tetszőleges \(\psi \in L^{q}\left( M \right)\) esetén
\(f\left( \phi \right): = {\int_{M}{\phi\psi}}\), \(\phi \in X\) mellett
\(f \in \left( {L^{p}\left( M \right)} \right)'\),
\(\left\| f \right\| = \left\| \psi \right\|\). Továbbá
\(\left( {L^{p}\left( M \right)} \right)'\) minden eleme ilyen alakú
\(p < \infty\) esetén.

\(\left. L^{q}\left( M \right) \ni \psi\leftrightarrow f \in \left( {L^{p}\left( M \right)} \right)' \right.\).
Könnyen belátható, hogy az eddigiek alapján \(\Phi\) bijekció, sőt,
\(\Phi\) lineáris. \(L^{p}\left( M \right)\) izomorf és izometrikus
(normatartó) \(L^{q}\left( M \right)\)-vel, ha \(p < \infty\).

\hypertarget{x-ter-mas-szoval-bidualis-reflexiv-ter}{%
\paragraph{\texorpdfstring{\(X''\) tér, más szóval biduális, reflexív
tér}{X'' tér, más szóval biduális, reflexív tér}}\label{x-ter-mas-szoval-bidualis-reflexiv-ter}}

\begin{definicio}

Definíció:\\
Legyen \(X\) normált tér. Ekkor definíció szerint
\(X": = \left( {X'} \right)'\).

\end{definicio}

\begin{allitas}

Állítás:\\
Ha \(X\) Hilbert tér, akkor \(X''\) izomorf, izometrikus az \(X\)
térrel.

\end{allitas}

\begin{definicio}

Definíció:\\
Legyen \(X\) Banach tér! Ha \(X''\) izomorf és izometrikus \(X\)-szel,
akkor \(X''\)-t reflexívnek nevezzük.

\end{definicio}

\begin{allitas}

Állítás:\\
Legyen \(X = L^{p}\left( M \right)\), ahol \(1 < p < \infty\)! Ekkor
\(L^{p}\left( M \right)\) reflexív.

\end{allitas}

Vizsgáljuk \(X''\)-t általános esetben, mikor \(X\) Banach tér!
Tekintsük egy tetszőleges, rögzített \(x \in X\) elemet, ehhez rendeljük
hozzá a következő, \(F_{x} \in X"\) elemet!
\(F_{x}\left( f \right): = f\left( x \right)\), \(\forall f \in X'\).
Ekkor \(F_{X}\) jól definiált funkcionál \(X'\)-n, nyilván lineáris,
korlátos is.
\(\left| {F_{x}\left( f \right)} \right| = \left| {f\left( x \right)} \right| \leq \left\| f \right\| \cdot \left\| x \right\|_{X}\),
\(\forall f \in X'\).
\(\left. \Rightarrow\left\| F_{x} \right\| \leq \left\| x \right\| \right.\).

\begin{allitas}

Állítás:\\
\(\left\| F_{x} \right\| = \left\| x \right\|\).

\end{allitas}

\begin{bizonyitas}

Bizonyítás:\\
(Definíció szerint
\(\left\| F_{x} \right\| = \sup\limits_{f \in X'}\left\{ {\left| {F_{x}\left( f \right)} \right| = \left| {f\left( x \right)} \right|:\left\| f \right\| = 1} \right\}\).)
Azt kellene belátni, hogy \(\exists f \in X':\left\| f \right\| = 1\),
melyre igaz, hogy
\(\left| {F_{x}\left( f \right)} \right| = \left\| x \right\|\) bármely
rögzített \(x\) esetén. Tekintsük a következő \(f_{0}\) funkcionált
\(X\) következő, 1 dimenziós alterén:
\(X_{0}: = \left\{ {\lambda x:\lambda \in {\mathbb{K}}} \right\}\), ahol
\(x \in X\) rögzített. Legyen
\(f_{0}\left( {\lambda x} \right): = \lambda\left\| x \right\|\).
\(f_{0}\) korlátos is,
\(\left. \left| {f_{0}\left( {\lambda x} \right)} \right| = \left| \lambda \right|\left\| x \right\| = \left\| {\lambda x} \right\| \cdot 1\Rightarrow\left\| f_{0} \right\| = 1 \right.\).
A Hahn-Banach tétel szerint az \(X_{0}\) altéren definiált \(f_{0}\)
korlátos lineáris funkcionál kiterjeszthető a korlátosság és linearitás
megtartásával az egész \(X\) térre úgy, hogy
\(\left\| f \right\| = \left\| f_{0} \right\|\) (ezt persze nem
bizonyítottuk). Jelölje ezt \(f\)! \(f \in X'\),
\(\left\| f \right\| = \left\| f_{0} \right\| = 1\). Erre
\(\left| {F_{x}\left( f \right)} \right| = \left| {f\left( x \right)} \right| = \left| {f\left( {1 \cdot x} \right)} \right| = f_{0}\left( {1 \cdot x} \right) = 1 \cdot \left\| x \right\| = \left\| x \right\|\).

\end{bizonyitas}

Általános esetben \(X''\) egy részhalmaza izomorf és izometrikus
\(X\)-szel. \(X''\)-nek lehetnek más elemei is (ha nem reflexív).

\hypertarget{gyenge-konvergencia}{%
\subsubsection{Gyenge konvergencia}\label{gyenge-konvergencia}}

\begin{definicio}

Definíció:\\
Legyenek \(X\), \(Y\) normált terek, és tfh
\(A_{j} \in L\left( {X,Y} \right),j \in {\mathbb{N}}\) (\(A_{j}\)
korlátos lineáris operátor \(X\)-n). Azt mondjuk, hogy ez az \(A_{j}\)
sorozat gyengén konvergál az \(A\) operátorhoz, ha \(\forall x \in X\)
elemre
\(\left. \left( {A_{j}x} \right)_{j \in {\mathbb{N}}}\rightarrow Ax \right.\)
(pontonkénti konvergencia). (\(Y\)-beli norma szerinti konvergencia).

\end{definicio}

\begin{allitas}

Állítás:\\
Ha \(\lim\left\| {A_{j} - A} \right\| = 0\), azaz
\(\left. \left( A_{j} \right)\rightarrow A \right.\) az
\(L\left( {X,Y} \right)\) norma szerint, akkor
\(\left. \left( A_{j} \right)\rightarrow A \right.\) gyengén, de
fordítva nem mindig igaz.

\end{allitas}

\begin{bizonyitas}

Bizonyítás:\\
Tfh \(\lim\left\| {A_{j} - A} \right\| = 0\). Ekkor
\[{\left\| {{A_j}x - Ax} \right\|_Y} = \left\| {\left( {{A_j} - A} \right)x} \right\| \leqslant \underbrace {\left\| {{A_j} - A} \right\|}_{ \to 0} \cdot \left\| x \right\| \to 0 .\]

\end{bizonyitas}

Speciális eset: \(Y = {\mathbb{K}}\), \(L\left( {X,Y} \right) = X'\).
\(\left. \left( f_{j} \right)\rightarrow f \right.\) gyengén \(X'\)-ben,
ha bármely rögzített \(x \in X\) esetén
\(\left. \left( {f_{j}\left( x \right)} \right)\rightarrow f\left( x \right) \right.\).

\begin{pelda}

Példa:\\
\(X'\)-beli gyengén konvergens sorozatra, amely norma szerint nem
konvergens. Legyen \(X\) szeparábilis, végtelen dimenziós Hilbert tér!
Legyen ebben egy \(y_{1},y_{2},...,y_{j},...\) ortonormált, teljes
rendszer!
\(f_{j}\left( x \right): = \left\langle {x,y_{j}} \right\rangle\). Ekkor
\(\left\langle {x,y_{j}} \right\rangle\) az \(x \in X\) elem \(j\)-edik
Fourier-egyeütthatója \(y_{j}\) ortonormált rendszer szerint,
\(c_{j}: = \left\langle {x,y_{j}} \right\rangle\). Tudjuk, hogy
\(\left. {\sum\limits_{j = 1}^{\infty}\left| c_{j} \right|^{2}} < \infty\Rightarrow\lim\left( c_{j} \right) = 0 \right.\),
azaz \(\lim\limits_{j\rightarrow\infty}f_{j}\left( x \right) = 0\),
\(\forall x \in X\). Más szóval \(\left( f_{j} \right)\) \(X\)'-beli
sorozat gyengén tart \(f = 0\) funkcionálhoz. Másrészt
\(\left\| f_{j} \right\| = \left\| y_{j} \right\|_{X} = 1\), így
\(\left( f_{j} \right)\) nem tart a norma szerint az \(f = 0\)
funkcionálhoz. (Bebizonyítható, hogy véges dimenzióban a gyenge
konvergencia egybeesik a norma szerinti konvergencia fogalmával.)

\end{pelda}

\begin{tetel}

Tétel:\\
Tfh \(A_{j} \in L\left( {X,Y} \right)\), ahol \(X\), \(Y\) Banach terek,
\(\left. \left( A_{j} \right)\rightarrow A \right.\) gyengén. Ekkor
\(\left( \left\| A_{j} \right\| \right)_{j \in {\mathbb{N}}}\) korlátos.
Ez a tétel következik az alábbi tételből.

\end{tetel}

\begin{tetel}

Egyenletes korlátosság tétele (Banach-Steinhaus-tétel, bizonyítás
nélkül):\\
Legyenek \(X\), \(Y\) Banach terek, \(A_{j} \in L\left( {X,Y} \right)\).
Ha az \(A_{j}\) operátor sorozat pontonként korlátos, azaz ha
\(\forall x \in X\) esetén
\(\left. \sup\limits_{j \in {\mathbb{N}}}\left\{ \left\| {A_{j}x} \right\| \right\} < \infty\Rightarrow\left( \left\| A_{j} \right\| \right) \right.\)
korlátos.

\end{tetel}

\begin{megjegyzes}

Megjegyzés (gyenge kompaktsági kritérium):\\
Tekintsük a \(X' = L\left( {X,{\mathbb{K}}} \right)\) speciális esetet
az egyszerűség kedvéért. Ha \(f_{j} \in X'\) korlátos sorozatot alkot
(\(X\) most Banach tér), akkor \(\left( f_{j} \right)\)-ból
kiválasztható egy gyengén konvergens részsorozat.

\end{megjegyzes}

\begin{ajanlo}

\begin{ajanlofig}

\href{https://xkcd.com}{\includegraphics[width=5.20833in,height=2.82292in]{wikipedian_protester.png}}

\end{ajanlofig}

Text

\end{ajanlo}

\hypertarget{gyenge-konvergencia-x-ben}{%
\subsubsection{\texorpdfstring{Gyenge konvergencia
\(X\)-ben}{Gyenge konvergencia X-ben}}\label{gyenge-konvergencia-x-ben}}

\begin{definicio}

Definíció:\\
Legyen \(X\) normált tér! Azt mondjuk, hogy egy
\(\left( x_{j} \right)_{j \in {\mathbb{N}}}\) \(X\)-beli sorozat gyengén
konvergál egy \(x \in X\) ponthoz, ha \(\forall f \in X'\) funkcionálra
\(\left. \left( {f\left( x_{j} \right)} \right)_{j \in {\mathbb{N}}}\rightarrow f\left( x \right) \right.\).

\end{definicio}

\begin{megjegyzes}

Megjegyzés:\\
Ha \(X\) reflexív Banach-tér, akkor minden korlátos \(X\)-beli
sorozatnak létezik gyengén konvergens részsorozata. Ugyanis ekkor
\(X = X'' = \left( {X'} \right)'\).

\end{megjegyzes}

\hypertarget{inverz-operator}{%
\paragraph{Inverz operátor}\label{inverz-operator}}

Emlékeztető: egy függvénynek létezik inverze, ha injektív. Tudjuk
továbbá, hogy egy \(\left. A:X\rightarrow Y \right.\) lineáris
operátornak létezik inverze (azaz injektív) \(\Leftrightarrow\) a magtér
csak a 0-ból áll, azaz
\(\left. Ax = 0_{Y}\Leftrightarrow x = 0_{X} \right.\). Továbbá, ha
\(A^{- 1}\) létezik, akkor \(A^{- 1}\) lineáris operátor. Egy \(A\)
operátor folytonos \(x_{0}\)-ban, ha
\(\left. \forall\varepsilon > 0\exists\rho > 0:\left\| {x - x_{0}} \right\|_{X} < \rho\Rightarrow\left\| {Ax - Ax_{0}} \right\|_{Y} < \varepsilon \right.\).

\textbf{Kérdés}: ha \(X\), \(Y\) normált terek,
\(\left. A:X\rightarrow Y \right.\) lineáris és injektív
\(\overset{?}{\Rightarrow}A^{- 1}\) korlátos is? Általában nem, akkor
sem, ha \(A\) korlátos.

\begin{tetel}

Nyílt leképezések tétele (bizonyítás nélkül):\\
Legyenek \(X\), \(Y\) Banach terek, \(\left. A:X\rightarrow Y \right.\)
korlátos lineáris operátor és \(R_{A} = Y\), vagyis ráképezés. Ekkor
\(A\) operátor \(X\) minden nyílt halmazát \(Y\) nyílt halmazába képezi.
Ebből következik:

\end{tetel}

\begin{tetel}

Tétel (Banach-tétel):\\
Legyenek \(X\), \(Y\) Banach terek, \(\left. A:X\rightarrow Y \right.\)
korlátos és lineáris, \(R_{A} = Y\) és \(A\) injektív! Ekkor \(A^{- 1}\)
korlátos (azaz folytonos).

\end{tetel}

\begin{bizonyitas}

Bizonyítás:\\
Legyen tetszőleges \(y_{0} \in Y = R_{A} = D_{A^{- 1}}\).
\(x_{0}: = A^{- 1}y_{0}\). Belátjuk, hogy az \(A^{- 1}\) folytonos
\(y_{0}\)-ban. Tekintsük \(x_{0} = A^{- 1}y_{0}\) egy tetszőleges
\(B_{r}\left( x_{0} \right)\) nyílt környezetét! Ennek képe is nyílt az
\(Y\)-ban az előbbi tétel szerint. Mivel
\(y_{0} \in A\left( {B_{r}\left( x_{0} \right)} \right)\), ami nyílt,
ezért \(y_{0}\)-nak van olyan környezete, melyre
\(B_{\rho}\left( y_{0} \right) \subset A\left( {B_{r}\left( x_{0} \right)} \right)\).
Ez azt jelenti, hogy ha
\(\left. y \in B_{\rho}\left( y_{0} \right)\Rightarrow A^{- 1}y \in B_{r}\left( x_{0} \right) \right.\).
Eszerint \(A^{- 1}\) folytonos \(y_{0}\)-ban.

\end{bizonyitas}

\hypertarget{zart-graf-grafikon-tetel}{%
\paragraph{Zárt gráf (grafikon) tétel}\label{zart-graf-grafikon-tetel}}

\begin{definicio}

Definíció:\\
Legyenek \(X\), \(Y\) normált terek, \(\left. A:M\rightarrow Y \right.\)
lineáris operátor, \(M \subset X\). Ekkor \(A\) operátor gráfja,
grafikonja az alábbi halmaz:
\(G_{A}: = \left\{ {\left( {x,Ax} \right):x \in M = D_{A}} \right\}\).

\end{definicio}

\begin{definicio}

Definíció:\\
Egy \(\left. A:M\rightarrow Y \right.\) lineáris operátort zártnak
nevezünk, ha a \(G_{A} \subset X \times Y\) zárt halmaz
\(X \times Y\)-ban.
\(X \times Y = \left\{ {\left( {x,y} \right):x \in X,y \in Y} \right\}\).

\end{definicio}

\begin{megjegyzes}

Megjegyzés:\\
A szorzattéren értelmezett műveletek:

\begin{itemize}
\tightlist
\item
  \(\left( {x_{1},y_{1}} \right) + \left( {x_{2},y_{2}} \right) = \left( {x_{1} + x_{2},y_{1} + y_{2}} \right)\)
\item
  \(\lambda\left( {x,y} \right) = \left( {\lambda x,\lambda y} \right)\)
\item
  \(\left\| \left( {x,y} \right) \right\|_{X \times Y}: = \sqrt{\left\| x \right\|^{2} + \left\| y \right\|^{2}}\),
  \(X \times Y\) normált tér tehát.
\end{itemize}

\end{megjegyzes}

Legyenek \(X\), \(Y\) normált terek, \(\left. A:M\rightarrow Y \right.\)
lineáris operátor, \(D_{A} = M \subset X\). \(A\) zárt
\(\Leftrightarrow\) ha minden
\(\left( x_{j} \right)_{j \in {\mathbb{N}}}\) \(M\)-beli sorozatra,
melyre \(\lim\left( x_{j} \right) = x \in X\) és
\(\exists\lim\left( {Ax_{j}} \right) = y \in Y\), akkor \(x \in M\) és
\(y = Ax\). Ezért ha \(A\) folytonos, akkor zárt is.

\begin{pelda}

Példa:\\
Példa zárt, lineáris, de nem folytonos (nem korlátos) operátorra:
\(X: = C\left\lbrack 0,1 \right\rbrack\),
\(M = D_{A} = C^{1}\left\lbrack 0,1 \right\rbrack\),
\(A\phi: = \phi^{\prime}\), vagyis a differenciáloperátor.
\(\left. \left( \phi_{j} \right)\rightarrow\phi \right.\) egyenletesen
(\(C\left\lbrack 0,1 \right\rbrack\)-beli konvergencia) és
\(\left. \left( {\phi^{\prime}}_{j} \right)\rightarrow\psi \right.\)
egyenletesen \(\left. \Rightarrow\psi = \phi^{\prime} \right.\), tehát
\(A\) valóban zárt, lineáris (de nem korlátos, így nem is folytonos, ezt
láttuk korábban).

\end{pelda}

\begin{tetel}

Tétel (zárt gráf tétel):\\
Legyenek \(X\), \(Y\) Banach terek, \(\left. A:X\rightarrow Y \right.\)
zárt, lineáris operátor (tehát \(D_{A} = X\)). Ekkor \(A\) folytonos
(korlátos).

\end{tetel}

\begin{bizonyitas}

Bizonyítás:\\
\(G_{A}: = \left\{ {\left( {x,Ax} \right):x \in D_{A} = X} \right\} \subset X \times Y\)
(utóbbi Banach-tér), ugyanis \(G_{A}\) zárt halmaz \(X \times Y\)-ban,
az \(X \times Y\) vektortenérnek altere:
\(\left( {x_{1},Ax_{1}} \right) + \left( {x_{2},Ax_{2}} \right) = \left( {x_{1} + x_{2},A\left( {x_{1} + x_{2}} \right)} \right) \in G_{A}\),
\(\lambda\left( {x,Ax} \right) = \left( {\lambda x,A\left( {\lambda x} \right)} \right) \in G_{A}\).
\(G_{A}\) az \(X \times Y\) Banach tér zárt lineáris altere
\(\left. \Rightarrow G_{A} \right.\) Banach-tér. Tekintsük a következő
két operátort: \(U\left( {x,Ax} \right): = x\),
\(V\left( {x,Ax} \right): = Ax\), ahol
\(\left( {x,Ax} \right) \in G_{A}\). Ekkor
\(\left. U:G_{A}\rightarrow X \right.\), \(R_{U} = X\),
\(\left. V:G_{A}\rightarrow Y \right.\). Most \(U\)-ra alkalmazható a
Banach tétel (az inverz operátor korlátosságáról): \(D_{U} = G_{A}\),
\(R_{U} = X\), \(U\) korlátos és injektív
\(\left. \Rightarrow U^{- 1}:X\rightarrow G_{A} \right.\) korlátos
(folytonos), \(A = VU^{- 1}\), mert
\(U^{- 1}x = \left( {x,Ax} \right)\),
\(V\left( {U^{- 1}\left( x \right)} \right) = V\left( {x,Ax} \right) = Ax\).
\(\left. V:G_{A}\rightarrow Y \right.\) korlátos
\(\left. \Rightarrow A = VU^{- 1} \right.\) is korlátos.

\end{bizonyitas}

\hypertarget{sajatertek-regularis-ertek-spektrum}{%
\subsection{Sajátérték, reguláris érték,
spektrum}\label{sajatertek-regularis-ertek-spektrum}}

Legyenek \(X\), \(Y\) normált terek, \(\left. A:M\rightarrow Y \right.\)
lineáris operátor, \(M \subset X\), \(b \in Y\) adott elem.

\begin{enumerate}
\def\labelenumi{\arabic{enumi}.}
\tightlist
\item
  Elsőfajú egyenlet: melyik az a \(x \in M = D_{A}:Ax = b\)?
\item
  Másodfajú egyenlet: legyen \(Y = X\). Melyik az a \(x \in X\), melyre
  \(\left( {\lambda I - A} \right)x = b\), ahol
  \(\lambda \in {\mathbb{K}}\), \(I\) az identitás. Ha
  \(\left( {\lambda I - A} \right)\) nem injektív, azaz nem létezik az
  inverzre, akkor \(\lambda\)-t az \(A\) operátor sajátértékének
  nevezzük. Ez azt jelenti, hogy
  \(\left. \exists x_{0} \neq 0:\left( {\lambda I - A} \right)x_{0} = 0\Leftrightarrow Ax_{0} = \lambda x_{0} \right.\).
\end{enumerate}

\begin{definicio}

Definíció:\\
Ha \(\exists\left( {\lambda I - A} \right)^{- 1}\), ez korlátos és
\(R_{\lambda I - A}\) értelmezési tartománya sűrű halmaz \(X\)-ben,
akkor \(\lambda\)-t reguláris értéknek nevezzük.

\end{definicio}

\begin{allitas}

Állítás:\\
Ha \(A\) zárt operátor, akkor reguláris érték esetén
\(D_{{({\lambda I - A})}^{- 1}} = X\), azaz \(R_{\lambda I - A} = X\).

\end{allitas}

\begin{megjegyzes}

Megjegyzés:\\
Ekkor reguláris értéke esetén \(\left( {\lambda I - A} \right)x = b\)
egyenletnek \(\forall b \in X\)-hez \(\exists!x\) megoldás, és \(x\)
folytonosan függ \(b\)-től, azaz
\(x = {\underbrace {\left( {\lambda I - A} \right)}_{{\text{folytonos}}}^{ - 1}}b\).

\end{megjegyzes}

\begin{definicio}

Definíció:\\
Az \(A\) operátor spektruma a reguláris értékek halmazának a
komplementere az alaptestben. A sajátértékek halmaza része a
spektrumnak.

\end{definicio}

\hypertarget{korlatos-linearis-operatorok-regularis-ertekei}{%
\subsubsection{Korlátos lineáris operátorok reguláris
értékei}\label{korlatos-linearis-operatorok-regularis-ertekei}}

\begin{tetel}

Tétel:\\
Legyen \(X\) Banach tér! Legyen \(\left. A:X\rightarrow X \right.\)
korlátos lineáris operátor. Ekkor
\(r_{\sigma}\left( A \right): = \lim\limits_{k\rightarrow\infty}\left\| A^{k} \right\|^{1/k}\),
ez létezik és véges. Ha \(\lambda \in {\mathbb{K}}\) számra teljesül,
hogy \(\left| \lambda \right| > r_{\sigma}\left( A \right)\), akkor
\(\lambda\) reguláris érték (\(A\)-ra nézve).

\end{tetel}

\begin{definicio}

Definíció:\\
\(r_{\sigma}\left( A \right)\) számot az \(A\) korlátos lineáris
operátor spektrálsugarának nevezzük.

\end{definicio}

\begin{megjegyzes}

Megjegyzések:

\begin{itemize}
\tightlist
\item
  \(A,B \in L\left( {X,X} \right)\) esetén
  \(\left\| {AB} \right\| \leq \left\| A \right\|\left\| B \right\|\),
  ugyanis
  \(\left\| {\left( {AB} \right)x} \right\| = \left\| {A\left( {Bx} \right)} \right\| \leq \left\| A \right\| \cdot \left\| {Bx} \right\| \leq \left\| A \right\| \cdot \left\| B \right\| \cdot \left\| x \right\|\)
  minden \(x\)-re,
  \(\left. \Rightarrow\left\| {AB} \right\| \leq \left\| A \right\|\left\| B \right\| \right.\)
\item
  \(\left\| A^{k} \right\| \leq \left\| A \right\|^{k}\).
  \(\left. \left\| A^{k} \right\|^{1/k} \leq \left( \left\| A \right\|^{k} \right)^{1/k} = \left\| A \right\|\Rightarrow r_{\sigma}\left( A \right) \leq \left\| A \right\| \right.\).
  Következmény: ha
  \(\left. \left| \lambda \right| > \left\| A \right\|\Rightarrow\lambda \right.\)
  reguláris érték.
\end{itemize}

\end{megjegyzes}

\begin{tetel}

Tétel (lemma 1):\\
Legyen \(Z\) Banach-tér, \(z_{k} \in Z\). Ha
\(\left. {\sum\limits_{k = 1}^{\infty}{\left\| z_{k} \right\| < \infty}}\Rightarrow{\sum\limits_{k = 1}^{\infty}z_{k}} \right.\)
konvergens \(Z\) Banach-téren.

\end{tetel}

\begin{bizonyitas}

Bizonyítás:\\
Legyen \(s_{j}: = {\sum\limits_{k = 1}^{j}z_{k}}\) részlet összeg!
\(\left\| {s_{j} - s_{l}} \right\| = \left\| {\sum\limits_{k = l + 1}^{j}z_{k}} \right\| \leq {\sum\limits_{k = l + 1}^{j}\left\| z_{k} \right\|} < \varepsilon\),
ha \(l,j > j_{0}\), tehát teljesül a Cauchy kritérium. Mivel \(Z\)
Banach-tér, azaz teljes normált tér, ezért minden Cauchy-sorozatnak van
határértéke \(Z\)-ben.

\end{bizonyitas}

\begin{tetel}

Tétel (lemma 2, bizonyítás nélkül):\\
Tfh \(B_{k} \in L\left( {X,X} \right)\),
\(\sum\limits_{k = 1}^{\infty}B_{k}\) konvergens
\(L\left( {X,X} \right)\)-en. Ekkor
\(\forall C \in L\left( {X,X} \right)\) operátorra
\(C{\sum\limits_{k = 1}^{\infty}B_{k}} = {\sum\limits_{k = 1}^{\infty}{CB_{k}}}\).
A bizonyítás egyszerű a részletösszegek segítségével.

\end{tetel}

\begin{ajanlo}

\begin{ajanlofig}

\href{https://xkcd.com}{\includegraphics[width=5.20833in,height=2.82292in]{wikipedian_protester.png}}

\end{ajanlofig}

Text

\end{ajanlo}

\begin{tetel}

Tétel:\\
Legyen \(X\) Banach-tér, \(\left. A:X\rightarrow X \right.\) korlátos,
lineáris operátor. Ekkor létezik és véges:
\(r_{\sigma}\left( A \right): = \lim\limits_{k\rightarrow\infty}\left\| A^{k} \right\|^{1/k}\).
Továbbá
\(\left. \left| \lambda \right| > r_{\sigma}\left( A \right)\Rightarrow\lambda \right.\)
reguláris érték,
\(\left( {\lambda I - A} \right)^{- 1} = \frac{1}{\lambda}\left( {I - \frac{1}{\lambda}A} \right)^{- 1} = \frac{1}{\lambda}{\sum\limits_{k = 0}^{\infty}{\frac{1}{\lambda^{k}}A^{k}}} = {\sum\limits_{k = 0}^{\infty}{\lambda^{- k - 1}A^{k}}}\).
Ez a sor -- a Neumann-sor -- \(L\left( {X,X} \right)\) normában
konvergens.

\end{tetel}

\begin{bizonyitas}

Bizonyítás:

\begin{enumerate}
\def\labelenumi{\arabic{enumi}.}
\tightlist
\item
  Jelöljük:
  \(r: = \inf\left\{ {\left\| A^{k} \right\|^{1/k}:k \in {\mathbb{N}}} \right\} \geq 0\),
  ez véges. Belátjuk, hogy
  \[r_{\sigma}\left( A \right) = \lim\limits_{k\rightarrow\infty}\left\| A^{k} \right\|^{1/k} = r = \inf\left\{ {\left\| A^{k} \right\|^{1/k}:k \in {\mathbb{N}}} \right\} \geq 0.\]
  Legyen \(\varepsilon > 0\) tetszőleges, ekkor az alsó határ
  definíciójából következik, hogy
  \(\exists m \in {\mathbb{N}}:r \leq \left\| A^{m} \right\|^{1/m} < r + \varepsilon\).
  Ezen \(m\) mellett válasszunk egy \(k > m\) számot, melyre
  \(k = pm + q\), ahol \(p \in {\mathbb{N}}\) és \(0 \leq q < m\) (ez
  \(k\)-nak \(m\)-vel vett maradékos osztása, \(q\) a maradéktag). Ekkor
  \(A^{k} = A^{pm + q} = \left( A^{p} \right)^{m} \cdot A^{q}\), így
  \[\left. \left\| A^{k} \right\| \leq \left\| A^{m} \right\|^{p} \cdot \left\| A \right\|^{q}\Rightarrow\left\| A^{k} \right\|^{1/k} \leq \left\| A^{m} \right\|^{p/k} \cdot \left\| A \right\|^{q/k} \leq \left( {r + \varepsilon} \right)^{mp/k}\left\| A \right\|^{q/k}. \right.\]
  Vegyük észre, hogy
  \(\lim\limits_{k\rightarrow\infty}\frac{mp}{k} = 1\), mert
  \(\lim\limits_{k\rightarrow\infty}\frac{q}{k} = 0\), így a fenti
  egyenlőtlenség jobb oldala
  \(\left. \rightarrow r + \varepsilon \right.\). Ebből következik, hogy
  \[\left. \exists k_{0}:k > k_{0}\Rightarrow r \leq \left\| A^{k} \right\|^{1/k} \leq r + 2\varepsilon\Rightarrow\lim\limits_{k\rightarrow\infty}\left\| A^{k} \right\|^{1/k} = r. \right.\]
\item
  Belátjuk, hogy a Neumann-sor \(L\left( {X,X} \right)\)-ben konvergens.
  Az 1. lemma szerint ehhez elég bizonyítani, hogy a sor tagjainak
  normáiból alkotott sor konvergens, azaz
  \({\sum\limits_{k = 0}^{\infty}\left\| {\lambda^{- k - 1}A^{k}} \right\|} < \infty\).
  Válasszunk egy olyan \(r_{1}\) számot, melyre
  \(\left| \lambda \right| > r_{1} > r_{\sigma}\left( A \right)\)! Mivel
  \(r_{\sigma}\left( A \right) = \lim\limits_{k\rightarrow\infty}\left\| A^{k} \right\|^{1/k}\)
  és \(r_{1} > r_{\sigma}\left( A \right)\), ezért
  \(\left. \exists k_{1} \in {\mathbb{N}}:k > k_{1}\Rightarrow r_{1} > \left\| A^{k} \right\|^{1/k} \right.\),
  így
  \(\left\| {\lambda^{- k - 1}A^{k}} \right\| = \frac{1}{\left| \lambda \right|^{k + 1}}\left\| A^{k} \right\| < \frac{1}{\left| \lambda \right|^{k + 1}}r_{1}^{k} = \frac{1}{\left| \lambda \right|}\left( \frac{r_{1}}{\left| \lambda \right|} \right)^{k}\).
  Ezeket összegezve \(k\) szerint egy mértani sort kapunk, melynek
  kvóciense \(0 < \frac{r_{1}}{\left| \lambda \right|} < 1\), így a sor
  konvergens, azaz
  \({\sum\limits_{k = 1}^{\infty}{\frac{1}{\left| \lambda \right|}\left( \frac{r_{k}}{\left| \lambda \right|} \right)^{k}}} < \infty\).
\item
  Jelöljük
  \(B: = {\sum\limits_{k = 0}^{\infty}{\lambda^{- k - 1}A^{k}}} \in L\left( {X,X} \right)\).
  Előbb láttuk, hogy ez konvergens. Ebből következni fog, hogy
  \(\left( {\lambda I - A} \right)^{- 1}\) létezik és egyenlő \(B\)-vel.
  A 2. lemmát felhasználva:
  \(\left( {\lambda I - A} \right)B = \lambda B - AB = \lambda{\sum\limits_{k = 0}^{\infty}{\lambda^{- k - 1}A^{k}}} - A{\sum\limits_{k = 0}^{\infty}{\lambda^{- k - 1}A^{k}}} = {\sum\limits_{k = 0}^{\infty}{\lambda^{- k}A^{k}}} - {\sum\limits_{k = 0}^{\infty}{\lambda^{- k - 1}A^{k + 1}}} = I\).
  Hasonlóképpen, \(B\left( {\lambda I - A} \right) = I\).
  Következtetésképpen \(\left( {\lambda I - A} \right)^{- 1}\) létezik
  és egyenlő \(B\)-vel.
\end{enumerate}

\end{bizonyitas}

\textbf{Következmény}:
\(\left| \lambda \right| > r_{\sigma}\left( A \right)\) esetén a
\(\left( {\lambda I - A} \right)x = b\) másodfajú egyenletnek létezik
egyetlen \(x\) megoldása, mégpedig
\[x = \left( {\lambda I - A} \right)^{- 1}b = \left( {{\sum\limits_{k = 0}^{\infty}\lambda^{- k - 1}}A^{k}} \right)b = {\sum\limits_{k = 0}^{\infty}{\left( {\lambda^{- k - 1}A^{k}} \right)b}} = {\sum\limits_{k = 0}^{\infty}{\lambda^{- k - 1}\left( {A^{k}b} \right)}},\]
ez a sor pedig \(X\) normában konvergens. A sor összege így is írható:
\(\frac{1}{\lambda}b + {\sum\limits_{k = 1}^{\infty}{\lambda^{- k - 1}A^{k}b}}\).
A fentiek még inkább érvényesek, ha
\(\left| \lambda \right| > \left\| A \right\|\).

\begin{allitas}

Állítás:\\
\(r_{\sigma}\left( A \right) = \sup\left\{ {\left| \lambda \right|:\lambda \in A_{\text{spektrum}}} \right\}\).

\end{allitas}

\hypertarget{peldak-alkalmazasok}{%
\subsubsection{Példák, alkalmazások}\label{peldak-alkalmazasok}}

\hypertarget{a-negyzetesen-integralhato-magu-integraloperatorok.}{%
\paragraph{A négyzetesen integrálható magú
integráloperátorok.}\label{a-negyzetesen-integralhato-magu-integraloperatorok.}}

Legyen \(M \subset {\mathbb{R}}^{n}\) egy Lebesgue szerint mérhető
halmaz, \(X: = L^{2}\left( M \right)\), ez ugye Hilbert tér. Legyen
\(\mathcal{K} \in L^{2}\left( {M \times M} \right)\) az úgynevezett
magfüggvény, s \(\phi \in L^{2}\left( M \right)\). Definiáljuk:
\(\psi\left( x \right): = {\int_{M}{\mathcal{K}\left( {x,y} \right)\phi\left( y \right)dy}}\).

\begin{allitas}

Állítás:\\
\(\psi \in L^{2}\left( M \right)\), továbbá a
\(K\left( \phi \right): = \psi\) képlettel értelmezett
\(\left. K:L^{2}\left( M \right)\rightarrow L^{2}\left( M \right) \right.\)
operátor lineáris, korlátos. A \(K\) operátort négyzetesen integrálató
magú integráloperátornak nevezzük.

\end{allitas}

\begin{bizonyitas}

Bizonyítás:\\
A Cauchy-Schwarz egyenlőtlenség szerint majdnem minden \(x\)-re
\[\left| {\psi\left( x \right)} \right| \leq {\int_{M}{\left| {\mathcal{K}\left( {x,y} \right)} \right| \cdot \left| {\phi\left( y \right)} \right|dy \leq \left\{ {\int_{M}{\left| {\mathcal{K}\left( {x,y} \right)} \right|^{2}dy}} \right\}^{1/2} \cdot \left\{ {\int_{M}{\left| {\phi\left( y \right)} \right|^{2}dy}} \right\}^{1/2}}}.\]
Mivel
\(\left. \mathcal{K} \in L^{2}\left( {M \times M} \right)\Rightarrow{\int_{M \times M}{\left| {\mathcal{K}\left( {x,y} \right)} \right|^{2}dxdy < \infty}} \right.\).
Fubini-tételt használva
\(\int_M {\underbrace {\int_M {{{\left| {\mathcal{K}\left( {x,y} \right)} \right|}^2}dy} }_{{\text{véges m.m. }}x{\text{-re}}} dx} < \infty\),
így
\[\left| {\psi\left( x \right)} \right|^{2} \leq {\int_{M}{\left| {\mathcal{K}\left( {x,y} \right)} \right|^{2}dy}} \cdot \left\lbrack {\int_{M}{\left| {\phi\left( y \right)} \right|^{2}dy}} \right\rbrack < \infty.\]Integrálva:
\[\left. {\int_{M}{\left| {\psi\left( x \right)} \right|^{2}dx}} \leq \left\lbrack {\int_{M}{{\int_{M}{\left| {\mathcal{K}\left( {x,y} \right)} \right|^{2}dy}}dx}} \right\rbrack \cdot \left\lbrack {\int_{M}{\left| {\phi\left( y \right)} \right|^{2}dy}} \right\rbrack < \infty\Rightarrow\psi \in L^{2}\left( M \right). \right.\]\(K\)
linearitása triviális. \(K\) korlátos, ugyanis
\[\left\| {K\phi} \right\|^{2}{}_{L^{2}{(M)}} = \left\| \psi \right\|^{2}{}_{L^{2}{(M)}} \leq \left\{ {\int_{M \times M}{\left| {\mathcal{K}\left( {x,y} \right)} \right|^{2}dxdy}} \right\} \cdot \left\| \phi \right\|^{2},\]
sőt:
\(\left\| K \right\| \leq \left\{ {\int_{M \times M}{\left| {\mathcal{K}\left( {x,y} \right)} \right|^{2}dxdy}} \right\}^{1/2} = \left\| \mathcal{K} \right\|_{L^{2}{({M \times M})}}\).

\end{bizonyitas}

\textbf{Következmény}:
\(\left| \lambda \right| > \left\| \mathcal{K} \right\|_{L^{2}{({M \times M})}}\)
esetén \(\lambda\) reguláris érték. Tudjuk, hogy
\(\left| \lambda \right| > r_{\sigma}\left( K \right)\) esetén
\(\lambda\) reguláris érték és
\(\left( {\lambda I - K} \right)^{- 1} = {\sum\limits_{k = 0}^{\infty}{\lambda^{- 1 - k}K^{k}}}\).

\textbf{Kérdés}: \(K\) integrál operátor hatványai hogyan számolhatók?

\begin{allitas}

Állítás:\\
Legyen \(\mathcal{K},\mathcal{L} \in L^{2}\left( {M \times M} \right)\)
és \(K,L\) a megfelelő integráloperátorok. Ekkor \(P: = KL\) szintén
négyzetesen integrálható magú operátor, amelynek magfüggvénye
\(\mathcal{P}\left( {x,y} \right): = {\int_{M}{\mathcal{K}\left( {x,t} \right)\mathcal{L}\left( {t,y} \right)dt}}\).

\end{allitas}

\begin{bizonyitas}

Bizonyítás:\\
\(\phi \in L^{2}\left( M \right)\) esetén \[\begin{aligned}
  \left( {P\phi } \right)\left( x \right) = \left[ {K\left( {L\phi } \right)} \right]\left( x \right) &  = \int_M {\mathcal{K}\left( {x,t} \right)\left[ {\int_M {\mathcal{L}\left( {t,y} \right)\phi \left( y \right)dy} } \right]dt}  \\ 
   &  = \int_M {\underbrace {\left[ {\int_M {\mathcal{K}\left( {x,t} \right)\mathcal{L}\left( {t,y} \right)dt} } \right]}_{\mathcal{P}\left( {x,y} \right)}\phi \left( y \right)dy} , \\ 
\end{aligned} \] ahol a Fubini-tételt ismét alkalmaztuk.
\(\mathcal{P} \in L^{2}\left( {M \times M} \right)\), merthogy
\(\left| {\mathcal{P}\left( {x,y} \right)} \right| \leq \left\{ {\int_{M}{\left| {\mathcal{K}\left( {x,t} \right)} \right|^{2}dt}} \right\}^{1/2}\left\{ {\int_{M}{\left| {\mathcal{L}\left( {t,y} \right)} \right|^{2}dy}} \right\}^{1/2}\),
így integrálva:
\[\int_{M \times M} {{{\left| {\mathcal{P}\left( {x,y} \right)} \right|}^2}dxdy}  \leqslant \underbrace {\int_M {\left[ {\int_M {{{\left| {\mathcal{K}\left( {x,t} \right)} \right|}^2}dt} } \right]dx} }_{ < \infty } \cdot \underbrace {\int_M {\left[ {\int_M {{{\left| {\mathcal{L}\left( {x,t} \right)} \right|}^2}dt} } \right]dy} }_{ < \infty } < \infty .\]

\end{bizonyitas}

\textbf{Következmény}:
\(\left( {K^{j}\phi} \right)\left( x \right) = {\int_{M}{\mathcal{K}_{j}\left( {x,y} \right)\phi\left( y \right)dy}}\),
\(j = 1,2,...\), ahol \[\begin{aligned}
  {\mathcal{K}_1}\left( {x,y} \right) & : = \mathcal{K}\left( {x,y} \right) \\ 
  {\mathcal{K}_2}\left( {x,y} \right) & : = \int_M {\mathcal{K}\left( {x,t} \right){\mathcal{K}_1}\left( {t,y} \right)dt}  \\ 
  {\mathcal{K}_j}\left( {x,y} \right) & : = \int_M {\mathcal{K}\left( {x,t} \right){\mathcal{K}_{j - 1}}\left( {t,y} \right)dt} . \\ 
\end{aligned} \] Ebből következik, hogy
\(\left( {\lambda I - K} \right)^{- 1}b = {\sum\limits_{j = 0}^{\infty}{\lambda^{- j - 1}K^{j}b}}\).
\[\begin{aligned}
  \left[ {{{\left( {\lambda I - K} \right)}^{ - 1}}b} \right]\left( x \right) &  = \left[ {\mathop \sum \limits_{j = 0}^\infty  {\lambda ^{ - j - 1}}{K^j}b} \right]\left( x \right) \\ 
   &  = \mathop \sum \limits_{j = 0}^\infty  {\lambda ^{ - j - 1}}\left( {{K^j}b} \right)\left( x \right) \\ 
   &  = \frac{{b\left( x \right)}}{\lambda } + \mathop \sum \limits_{j = 1}^\infty  {\lambda ^{ - j - 1}}\int_M {{\mathcal{K}_j}\left( {x,y} \right)b\left( y \right)dy}  \\ 
   & \mathop  = \limits^? \frac{{b\left( x \right)}}{\lambda } + \int_M {\underbrace {\left[ {\mathop \sum \limits_{j = 1}^\infty  {\lambda ^{ - j - 1}}{\mathcal{K}_j}\left( {x,y} \right)} \right]}_{ \in {L^2}\left( {M \times M} \right)}b\left( y \right)dy} . \\ 
\end{aligned} \] A sor \(L^{2}\left( M \right)\) normában konvergál. Az
utolsó egyenlőséget a következő órán látjuk be.

\begin{ajanlo}

\begin{ajanlofig}

\href{https://xkcd.com}{\includegraphics[width=5.20833in,height=2.82292in]{wikipedian_protester.png}}

\end{ajanlofig}

Text

\end{ajanlo}

A korábbiak szerint \(\left( {\lambda I - A} \right)x = b\) egyenletnek
van egyértelmű megoldása \(x\)-re és
\(x = {\sum\limits_{k = 0}^{\infty}{\lambda^{- k - 1}\left( {A^{k}b} \right)}}\),
ha \(\lambda\) reguláris érték, ugyanis ekkor a jobb oldal konvergens
\(X \ni x\)-ben.

Az előző példában \(X: = L^{2}\left( M \right)\) volt, (ahol
\(M \subset {\mathbb{R}}^{n}\) mérhető halmaz,) illetve
\(\mathcal{K} \in L^{2}\left( {M \times M} \right)\) , és
\[\psi\left( x \right): = \left( {K\phi} \right)\left( x \right) = {\int_{M}{\mathcal{K}\left( {x,y} \right)\phi\left( y \right)dy}},\]
ahol
\(\left. K:L^{2}\left( M \right)\rightarrow L^{2}\left( M \right) \right.\)
korlátos lineáris operátor és
\(r_{\sigma}\left( K \right) \leq \left\| K \right\| \leq \left\| \mathcal{K} \right\|_{L^{2}{({M \times M})}}\).

\(\left( {\lambda I - K} \right)\phi = b\),
\(b \in L^{2}\left( M \right)\) adott esetén mi a megoldás
\(\phi \in L^{2}\left( M \right)\)-re? Az egyenlet ekvivalens:
\(\lambda\phi\left( x \right) - {\int_{M}{\mathcal{K}\left( {x,y} \right)\phi\left( y \right)dy}} = b\left( x \right)\)
majdnem minden \(x \in M\)-re. Ha
\(\left. \left| \lambda \right| > r_{\sigma}\left( K \right)\Rightarrow\phi = {\sum\limits_{j = 0}^{\infty}{\lambda^{- j - 1}K^{j}b}} = \frac{b}{\lambda} + {\sum\limits_{j = 1}^{\infty}{\lambda^{- j - 1}K^{j}b}} \right.\).
\(\left( {K^{j}b} \right)\left( x \right) = {\int_{M}{\mathcal{K}_{j}\left( {x,y} \right)b\left( y \right)dy}}\),
\(\mathcal{K}_{j}\left( {x,y} \right) = {\int_{M}{\mathcal{K}_{j - 1}\left( {x,t} \right)\mathcal{K}\left( {t,y} \right)dt}}\)
és \(\mathcal{K}_{1} = \mathcal{K}\). Így
\[\phi \left( x \right) = \frac{{b\left( x \right)}}{x} + \mathop \sum \limits_{j = 1}^\infty  {\lambda ^{ - j - 1}}\mathop \smallint \limits_M {\mathcal{K}_j}\left( {x,y} \right)b\left( y \right)dy = \frac{{b\left( x \right)}}{\lambda } + \int_M {\underbrace {\left[ {\mathop \sum \limits_{j = 1}^\infty  {\lambda ^{ - j - 1}}{\mathcal{K}_j}\left( {x,y} \right)} \right]}_{\substack{{R_\lambda }\left( {x,y} \right) \in {L^2}\left( {M \times M} \right) \\ {\text{ rezolv. op. magfgve}}}}b\left( y \right)dy} .\]
A sor \(L^{2}\left( {M \times M} \right)\)-ben konvergens, ha
\(\left| \lambda \right| > r_{\sigma}\left( \mathcal{K} \right)\).

A bizonyítás alapja:
\(\left. \mathcal{K}_{j}\left( {x,y} \right) = {\int_{M}{\mathcal{K}_{j - 1}\left( {x,t} \right)\mathcal{K}\left( {t,y} \right)dt}}\Rightarrow K^{j - 1} \right.\)
operátor alkalmazva
\(\left. t\mapsto\mathcal{K}\left( {t,y} \right) \right.\) függvényre
(\(y\) rögzített): \[\begin{gathered}
  {\left\{ {\int_M {{{\left| {{\mathcal{K}_j}\left( {x,y} \right)} \right|}^2}dx} } \right\}^{1/2}} \leqslant \left\| {{\mathcal{K}^{j - 1}}} \right\|{\left\{ {\int_M {{{\left| {\mathcal{K}\left( {t,y} \right)} \right|}^2}dt} } \right\}^{1/2}} \\ 
   \Downarrow  \\ 
  \int_M {{{\left| {{\mathcal{K}_j}\left( {x,y} \right)} \right|}^2}dx}  \leqslant {\left\| {{K^{j - 1}}} \right\|^2}\int_M {{{\left| {\mathcal{K}\left( {t,y} \right)} \right|}^2}dt} . \\ 
\end{gathered} \] Integrálva \(y\) szerint:
\[{\int_{M \times M}{\left| {\mathcal{K}_{j}\left( {x,y} \right)} \right|^{2}dxdy}} \leq \left\| K^{j - 1} \right\|^{2}{\int_{M \times M}{\left| {\mathcal{K}\left( {t,y} \right)} \right|^{2}dtdy}}.\]
Nézzük:
\[\int_{M \times M} {\frac{1}{{{{\left| \lambda  \right|}^{2\left( {j + 1} \right)}}}}{{\left| {{K_j}\left( {x,y} \right)} \right|}^2}dxdy}  \leqslant \underbrace {\frac{1}{{\left| {{\lambda ^{2\left( {j + 1} \right)}}} \right|}}{{\left\| {{K^{j - 1}}} \right\|}^2}}_{\sum\limits_{j = 1}^\infty  {{\text{ sor konv}}{\text{. ha }}\left| \lambda  \right| > {r_\sigma }\left( K \right)} } \cdot \left\| K \right\|_{{L^2}\left( {M \times M} \right)}^2,\]így
a bal oldalból képzett számsor (ami \(\geq 0\)) is konvergens.

\hypertarget{folytonos-magu-integraloperatorok.}{%
\paragraph{Folytonos magú
integráloperátorok.}\label{folytonos-magu-integraloperatorok.}}

Legyen \(\Omega \subset {\mathbb{R}}^{n}\) korlátos tartomány (azaz
nyílt és összefüggő), \(X: = C\left( \overline{\Omega} \right)\),
\(\left. \overline{\Omega}\rightarrow{\mathbb{K}} \right.\) folytonos
függvények (a felülvonás a lezárást jelenti), tehát
\(C\left( \overline{\Omega} \right)\) az \(\Omega\) korlátos tartomány
lezárásán értelmezett folytonos függvények tere a
\(\left\| \phi \right\| = \sup\limits_{\Omega}\left| \phi \right|\)
normával. Legyen
\(\mathcal{K} \in C\left( {\overline{\Omega} \times \overline{\Omega}} \right)\),
\(\psi\left( x \right): = \left( {K\phi} \right)\left( x \right): = {\int_{\overline{\Omega}}{\mathcal{K}\left( {x,y} \right)\phi\left( y \right)}}dy\).

\begin{allitas}

Állítás:\\
\(\left. K:C\left( \overline{\Omega} \right)\rightarrow C\left( \overline{\Omega} \right) \right.\)
korlátos, lineáris operátor.

\end{allitas}

\begin{bizonyitas}

Bizonyítás:\\
\(\left| {\psi\left( x \right)} \right| = \left| {\int_{\overline{\Omega}}{\mathcal{K}\left( {x,y} \right)\phi\left( y \right)dy}} \right| \leq {\int_{\overline{\Omega}}{\left| {\mathcal{K}\left( {x,y} \right)} \right| \cdot \left| {\phi\left( y \right)} \right|dy}} \leq \left\| \phi \right\|{\int_{\overline{\Omega}}{\left| {\mathcal{K}\left( {x,y} \right)} \right|dy}} \leq \left\| \phi \right\|\sup\limits_{x \in \overline{\Omega}}{\int_{\overline{\Omega}}{\left| {\mathcal{K}\left( {x,y} \right)} \right|dy}}\).
Itt is igaz:
\(\left( {K^{j}\phi} \right)\left( x \right) = {\int_{\overline{\Omega}}{\mathcal{K}_{j}\left( {x,y} \right)\phi\left( y \right)dy}}\).
\(\mathcal{K}_{j}\left( {x,y} \right) = {\int_{\overline{\Omega}}{\mathcal{K}_{j - 1}\left( {x,t} \right)\mathcal{K}\left( {t,y} \right)dt}}\),
\(K_{j}\) folytonos.

\end{bizonyitas}

\hypertarget{egy-specialis-eset}{%
\paragraph{Egy speciális eset}\label{egy-specialis-eset}}

Az előbbi spec esete:
\(\overline{\Omega} = \left\lbrack {a,b} \right\rbrack \subset {\mathbb{R}}\),
ekkor
\(\mathcal{K} \in C\left( {\left\lbrack {a,b} \right\rbrack \times \left\lbrack {a,b} \right\rbrack} \right)\),
továbbá \(\mathcal{K}\left( {x,y} \right) = 0\), ha \(y > x\).
\(\left( {K\phi} \right)\left( x \right): = {\int\limits_{a}^{b}{\mathcal{K}\left( {x,y} \right)\phi\left( y \right)dy}} = {\int\limits_{a}^{x}{\mathcal{K}\left( {x,y} \right)\phi\left( y \right)dy}}\)
Voltera típusú operátor. Erre is igaz, hogy
\(\left. \mathcal{K}:C\left\lbrack {a,b} \right\rbrack\rightarrow C\left\lbrack {a,b} \right\rbrack \right.\)
folytonos lineáris operátor.

\begin{allitas}

Állítás:\\
\(r_{\sigma}\left( K \right) = 0\), így \(\lambda \neq 0\) esetén
\(\lambda\) reguláris érték, azaz létezik egyértelmű megoldása a
\(\lambda\phi\left( x \right) - {\int\limits_{a}^{x}{\mathcal{K}\left( {x,y} \right)\phi\left( y \right)dy}} = b\left( x \right)\)
másodfajú egyenletlnek bármely folytonos \(b\left( x \right)\) esetén.

\end{allitas}

\begin{bizonyitas}

Bizonyítás:\\
\(\mathcal{K}_{j}\left( {x,y} \right) = {\int\limits_{a}^{b}{\mathcal{K}_{j - 1}\left( {x,t} \right)\mathcal{K}\left( {t,y} \right)dt}}\),
speciálisan
\[{\mathcal{K}_2}\left( {x,y} \right) = \int\limits_a^b {\underbrace {\mathcal{K}\left( {x,t} \right)}_{0{\text{\;ha\;}}t > x}\underbrace {\mathcal{K}\left( {t,y} \right)}_{0{\text{\;ha\;}}y > t}dt}  = \int\limits_y^x {\mathcal{K}\left( {x,t} \right)\mathcal{K}\left( {t,y} \right)dt},\]
mert csak \(y \leq t \leq x\) esetén nem 0 az integrandus. Így
\(\mathcal{K}_{2}\left( {x,y} \right) = 0\), ha \(y > x\).
\(\mathcal{K}_{3}\left( {x,y} \right) = {\int\limits_{y}^{x}{\mathcal{K}_{2}\left( {x,t} \right)\mathcal{K}\left( {t,y} \right)dt}} = 0\)
ha \(y > x\). Ekkor
\(\left\| K \right\| \leq \sup\limits_{x \in {\lbrack{a,b}\rbrack}}{\int\limits_{a}^{b}{\left| {\mathcal{K}\left( {x,y} \right)} \right|dy}} \leq \alpha\left( {b - a} \right)\),
ugyanis
\(\left. \mathcal{K} \in C\left( {\left\lbrack {a,b} \right\rbrack \times \left\lbrack {a,b} \right\rbrack} \right)\Rightarrow\mathcal{K} \right.\)
korlátos és így
\(\left| {\mathcal{K}\left( {x,y} \right)} \right| \leq \alpha,\quad\forall x,y \in \left\lbrack {a,b} \right\rbrack\).

\end{bizonyitas}

\(\left\| K^{2} \right\| \leq \sup\limits_{x \in {\lbrack{a,b}\rbrack}}{\int\limits_{a}^{b}{\left| {\mathcal{K}_{2}\left( {x,y} \right)} \right|dy}} = \sup\limits_{x \in {\lbrack{a,b}\rbrack}}{\int\limits_{a}^{x}{\left| {\mathcal{K}_{2}\left( {x,y} \right)} \right|dy}}\).
Az integrandusra
\[\left| {{\mathcal{K}_2}\left( {x,y} \right)} \right| = \left| {\int\limits_y^x {\mathcal{K}\left( {x,t} \right)\mathcal{K}\left( {t,y} \right)dt} } \right| \leqslant \int\limits_y^x {\underbrace {\left| {\mathcal{K}\left( {x,t} \right)} \right|}_{ < \alpha }\underbrace {\left| {\mathcal{K}\left( {t,y} \right)} \right|}_{ \leqslant \alpha }dt}  \leqslant {\alpha ^2}\left( {x - y} \right),\]
ha \(x > y\). Így \[\begin{aligned}
  \left\| {{K^2}} \right\| &  \leqslant \mathop {\sup }\limits_{x \in \left[ {a,b} \right]} \int\limits_a^x {\left| {{\mathcal{K}_2}\left( {x,y} \right)} \right|dy}  \\ 
   &  \leqslant \mathop {\sup }\limits_{x \in \left[ {a,b} \right]} \int\limits_a^x {{\alpha ^2}\left( {x - y} \right)dy}  \\ 
   &  = {\alpha ^2}\mathop {\sup }\limits_{x \in \left[ {a,b} \right]} \left[ { - \frac{{{{\left( {x - y} \right)}^2}}}{2}} \right]_{y = a}^x \\ 
   &  = {\alpha ^2}\mathop {\sup }\limits_{x \in \left[ {a,b} \right]} \frac{{{{\left( {x - a} \right)}^2}}}{2} = {\alpha ^2}\frac{{{{\left( {b - a} \right)}^2}}}{2}. \\ 
\end{aligned} \]

\(\left\| K^{3} \right\|\)-re hasonló módon járunk el. Ekkor
\[\left| {{\mathcal{K}_3}\left( {x,y} \right)} \right| = \left| {\int\limits_y^x {\mathcal{K}\left( {x,t} \right){\mathcal{K}_2}\left( {t,y} \right)dt} } \right| \leqslant \int\limits_y^x {\underbrace {\left| {\mathcal{K}\left( {x,t} \right)} \right|}_{ \leqslant \alpha }\underbrace {\left| {{\mathcal{K}_2}\left( {t,y} \right)} \right|}_{ \leqslant {\alpha ^2}\left( {t - y} \right)}dt}  \leqslant {\alpha ^3}\frac{{{{\left( {x - y} \right)}^2}}}{2}.\]
Így
\(\left\| K^{3} \right\| \leq \sup\limits_{x \in {\lbrack{a,b}\rbrack}}{\int\limits_{a}^{x}{\left| {\mathcal{K}_{3}\left( {x,y} \right)} \right|dy}} \leq \sup\limits_{x \in {\lbrack{a,b}\rbrack}}{\int\limits_{a}^{x}{\alpha^{3}\frac{\left( {x - y} \right)^{2}}{2}dy}} = \alpha^{3}\sup\limits_{x \in {\lbrack{a,b}\rbrack}}\frac{\left( {x - a} \right)^{3}}{3!} \leq \alpha^{3}\frac{\left( {b - a} \right)^{3}}{3!}\).
Teljes indukcióval bizonyítható, hogy
\(\left. \left\| K^{j} \right\| \leq \alpha^{j}\frac{\left( {b - a} \right)^{j}}{j!}\Rightarrow\left\| K^{j} \right\|^{1/j} = \alpha\frac{b - a}{\left( {j!} \right)^{1/j}}\rightarrow 0 \right.\),
ha \(\left. j\rightarrow\infty \right.\).

\hypertarget{hilbert-ter-operatorai}{%
\section{Hilbert tér operátorai}\label{hilbert-ter-operatorai}}

\hypertarget{az-adjungalt-operator}{%
\subsection{Az adjungált operátor}\label{az-adjungalt-operator}}

Legyen \(X\) Hilbert tér, \(\left. A:D_{A}\rightarrow X \right.\)
lineáris operátor, ahol \(D_{A}\) az \(A\)-nak az értelmezési
tartománya, \(D_{A} \subset X\), \(y \in X\) elem.

\textbf{Kérdés}: létezik-e illetve hány \(y^{*} \in X\) létezik, melyre
\(\left\langle {Ax,y} \right\rangle = \left\langle {x,y^{*}} \right\rangle\),
\(\forall x \in D_{A}\) esetén? Mi az egyértelműség feltétele?

\begin{allitas}

Állítás:\\
Legfeljebb egy \(y^{*}\) létezik
\(\left. \Leftrightarrow\overline{D_{A}} = X \right.\), vagyis ha az
értelmezési tartomány sűrű \(X\)-ben.

\end{allitas}

\begin{bizonyitas}

Bizonyítás:\\
Legfeljebb egy \(y^{*}\) létezik \(\Leftrightarrow\) hogy ha
\(\left\langle {x,y^{*}} \right\rangle = \left\langle {x,\widetilde{y}} \right\rangle\),
\(\forall x \in D_{A}\)-ból következik, hogy \(y^{*} = \widetilde{y}\).
\(\left\langle {x,y^{*}} \right\rangle = \left\langle {x,\widetilde{y}} \right\rangle\),
\(\forall x \in D_{A}\) pontosan azt jelenti, hogy
\(\left\langle {x,y^{*} - \widetilde{y}} \right\rangle = 0\),
\(\forall x \in D_{A}\). Ebből következik:
\(\left. y^{*} = \widetilde{y}\Leftrightarrow\overline{D_{A}} = X \right.\).
(Felhasználjuk, hogy a skalárszorzat folytonosan függ a tényezőktől.)

\end{bizonyitas}

\begin{definicio}

Definíció:\\
Legyen \(X\) Hilbert tér, \(\left. A:D_{A}\rightarrow X \right.\)
lineáris operátor, \(\overline{D_{A}} = X\). Ekkor \(A\) operátor
adjungáltját, \(A^{*}\) operátort így értelmezzük:
\[D_{A^{*}}: = \left\{ {y \in X:\exists y^{*} \in X:\left\langle {Ax,y} \right\rangle = \left\langle {x,y^{*}} \right\rangle\ \forall x \in D_{A}} \right\}\]
és \(A^{*}\left( y \right): = y^{*}\).

\end{definicio}

\begin{megjegyzes}

Megjegyzés:\\
\(0 \in D_{A^{*}}\), ugyanis
\(\left\langle {Ax,0} \right\rangle = \left\langle {x,0} \right\rangle = 0\),
\(\forall x \in D_{A}\).

\end{megjegyzes}

\begin{allitas}

Állítás:\\
\(A^{*}\) lineáris operátor.

\end{allitas}

\begin{bizonyitas}

Bizonyítás:\\
Legyen \(y_{1},y_{2} \in D_{A^{*}}\)! Ekkor
\(\left\langle {Ax,y_{1}} \right\rangle = \left\langle {x,A^{*}\left( y_{1} \right)} \right\rangle\),
\(\forall x \in D_{A}\) és
\(\left\langle {Ax,y_{2}} \right\rangle = \left\langle {x,A^{*}\left( y_{2} \right)} \right\rangle\),
\(\forall x \in D_{A}\). Így
\(\left\langle {Ax,y_{1}} \right\rangle + \left\langle {Ax,y_{2}} \right\rangle = \left\langle {x,A^{*}\left( y_{1} \right)} \right\rangle + \left\langle {x,A^{*}\left( y_{2} \right)} \right\rangle\).
\(\left\langle {Ax,y_{1} + y_{2}} \right\rangle = \left\langle {x,A^{*}\left( y_{1} \right) + A^{*}\left( y_{2} \right)} \right\rangle\),
\(\forall x \in D_{A}\). Ebből következik, hogy
\(A^{*}\left( {y_{1} + y_{2}} \right) = A^{*}\left( y_{1} \right) + A^{*}\left( y_{2} \right)\).
Hasonlóan igazolható
\(A^{*}\left( {\lambda g} \right) = \lambda A^{*}\left( g \right)\).

\end{bizonyitas}

\begin{tetel}

Tétel:\\
Legyen \(\left. A:X\rightarrow X \right.\) korlátos lineáris operátor.
Ekkor \(\left. A^{*}:X\rightarrow X \right.\) korlátos lineáris operátor
és \(\left\| A^{*} \right\| = \left\| A \right\|\).

\end{tetel}

\begin{bizonyitas}

Bizonyítás:\\
Tekintsünk tetszőleges, rögzített \(y \in X\) elemet! Ekkor
\(f\left( x \right): = \left\langle {Ax,y} \right\rangle\), \(f\)
lineáris funkcionál korlátos is:
\(\left| {f\left( x \right)} \right| = \left| \left\langle {Ax,y} \right\rangle \right| \leq \left\| {Ax} \right\| \cdot \left\| y \right\| \leq \left\| A \right\| \cdot \left\| x \right\| \cdot \left\| y \right\| = \left( {\left\| A \right\|\left\| y \right\|} \right) \cdot \left\| x \right\|\),
így
\(\left\| f \right\| \leq \left\| A \right\| \cdot \left\| y \right\|\).
A Riesz-tételből most következik, hogy
\(\exists!y^{*} \in X:f\left( x \right) = \left\langle {x,y^{*}} \right\rangle\),
azaz
\(\left\langle {Ax,y} \right\rangle = \left\langle {x,y^{*}} \right\rangle\),
\(\forall x \in X\)-re. Így \(D_{A^{*}} = X\), \(A^{*}y = y^{*}\).
Továbbá
\(\left\| {A^{*}y} \right\| = \left\| y^{*} \right\| = \left\| f \right\| \leq \left\| A \right\| \cdot \left\| y \right\|\),
ezért \(A^{*}\) korlátos és
\(\left\| A^{*} \right\| \leq \left\| A \right\|\). Az egyenlőség abból
fog következni, hogy
\(\left. \left( A^{*} \right)^{*} = A\Rightarrow\left\| A \right\| = \left\| \left( A^{*} \right)^{*} \right\| \leq \left\| A^{*} \right\| \right.\).

\end{bizonyitas}

\begin{ajanlo}

\begin{ajanlofig}

\href{https://xkcd.com}{\includegraphics[width=5.20833in,height=2.82292in]{wikipedian_protester.png}}

\end{ajanlofig}

Text

\end{ajanlo}

Legyen \(\left. A:X\rightarrow X \right.\) korlátos lineáris operátor!
Láttuk már, hogy \(\left. A^{*}:X\rightarrow X \right.\) operátor
korlátos és lineáris, és
\(\left\| A^{*} \right\| \leq \left\| A \right\|\).

\begin{tetel}

Tétel:\\
Legyenek \(\left. A,B:X\rightarrow X \right.\) korlátos lineáris
operátor! Ekkor

\begin{enumerate}
\def\labelenumi{\arabic{enumi}.}
\tightlist
\item
  \(\left( {A + B} \right)^{*} = A^{*} + B^{*}\)
\item
  \(\left( {\lambda A} \right)^{*} = \overline{\lambda}A^{*}\)
\item
  \(\left( A^{*} \right)^{*} = A\)
\item
  \(I = I^{*}\), \(0^{*} = 0\)
\item
  \(\left( {AB} \right)^{*} = B^{*}A^{*}\).
\end{enumerate}

\end{tetel}

\begin{bizonyitas}

Bizonyítás:\\
Legyenek \(x,y \in X\) !

\begin{enumerate}
\def\labelenumi{\arabic{enumi}.}
\tightlist
\item
  \(\left\langle {\left( {A + B} \right)x,y} \right\rangle = \left\langle {Ax + Bx,y} \right\rangle = \left\langle {Ax,y} \right\rangle + \left\langle {Bx,y} \right\rangle = \left\langle {x,A^{*}y} \right\rangle + \left\langle {x,B^{*}y} \right\rangle =\)
  \(= \left\langle {x,A^{*}y + B^{*}y} \right\rangle = \left\langle {x,\left( {A^{*} + B^{*}} \right)y} \right\rangle\)
\item
  \(\left\langle {Ax,y} \right\rangle = \left\langle {x,A^{*}y} \right\rangle = \overline{\left\langle {A^{*}y,x} \right\rangle} = \overline{\left\langle {y,\left( A^{*} \right)^{*}x} \right\rangle} = \left\langle {\left( A^{*} \right)^{*}x,y} \right\rangle\),tehát
  \(Ax = \left( A^{*} \right)^{*}x\),
  \(\left. \forall x \in X\Rightarrow A = \left( A^{*} \right)^{*} \right.\),így
  \(\left\| A^{*} \right\| \leq \left\| \left( A^{*} \right)^{*} \right\| = \left\| A \right\|\),így
  az előző tétellel együtt:
  \(\left\| A \right\| = \left\| A^{*} \right\|\).
\item
  \(\left\langle {x,\left( {AB} \right)^{*}y} \right\rangle = \left\langle {ABx,y} \right\rangle = \left\langle {Bx,A^{*}y} \right\rangle = \left\langle {x,B^{*}A^{*}y} \right\rangle\)
\end{enumerate}

\end{bizonyitas}

\begin{megjegyzes}

Megjegyzés:\\
Mi a helyezet a lineáris operátorok esetén (ha nem korlátos)?
\(D_{A},D_{B} \subset X\), \(\overline{D_{A}} = \overline{D_{B}} = X\).

\end{megjegyzes}

Jelölés: ha \(A^{*}x = Ax\), \(\forall x \in D_{A}\),
\(D_{A} \subset D_{A^{*}}\), akkor \(A^{*}\) kiterjesztése \(A\)-nak s
ezt így jelöljük: \(A \subset A^{*}\). Ezzel a jelöléssel:
\(\left( {A + B} \right)^{*} \supset A^{*} + B^{*}\) és
\(D_{A^{*} + B^{*}} = D_{A^{*}} \cap D_{B^{*}}\). Ugyanis
\(\forall y \in \left( {D_{A^{*}} \cap D_{B^{*}}} \right)\) esetén
\(\left\langle {\left( {A + B} \right)x,y} \right\rangle = \left\langle {x,\left( {A^{*} + B^{*}} \right)y} \right\rangle\),
\(\forall x \in \left( {D_{A} \cap D_{B}} \right)\).\\
Továbbá \(\left( {\lambda A} \right)^{*} = \overline{\lambda}A^{*}\),
\(\left( {AB} \right)^{*} \supset B^{*}A^{*}\),
\(\left( A^{*} \right)^{*} \supset A\) és 1
\(\left. A \subset B\Rightarrow A^{*} \supset B^{*} \right.\).

\begin{pelda}

Példák:\\
\(X: = {\mathbb{K}}^{n}\). Tudjuk, hogy ekkor minden lineáris operátor
korlátos.
\(\left. A:{\mathbb{K}}^{n}\rightarrow{\mathbb{K}}^{n} \right.\)
lineáris korlátos operátor. Tudjuk, hogy \(A\) reprezentálható egy
\(\mathcal{A}\) (valós vagy komplex elemekből alkotott),
\(n \times n\)-es mátrixszal úgy, hogy \(\mathcal{A}x = Ax\). Ekkor
\(\left. A^{*}:{\mathbb{K}}^{n}\rightarrow{\mathbb{K}}^{n} \right.\)
korlátos lineáris operátor. Kérdés: mi a lesz ennek a mátrixa?
\[\mathcal{A} = \left( \begin{array}{lll}
a_{11} & \cdots & a_{1n} \\
 \vdots & \ddots & \vdots \\
a_{n1} & \cdots & a_{nn} \\
\end{array} \right)\text{,\quad}a_{jk} \in {\mathbb{K}}\] Ekkor
\(x,y \in {\mathbb{K}}^{n}\) esetén
\(\left\langle {\mathcal{A}x,y} \right\rangle = {\sum\limits_{j = 1}^{n}{\left\lbrack {\sum\limits_{k = 1}^{n}{a_{jk}x_{k}}} \right\rbrack\overline{y_{j}}}} = {\sum\limits_{k = 1}^{n}x_{k}}\left\lbrack {\sum\limits_{j = 1}^{n}{a_{jk}\overline{y_{j}}}} \right\rbrack = {\sum\limits_{k = 1}^{n}{x_{k}\overline{\left\lbrack {\underset{j=1}{\overset{n}{\sum}}{\overline{a_{jk}}y_{j}}} \right\rbrack}}} = {\sum\limits_{k = 1}^{n}{x_{k}\overline{\left\lbrack {\underset{j=1}{\overset{n}{\sum}}{a_{kj}^{*}y_{j}}} \right\rbrack} = \left\langle {x,\mathcal{A}^{*}y} \right\rangle}}\),
vagyis \(a_{kj}^{*} = \overline{a_{jk}}\), vagyis
\(\mathcal{A}^{*} = \left( \begin{array}{lll} a_{11}^{*} & \cdots & a_{1n}^{*} \\  \vdots & \ddots & \vdots \\ a_{n1}^{*} & \cdots & a_{nn}^{*} \\ \end{array} \right) = \left( \begin{array}{lll} \overline{a_{11}} & \cdots & \overline{a_{n1}} \\  \vdots & \ddots & \vdots \\ \overline{a_{1n}} & \cdots & \overline{a_{nn}} \\ \end{array} \right)\).

\end{pelda}

\hypertarget{negyzetesen-integralhato-magu-integral-operatorok-valos-vagy-komplex-fuggvenyeken}{%
\subsubsection{\texorpdfstring{\textbf{Négyzetesen integrálható magú
integrál operátorok valós vagy komplex
függvényeken}}{Négyzetesen integrálható magú integrál operátorok valós vagy komplex függvényeken}}\label{negyzetesen-integralhato-magu-integral-operatorok-valos-vagy-komplex-fuggvenyeken}}

Legyen \(X: = L^{2}\left( M \right),M \subset {\mathbb{R}}^{n}\) mérhető
halmaz és
\[\mathcal{K} \in L^{2}\left( {M \times M} \right),\quad\left( {K\phi} \right)\left( x \right): = {\int_{M}{\mathcal{K}\left( {x,y} \right)\phi\left( y \right)dy}}.\]
Tudjuk, hogy
\(\left. K:L^{2}\left( M \right)\rightarrow L^{2}\left( M \right) \right.\)
lineáris operátor, node mi \(K^{*}\)? Legyen
\(\phi,\psi \in L^{2}\left( M \right)\), ekkor \[\begin{aligned}
  \left\langle {K\phi ,\psi } \right\rangle  &  = \int_M {\left( {K\phi } \right)\left( x \right)\overline {\psi \left( x \right)} dx}  \\ 
   &  = \int_M {\left[ {\int_M {\mathcal{K}\left( {x,y} \right)\phi \left( y \right)dy} } \right]\overline {\psi \left( x \right)} dx}  \\ 
  {\text{(Fubini) }} &  = \int_M {\phi \left( y \right)\left[ {\int_M {\mathcal{K}\left( {x,y} \right)\overline {\psi \left( x \right)} dx} } \right]dy}  \\ 
   &  = \int_M {\phi \left( y \right)\overline {\left[ {\int_M {\overline {\mathcal{K}\left( {x,y} \right)} \psi \left( x \right)dx} } \right]} dy}  \\ 
  {\text{(}}x \leftrightarrow y{\text{ csere) }} &  = \int_M {\phi \left( x \right)\left[ {\overline {\int_M {\underbrace {\overline {\mathcal{K}\left( {y,x} \right)} }_{: = {\mathcal{K}^*}\left( {x,y} \right)}\psi \left( y \right)dy} } } \right]dx}  \\ 
   &  = \int_M {\phi \left( x \right)\left[ {\overline {\int_M {{\mathcal{K}^*}\left( {x,y} \right)\psi \left( y \right)dx} } } \right]dy} . \\ 
\end{aligned} \] A bevezetett jelöléssel összhangban
\(\left( {K^{*}\psi} \right)\left( x \right): = {\int_{M}{\mathcal{K}^{*}\left( {x,y} \right)\psi\left( y \right)dy}}\),
így az korábbiakkal együtt:
\(\left\langle {K\phi,\psi} \right\rangle = {\int_{M}{\phi\left( x \right)\overline{\left( {K^{*}\psi} \right)\left( x \right)}dx}} = \left\langle {\phi,K^{*}\psi} \right\rangle\).

\begin{allitas}

Állítás:\\
Tetszőleges \(A\) lineáris operátor esetén (melyre \(D_{A} \subset X\),
\(\overline{D_{A}} = X\)) \(A^{*}\) zárt operátor.

\end{allitas}

\begin{bizonyitas}

Bizonyítás:\\
Azt kellene belátni, hogy ha \(y_{j} \in D_{A^{*}}\),
\(\left. \left( y_{j} \right)_{j \in {\mathbb{N}}}\rightarrow y \right.\)
\(X\)-ben, továbbá
\(\left. \left( {A^{*}y_{j}} \right)\rightarrow z \right.\) \(X\)-ben
\(\left. \Rightarrow y \in D_{A^{*}} \right.\) és \(A^{*}y = z\).
Tudtuk, hogy
\(\left\langle {Ax,y_{j}} \right\rangle = \left\langle {x,A^{*}y_{j}} \right\rangle\),
\(\forall x \in D_{A},\forall j\), így
\(\left. j\rightarrow\infty \right.\) esetén
\(\left\langle {Ax,y} \right\rangle = \left\langle {x,z} \right\rangle\),
\(\forall x \in D_{A}\). Ez azt jelenti, hogy \(y \in D_{A^{*}}\) és
\(z = A^{*}y\).

\end{bizonyitas}

\begin{tetel}

Tétel:\\
Legyen \(X\) Hilbert tér, \(\left. A:X\rightarrow X \right.\) korlátos
lineáris operátor és \(\lambda \in {\mathbb{K}}\). Ekkor
\({\overline{R_{({\lambda I-A})}}}^{\bot} = S_{\overline{\lambda}}\left( A^{*} \right): = \left\{ {x \in X:\left( {\overline{\lambda}I - A^{*}} \right)x = 0} \right\}\),
ahol \(R\) az értékkészletet jelöli.

\end{tetel}

\begin{bizonyitas}

Bizonyítás:\\
Világos, hogy \(R_{({\lambda I - A})}\) lineáris altér, ezért
\(\overline{R_{({\lambda I-A})}}\) zárt altér. Másrészt
\(S_{\overline{\lambda}}\left( A^{*} \right)\) is zárt altér. Az
\(S_{\overline{\lambda}}\left( A^{*} \right)\) halmaz azért zárt, mert
\(A^{*}\) folytonos lineáris operátor.

\begin{itemize}
\tightlist
\item
  Először tfh \(y \in {\overline{R_{\lambda I-A}}}^{\bot}\), ekkor
  \(0 = \left\langle {\underbrace {\left( {\lambda I - A} \right)x}_{ \in {R_{\lambda I - A}} \subset \overline {{R_{\lambda I - A}}} },y} \right\rangle = \left\langle {x,{{\left( {\lambda I - A} \right)}^ * }y} \right\rangle\),
  ez igaz
  \(\forall x \in X \Rightarrow \underbrace {{{\left( {\lambda I - A} \right)}^ * }}_{ = \bar \lambda I - {A^ * }}y = 0\),
  vagyis \(y \in S_{\overline{\lambda}}\left( A^{*} \right)\).
\item
  tfh \(y \in S_{\overline{\lambda}}\left( A^{*} \right)\),azaz
  \(\left( {\overline{\lambda}I - A^{*}} \right)y = 0\),
  \(\forall x \in X\),így
  \(\left\langle {\left( {\lambda I - A} \right)x,y} \right\rangle = \left\langle {x,\left( {\lambda I - A} \right)^{*}y} \right\rangle = 0\),vagyis
  \(y\bot R_{\lambda I - A}\) minden elemére
  \(\left. \Rightarrow y\bot\overline{R_{\lambda I-A}} \right.\) minden
  elemére.
\end{itemize}

\end{bizonyitas}

\begin{megjegyzes}

Megjegyzés:\\
Spec eset, mikor \(R_{\lambda I - A}\) zárt halmaz, azaz
\(R_{\lambda I - A} = \overline{R_{\lambda I-A}}\). Ekkor a fenti
tételből következik: \(\left( {\lambda I - A} \right)x = b\) másodfajú
egyenletnek létezik \(x \in X\) megoldása pontosan akkor, ha
\(b \in R_{\lambda I - A} = S_{\overline{\lambda}}\left( A^{*} \right)^{\bot}\),
azaz \(\left\langle {b,y} \right\rangle = 0\) a
\(\left( {\lambda I - A} \right)^{*}y = 0\) egyenlet \(\forall y \in X\)
megoldására. Később látni fogjuk, hogy ha \(A\) ún. kompakt lineáris
operátor, akkor \(\lambda \neq 0\) esetén az \(R_{\lambda I - A}\) zárt
halmaz.

\end{megjegyzes}

\hypertarget{szimmetrikus-es-onadjungalt-operatorok}{%
\subsubsection{Szimmetrikus és önadjungált
operátorok}\label{szimmetrikus-es-onadjungalt-operatorok}}

\begin{definicio}

Definíció:\\
Legyen \(X\) Hilbert tér, \(D_{A} \subset X\) és
\(\overline{D_{A}} = X\) és \(\left. A:D_{A}\rightarrow X \right.\)
lineáris operátor. Ekkor \(A\)-t önadjungáltnak nevezzük, ha
\(A^{*} = A\) (ekkor ugyanott vannak értelmezve, \(D_{A^{*}} = D_{A}\)).

\end{definicio}

\begin{definicio}

Definíció:\\
Legyen \(X\) Hilbert tér, \(D_{A} \subset X\) és
\(\overline{D_{A}} = X\) és \(\left. A:D_{A}\rightarrow X \right.\)
lineáris operátor. Ekkor \(A\)-t szimmetrikusnak nevezzük, ha
\(A \subset A^{*}\). Tehát minden önadjungált operátor egyúttal
szimmetrikus is.

\end{definicio}

\begin{megjegyzes}

Megjegyzés:\\
Ekvivalens definíció: \(A\) szimmetrikus, ha
\(\left\langle {Ax,y} \right\rangle = \left\langle {x,Ay} \right\rangle\),
\(\forall x,y \in D_{A}\).

\end{megjegyzes}

\begin{pelda}

Példa:\\
Ha \(X = {\mathbb{K}}^{n}\), akkor
\(\left. A:{\mathbb{K}}^{n}\rightarrow{\mathbb{K}}^{n} \right.\)-nak
megfelel egy \(\mathcal{A}\) mátrix. Tudjuk, hogy \(A^{*}\) mátrixa
\(\mathcal{A}^{*}\), melynek elemei \(a_{jk}^{*} = \overline{a_{kj}}\).
Ekkor \(A\) önadjungált
\(\left. \Leftrightarrow a_{jk}^{*} = a_{jk} \right.\), azaz
\(a_{jk} = \overline{a_{kj}}\).

\end{pelda}

\begin{pelda}

Példa:\\
Legyen \(X: = L^{2}\left( M \right)\), \(M \subset {\mathbb{R}}^{n}\)
mérhető halmaz,
\(\left( {K\phi} \right)\left( x \right): = {\int_{M}{\mathcal{K}\left( {x,y} \right)\phi\left( y \right)dy}}\)
korlátos operátor, ahol
\(\mathcal{K} \in L^{2}\left( {M \times M} \right)\). Ekkor
\(\left( {K^{*}\phi} \right)\left( x \right) = {\int_{M}{\mathcal{K}^{*}\left( {x,y} \right)\phi\left( y \right)dy}}\),
vagyis
\(\mathcal{K}^{*}\left( {x,y} \right) = \overline{\mathcal{K}\left( {y,x} \right)}\).
\(K\) önadjugnált pontosan akkor, ha
\(\mathcal{K}\left( {x,y} \right) = \overline{\mathcal{K}\left( {y,x} \right)}\)
majdnem minden \(x,y \in M\).

\end{pelda}

\begin{pelda}

Példa:\\
Legyen \(X: = L^{2}\left( 0,1 \right)\),
\(\left( {A\phi} \right)\left( t \right): = \phi^{''}\left( t \right)\),
midőn \(t \in \left\lbrack 0,1 \right\rbrack\), vagyis legyen \(A\) a
második derivált operátor (ami lineáris)!
\(D_{A}: = \left\{ {\phi \in C^{2}\left\lbrack 0,1 \right\rbrack:\phi\left( 0 \right) = 0,\phi\left( 1 \right) = 0} \right\}\),
erre belátható, hogy \(\overline{D_{A}} = L^{2}\left( 0,1 \right)\).

\end{pelda}

\begin{allitas}

Állítás:\\
\(A\) szimmetrikus operátor (de nem önadjungált). Ennek igazolásához
tekintsünk \(\phi,\psi \in D_{A}\) tetszőleges függvényeket, ekkor
parciális integrálással: \[\begin{aligned}
  \left\langle {A\phi ,\psi } \right\rangle  &  = \int\limits_0^1 {\left( {A\phi \left( t \right)} \right)\psi \left( t \right)dt}  \\ 
   &  = \int\limits_0^1 {\phi ''\left( t \right)\psi \left( t \right)dt}  \\ 
   &  = \left[ {\phi '\left( t \right)\psi \left( t \right)} \right]_0^1 - \int\limits_0^1 {\phi '\left( t \right)\psi '\left( t \right)dt}  \\ 
   &  =  - \left[ {\phi \left( y \right)\psi '\left( t \right)} \right]_0^1 + \int\limits_0^1 {\phi \left( t \right)\psi ''\left( t \right)dt}  \\ 
   &  = \left\langle {\phi ,A\psi } \right\rangle . \\ 
\end{aligned} \]

\end{allitas}

\begin{ajanlo}

\begin{ajanlofig}

\href{https://xkcd.com}{\includegraphics[width=5.20833in,height=2.82292in]{wikipedian_protester.png}}

\end{ajanlofig}

Text

\end{ajanlo}

\begin{allitas}

Állítás:\\
Legyen \(X\) komplex Hilbert tér! Ha \(D_{A} \subset X\),
\(\left. A:D_{A}\rightarrow X \right.\) szimmetrikus operátor, akkor
\(\left\langle {Ax,x} \right\rangle\) értéke valós
\(\forall x \in {\mathbb{D}}_{A}\) esetén.

\end{allitas}

\begin{bizonyitas}

Bizonyítás:\\
Mivel \(A\) szimmetrikus, ezért
\(\left\langle {Ax,x} \right\rangle = \left\langle {x,Ax} \right\rangle\),
\(\forall x \in D_{A}\), másrészt a skaláris szorzat tulajdonságából
következően:
\(\left. \left\langle {Ax,x} \right\rangle = \overline{\left\langle {x,Ax} \right\rangle}\Rightarrow\left\langle {x,Ax} \right\rangle = \overline{\left\langle {x,Ax} \right\rangle}\Rightarrow\left\langle {x,Ax} \right\rangle \right.\)
valós, így \(\left\langle {Ax,x} \right\rangle\) is valós.

\end{bizonyitas}

\begin{megjegyzes}

Megjegyzés:\\
Bebizonyítható, hogy ha \(X\) komplex Hilbert tér és
\(\left\langle {Ax,x} \right\rangle\) valós
\(\left. \forall x \in D_{A}\Rightarrow A \right.\) szimmetrikus.

\end{megjegyzes}

\begin{tetel}

Tétel:\\
Legyen \(X\) Hilbert tér (lehet valós is). Ha \(D_{A} \subset X\),
\(\left. A:D_{A}\rightarrow X \right.\) szimmetrikus operátor, akkor
\(A\) minden sajátértéke valós és a különböző sajátértékekhez tartozó
sajátelemek ortogonálisak.

\end{tetel}

\begin{bizonyitas}

Bizonyítás:

\begin{itemize}
\tightlist
\item
  Tfh \(Ax = \lambda x\) valamely \(0 \neq x \in D_{A}\) elemre,
  \(\lambda \in {\mathbb{K}}\). Ekkor
  \(\underbrace {\left\langle {Ax,x} \right\rangle }_{{\text{valós}}} = \left\langle {\lambda x,x} \right\rangle = \lambda {\left\| x \right\|^2}\).
  A norma értéke valós, így a sajátérték is az, mert szorzatuk valós.
\item
  Tfh \[\begin{array}{l}
  {Ax_{1} = \lambda_{1}x_{1}} \\
  {Ax_{2} = \lambda_{2}x_{2}} \\
  \end{array}\] és \(\lambda_{1} \neq \lambda_{2}\) valós sajátértékek.
  Szorozzuk skalárisan jobbról előbbit \(x_{2}\)-vel!
  \[\left\langle {Ax_{1},x_{2}} \right\rangle = \left\langle {\lambda_{1}x_{1},x_{2}} \right\rangle = \lambda_{1}\left\langle {x_{1},x_{2}} \right\rangle\]illetve
  \[\left\langle {Ax_{1},x_{2}} \right\rangle = \left\langle {x_{1},Ax_{2}} \right\rangle = \left\langle {x_{1},\lambda_{2}x_{2}} \right\rangle = \lambda_{2}\left\langle {x_{1},x_{2}} \right\rangle,\]így
  \[{\lambda _1}\left\langle {{x_1},{x_2}} \right\rangle  = {\lambda _2}\left\langle {{x_1},{x_2}} \right\rangle  \Leftrightarrow \underbrace {\left( {{\lambda _1} - {\lambda _2}} \right)}_{ \ne 0 \Leftarrow {\lambda _1} \ne {\lambda _2}}\left\langle {{x_1},{x_2}} \right\rangle  = 0 \Rightarrow \left\langle {{x_1},{x_2}} \right\rangle  = 0.\]
\end{itemize}

\end{bizonyitas}

\begin{tetel}

Tétel:\\
Legyen \(X\) Hilbert tér, \(\left. A:X\rightarrow X \right.\) korlátos
önadjungált operátor. Ekkor
\(\left\| A \right\| = \sup\left\{ {\left| \left\langle {Ax,x} \right\rangle \right|:x \in X,\left\| x \right\| = 1} \right\}\).

\end{tetel}

\begin{bizonyitas}

Bizonyítás:

Az operátor norma definíciója szerint
\(\left\| A \right\| = \sup\left\{ {\left\| {Ax} \right\|:x \in X,\left\| x \right\| = 1} \right\}\)
. Ezért egyrészt a Cauchy-Schwarz egyenlőtlenségből
\(\left| \left\langle {Ax,x} \right\rangle \right| \leq \left\| {Ax} \right\| \cdot \left\| x \right\| \leq \left\| A \right\| \cdot \left\| x \right\|^{2} = \left\| A \right\|\)
, ha \(\left\| x \right\| = 1\) . Jelöljük:
\(\alpha: = \sup\left\{ {\left| \left\langle {Ax,x} \right\rangle \right|:x \in X,\left\| x \right\| = 1} \right\}\)
. Az előbbiek szerint \(\alpha \leq \left\| A \right\|\) . Belátjuk a
fordított egyenlőtlenséget. Tetszőleges \(x,y \in X\) elemekre
\[\begin{aligned}
    \left\langle {A\left( {x + y} \right),x + y} \right\rangle  &  = \left\langle {Ax + Ay,x + y} \right\rangle  \\ 
    &  = \left\langle {Ax,x} \right\rangle  + \underbrace {\left\langle {Ay,x} \right\rangle }_{ = \langle y,Ax\rangle  = \overline {\langle Ax,y\rangle } } + \left\langle {Ax,y} \right\rangle  + \left\langle {Ay,y} \right\rangle  \\ 
    &  = \left\langle {Ax,x} \right\rangle  + \left\langle {Ay,y} \right\rangle  + 2\Re \left\langle {Ax,y} \right\rangle . \\ 
    \end{aligned} \]

Hasonlóképpen:
\[\left\langle {A\left( {x - y} \right),x - y} \right\rangle = \left\langle {Ax,x} \right\rangle + \left\langle {Ay,y} \right\rangle - 2\Re\left\langle {Ax,y} \right\rangle.\]
A kapott 1. egyenlőségből a 2-at kivonva: \[\begin{aligned}
  4\Re \left\langle {Ax,y} \right\rangle  &  = \left\langle {A\left( {x + y} \right),x + y} \right\rangle  - \left\langle {A\left( {x - y} \right),x - y} \right\rangle  \\ 
   &  \leqslant \left| {\left\langle {A\left( {x + y} \right),x + y} \right\rangle } \right| + \left| {\left\langle {A\left( {x - y} \right),x - y} \right\rangle } \right| \\ 
   &  \leqslant \alpha {\left\| {x + y} \right\|^2} + \alpha {\left\| {x - y} \right\|^2} \\ 
   &  = \alpha \left( {{{\left\| x \right\|}^2} + 2{{\left\langle {x,y} \right\rangle }^2} + {{\left\| y \right\|}^2} + {{\left\| x \right\|}^2} - 2{{\left\langle {x,y} \right\rangle }^2} + {{\left\| y \right\|}^2}} \right) \\ 
   &  = 2\alpha \left( {{{\left\| x \right\|}^2} + {{\left\| y \right\|}^2}} \right) \\ 
   \Downarrow  \\ 
  \Re \left\langle {Ax,y} \right\rangle  &  \leqslant \frac{\alpha }{2}\left( {{{\left\| x \right\|}^2} + {{\left\| y \right\|}^2}} \right). \\ 
\end{aligned} \]

Tetszőleges \(\lambda > 0\) számra: \[\begin{aligned}
  \underbrace {{{\left\| {Ax} \right\|}^2}}_{ \in \mathbb{R}_0^ + } &  = \left\langle {Ax,Ax} \right\rangle  \\ 
   &  = \left\langle {A\underbrace {\left( {\lambda x} \right)}_{: = f},\underbrace {Ax/\lambda }_{: = g}} \right\rangle  \\ 
   &  = \underbrace {\left\langle {Af,g} \right\rangle }_{ \geqslant 0} \\ 
   &  = \Re \left\langle {Af,g} \right\rangle  \\ 
   &  \leqslant \frac{\alpha }{2}\left[ {{{\left\| f \right\|}^2} + {{\left\| g \right\|}^2}} \right] \\ 
   &  = \frac{\alpha }{2}\left[ {{{\left\| {\lambda x} \right\|}^2} + {{\left\| {\frac{{Ax}}{\lambda }} \right\|}^2}} \right]. \\ 
   &  = \frac{\alpha }{2}\left[ {{\lambda ^2}{{\left\| x \right\|}^2} + \frac{{{{\left\| {Ax} \right\|}^2}}}{{{\lambda ^2}}}} \right] \\ 
\end{aligned} \] Válasszuk:
\(\lambda^{2}: = \frac{\left\| {Ax} \right\|}{\left\| x \right\|}\) ,
ekkor \(\lambda > 0\) teljesül (feltéve, hogy \(Ax \neq 0\) ), és
\[\begin{aligned}
  {\left\| {Ax} \right\|^2} &  \leqslant \frac{\alpha }{2}\left[ {\frac{{\left\| {Ax} \right\|}}{{\left\| x \right\|}}{{\left\| x \right\|}^2} + \frac{{\left\| x \right\|}}{{\left\| {Ax} \right\|}}{{\left\| {Ax} \right\|}^2}} \right] \\ 
   &  = \frac{\alpha }{2}\left[ {\left\| {Ax} \right\| \cdot \left\| x \right\| + \left\| x \right\| \cdot \left\| {Ax} \right\|} \right] \\ 
   &  = \alpha \left\| {Ax} \right\| \cdot \left\| x \right\| .\\ 
\end{aligned} \] \(\left\| {Ax} \right\| = 0\) triviális esetet kivéve
osztva \(\left\| {Ax} \right\| > 0\) -val:
\(\left\| {Ax} \right\| \leq \alpha \cdot \left\| x \right\|\) . Ez igaz
\(\left\| {Ax} \right\| = 0\) esetén is persze. Tehát
\(\left\| A \right\| \leq \alpha\) . Előbb azt kaptuk, hogy
\(\alpha \leq \left\| A \right\|\) , így a mostanival együtt:
\(\left\| A \right\| = \alpha\) .

\end{bizonyitas}

\begin{allitas}

Állítás:\\
Vezessük be \[\begin{aligned}
  M: =  & \sup \left\{ {\left\langle {Ax,x} \right\rangle :x \in X,\left\| x \right\| = 1} \right\} \\ 
  m: =  & \inf \left\{ {\left\langle {Ax,x} \right\rangle :x \in X,\left\| x \right\| = 1} \right\} \\ 
\end{aligned} \] jelöléseket! (Ekkor a fentiek miatt
\(\left\lbrack {m,M} \right\rbrack \subset \left\lbrack {- \left\| A \right\|,\left\| A \right\|} \right\rbrack\),
és \(\max\left\{ {\left| m \right|,M} \right\} = \left\| A \right\|\)).
Az \(A\) önadjungált korlátos operátor spektruma
\(\subset \left\lbrack {m,M} \right\rbrack\), más szóval, ha
\(\lambda \in {\mathbb{K}}\)-ra
\(\left. \lambda \notin \left\lbrack {m,M} \right\rbrack\Rightarrow\lambda \right.\)
reguláris érték \(A\)-ra.

\end{allitas}

\begin{megjegyzes}

Megjegyzés:\\
Azt eddig is tudtuk, hogy
\(\left| \lambda \right| > \left\| A \right\|\) esetén \(\lambda\)
reguláris érték (ha \(A\) korlátos). Azt is tudtuk, hogy ha \(A\)
szimmetrikus és \(\left. \Im\lambda \neq 0\Rightarrow\lambda \right.\)
nem lehet sajátérték.

\end{megjegyzes}

\begin{definicio}

Definíció:\\
Legyen \(\left. A:D_{A}\rightarrow X \right.\) lineáris operátor,
\(D_{A} \subset X\), \(\overline{{\mathbb{D}}_{A}} = X\). Ha
\(\left\langle {Ax,x} \right\rangle \geq 0\), \(\forall x \in D_{A}\),
akkor \(A\)-t pozitív operátornak nevezzük (konzekvensen pozitív
szemidefinitnek kéne nevezni).

\end{definicio}

\begin{allitas}

Állítás:\\
Ha \(A\) pozitív, akkor \(A\) minden sajátértéke \(\geq 0\).

\end{allitas}

\begin{bizonyitas}

Bizonyítás:\\
\(\left. Ax = \lambda x\Rightarrow 0 \leq \left\langle {Ax,x} \right\rangle = \left\langle {\lambda x,x} \right\rangle = \lambda\left\| x \right\|^{2}\Rightarrow\lambda \geq 0 \right.\),
ha \(\left\| x \right\|^{2} \neq 0\).

\end{bizonyitas}

\hypertarget{izometrikus-es-uniter-operatorok}{%
\subsubsection{Izometrikus és unitér
operátorok}\label{izometrikus-es-uniter-operatorok}}

\begin{definicio}

Definíció:\\
Legyen \(X\) Hilbert tér! Az \(\left. A:X\rightarrow X \right.\)
operátort izometrikusnak nevezzük, ha
\(\left\| {Ax} \right\| = \left\| x \right\|\), \(\forall x \in X\).
Ekkor látható, hogy \(A\) korlátos és \(\left\| A \right\| = 1\).

\end{definicio}

\begin{allitas}

Állítás:\\
Ha \(A\) izometrikus, akkor távolság és skalárszorzattartó (szögtartó).

\end{allitas}

\begin{bizonyitas}

Bizonyítás:

\begin{itemize}
\tightlist
\item
  \(\left\| {Ax - Ay} \right\| = \left\| {A\left( {x - y} \right)} \right\| = \left\| {x - y} \right\|\).
\item
  Belátjuk a skalárszorzattartást valós \(X\) Hilbert tér esetén.
  \[\begin{aligned}
    I:\quad  & {\left\| {x + y} \right\|^2} = {\left\| x \right\|^2} + 2\left\langle {x,y} \right\rangle  + {\left\| y \right\|^2} \\ 
    II:\quad  & {\left\| {x - y} \right\|^2} = {\left\| x \right\|^2} - 2\left\langle {x,y} \right\rangle  + {\left\| y \right\|^2} \\ 
    I - II:\quad  & {\left\| {x + y} \right\|^2} - {\left\| {x - y} \right\|^2} = 4\left\langle {x,y} \right\rangle  \Rightarrow \left\langle {x,y} \right\rangle  = \frac{{{{\left\| {x + y} \right\|}^2} - {{\left\| {x - y} \right\|}^2}}}{4} \\ 
  \end{aligned} \] Így \[\begin{aligned}
    \left\langle {Ax,Ay} \right\rangle  &  = \frac{1}{4}\left( {{{\left\| {Ax + Ay} \right\|}^2} - {{\left\| {Ax - Ay} \right\|}^2}} \right) \\ 
     &  = \frac{1}{4}\left( {{{\left\| {A\left( {x + y} \right)} \right\|}^2} - {{\left\| {A\left( {x - y} \right)} \right\|}^2}} \right) \\ 
     &  = \frac{1}{4}\left( {{{\left\| {x + y} \right\|}^2} - {{\left\| {x - y} \right\|}^2}} \right) \Rightarrow \left\langle {x,y} \right\rangle  = \left\langle {Ax,Ay} \right\rangle . \\ 
  \end{aligned} \]
\item
  Komplex esetben
  \(\left\langle {x,y} \right\rangle = \frac{1}{4}\left\lbrack {\left\| {x + y} \right\|^{2} - \left\| {x - y} \right\|^{2} + i\left\| {x + iy} \right\|^{2} - i\left\| {x - iy} \right\|^{2}} \right\rbrack\),
  így kicsit hosszabb a bizonyítás.
\end{itemize}

\end{bizonyitas}

\textbf{Következemény}: ha \(\left. A:X\rightarrow X \right.\)
izometrikus operátor és \(\left( {x_{1},x_{2},...} \right)\) ortonormált
rendszer, akkor \(\left( {Ax_{1},Ax_{2},...} \right)\) is ortonormált
rendszer.

\textbf{Kérdés}: ha \(\left( {x_{1},x_{2},...} \right)\) teljes
ortonormált rendszer, akkor ebből következik-e, hogy
\(\left( {Ax_{1},Ax_{2},...} \right)\) is teljes ortonormált rendszer?
Általában sajnos nem.

\begin{pelda}

Példa:\\
Legyen \(X\) végtelen dimenziós, szeparábilis Hilbert tér és
\(\left( {x_{1},x_{2},...,x_{k},...} \right)\) teljes ortonormált
rendszer. Értelmezzük \(A\)-t! Egy \(x \in X\) elemet fejtsük
Fourier-sorba!
\(x = {\sum\limits_{k = 1}^{\infty}{c_{k}x_{k}}} = c_{1}x_{1} + c_{2}x_{2} + ...\),
\(Ax: = {\sum\limits_{k = 1}^{\infty}{c_{k}x_{k + 1}}} = c_{1}x_{2} + c_{2}x_{3} + ...\).
Ez egy jól definiált lineáris operátor. Tudjuk, hogy
\(\left\| {Ax} \right\|^{2} = {\sum\limits_{k = 1}^{\infty}\left| c_{k} \right|^{2}} = \left\| x \right\|^{2}\),
tehát \(A\) izometrikus. Láthatjuk, hogy így
\(\left( {Ax_{1} = x_{2},Ax_{2} = x_{3},...} \right)\) nem teljes. Az is
kiolvasható \(A\) definíciójából, hogy
\(R_{A} = \overline{\mathcal{L}\left( {x_{2},x_{3},...} \right)}\) az
\(X\)-nek valódi altere, így \(R_{A} \neq X\).

\end{pelda}

\begin{definicio}

Definíció:\\
\(\left. A:X\rightarrow X \right.\) izometrikus operátort unitérnek
nevezzük, ha \(R_{A} = X\).

\end{definicio}

\begin{tetel}

Tétel:\\
Egy \(\left. A:X\rightarrow X \right.\) korlátos operátor unitér
\(\left. \Leftrightarrow\exists A^{- 1} = A^{*} \right.\).

\end{tetel}

\begin{bizonyitas}

Bizonyítás:

\begin{itemize}
\tightlist
\item
  \(\Rightarrow\) irányban: tfh \(A\) unitér. Ekkor \(A\) korlátossága
  lévén \(A^{*}\) értelmezve van \(X\)-n, továbbá
  \(\left\| {Ax} \right\| = \left\| x \right\|\),
  \(\left. \forall x \in X\Rightarrow A \right.\) injektív
  \(\left. \Rightarrow A^{- 1} \right.\) is létezik. Belátjuk, hogy
  \(A^{*} = A^{- 1}\). Egyrészt \(D_{A^{- 1}} = R_{A} = X\), mivel \(A\)
  unitér. Ekkor \(\forall x,y \in X\) elemre
  \(\left. \left\langle {x,y} \right\rangle = \left\langle {Ax,Ay} \right\rangle = \left\langle {x,A^{*}Ay} \right\rangle\Rightarrow y = A^{*}Ay \right.\),
  \(\left. \forall y \in X\Rightarrow A^{*}A = I\Rightarrow A^{*}AA^{- 1} = A^{- 1}\Rightarrow A^{*} = A^{- 1} \right.\)
\item
  \(\Leftarrow\) irányban: tfh \(A^{*} = A^{- 1}\). Ekkor mivel
  \(\left. D_{A^{*}} = X\Rightarrow R_{A} = D_{A^{- 1}} = X \right.\),
  továbbá
  \(\left\| {Ax} \right\|^{2} = \left\langle {Ax,Ax} \right\rangle = \left\langle {x,A^{*}Ax} \right\rangle = \left\langle {x,Ix} \right\rangle = \left\| x \right\|^{2}\),
  tehát \(A\) izometrikus is.
\end{itemize}

\end{bizonyitas}

\begin{ajanlo}

\begin{ajanlofig}

\href{https://xkcd.com}{\includegraphics[width=5.20833in,height=2.82292in]{wikipedian_protester.png}}

\end{ajanlofig}

Text

\end{ajanlo}

\begin{allitas}

Állítás:\\
Ha \(A\) unitér, akkor teljes ortonormált rendszer képe szintén teljes
ortonormált rendszer.

\end{allitas}

\begin{pelda}

Példák unitér operátorokra:

\begin{enumerate}
\def\labelenumi{\arabic{enumi}.}
\tightlist
\item
  Triviális példa az identitás.
\item
  \(X: = {\mathbb{K}}^{n}\). Tudjuk, hogy egy
  \(\left. A:{\mathbb{K}}^{n}\rightarrow{\mathbb{K}}^{n} \right.\)
  lineáris korlátos operátor megadható egy \(\mathcal{A}\) négyzetes
  mátrixszal, \[\begin{aligned}
    {\mathcal{A}} =  & \left( {\begin{array}{*{20}{c}}
    {{a_{11}}}& \cdots &{{a_{1n}}} \\ 
     \vdots & \ddots & \vdots  \\ 
    {{a_{n1}}}& \cdots &{{a_{nn}}} 
  \end{array}} \right) = \left( {\begin{array}{*{20}{c}}
    {{{\mathbf{a}}_1}} \\ 
     \vdots  \\ 
    {{{\mathbf{a}}_n}} 
  \end{array}} \right), \\ 
    {{\mathcal{A}}^*} =  & \left( {{\mathbf{\bar a}}_1^T,{\mathbf{\bar a}}_2^T,...,{\mathbf{\bar a}}_n^T} \right). \\ 
  \end{aligned}\] A leképzés pontosan akkor unitér, ha
  \[{A^*} = {A^{ - 1}} \Leftrightarrow A{A^*} = I = {A^*}A \Leftrightarrow \mathcal{A}{\mathcal{A}^*} = \mathcal{I} = {\mathcal{A}^*}\mathcal{A}.\]
  \(\mathcal{A}\mathcal{A}^{*}\) elemei:
  \(\mathbf{a}_{j}\mathbf{\overline{a}}_{k}^{T} = \left\langle {a_{j},a_{k}} \right\rangle_{{\mathbb{K}}^{n}} = \delta_{jk} = \left\{ \begin{array}{l} {1\text{~ha~}j = k} \\ {0\text{~ha~}j \neq k} \\ \end{array} \right.\)\\
  A sorvektorok tehát ortonormáltak, belátható az
  \(\mathcal{A}^{*}\mathcal{A} = \mathcal{I}\) egyenletből, hogy az
  oszlopvektorok is. Az ilyen -- unitér operátorokat megadó --
  mátrixokat ortogonális mátrixoknak is nevezzük.
\item
  Fourier-operátor (Fourier-transzformáció):
  \(X: = L^{2}\left( {\mathbb{R}} \right)\) Hilbert tér! Az
  \(\mathcal{F}\) fourier operátort így értelmezzük az
  \(\phi \in L^{2}\left( {\mathbb{R}} \right) \cap L^{1}\left( {\mathbb{R}} \right) \subset L^{2}\left( {\mathbb{R}} \right)\)
  függvényeken:
  \(\left\lbrack {\mathcal{F}\left( \phi \right)} \right\rbrack\left( x \right): = \frac{1}{\sqrt{2\pi}}{\int\limits_{- \infty}^{\infty}{e^{- ixy}\phi\left( y \right)dy}}\).
  Látható, hogy ennek csak akkor van értelme, ha
  \(\phi\left( y \right)\) integrálható. Tudjuk, hogy
  \(\left| {e^{- ixy}\phi\left( y \right)} \right| = \left| {\phi\left( y \right)} \right|\),
  mert \(\left| e^{- ixy} \right| = 1\).
  \(\phi \in L^{2}\left( {\mathbb{R}} \right)\) esetén
  \(\left\lbrack {\mathcal{F}\left( \phi \right)} \right\rbrack\left( x \right) = \lim\limits_{N\rightarrow\infty}\frac{1}{\sqrt{2\pi}}{\int\limits_{- N}^{N}{e^{- ixy}\phi\left( y \right)dy}}\)
  az \(L^{2}\left( {\mathbb{R}} \right)\) normával.
\end{enumerate}

\end{pelda}

\begin{tetel}

Tétel:\\
Az
\(\left. \mathcal{F}:L^{2}\left( {\mathbb{R}} \right)\rightarrow L^{2}\left( {\mathbb{R}} \right) \right.\)
operátor unitér, \(\mathcal{F}^{- 1} = \mathcal{F}^{*}\) a következő
képlettel adható meg:
\(\left\lbrack {\mathcal{F}^{- 1}\left( \psi \right)} \right\rbrack\left( y \right) = \lim\limits_{N\rightarrow\infty}\frac{1}{\sqrt{2\pi}}{\int\limits_{- N}^{N}{e^{ixy}\psi\left( x \right)dx}}\),
ahol a limesz \(L^{2}\left( {\mathbb{R}} \right)\) norma szerinti.

\end{tetel}

\begin{bizonyitas}

Bizonyítás (vázlatos):

\begin{enumerate}
\def\labelenumi{\arabic{enumi}.}
\tightlist
\item
  Először értelmezzük \(\mathcal{F}\)-et a következő spec. alakú lépcsős
  függvényeken:
  \(\phi_{\alpha}\left( x \right): = \begin{cases} 1 & {\text{~ha~}x\ 0\text{~és~}\alpha\text{~között~van}} \\ 0 & \text{egyébként} \\ \end{cases}\).\\
  Egyszerű számolással
  \(\left( {\mathcal{F}\phi_{\alpha}} \right)\left( x \right) = \frac{1}{\sqrt{2\pi}}\frac{1 - e^{- i\alpha x}}{ix}\).
  Bevezetjük a \(\mathcal{G}\) operátort
  \(\phi \in L^{2}\left( {\mathbb{R}} \right) \cap L\left( {\mathbb{R}} \right)\)
  függvényekre:
  \(\left( {\mathcal{G}\phi} \right)\left( x \right): = \frac{1}{\sqrt{2\pi}}{\int\limits_{- \infty}^{\infty}{e^{ixy}\phi\left( y \right)dy}}\).
  Hasonlóan adódik:
  \(\left( {\mathcal{G}\phi_{\alpha}} \right)\left( x \right) = \frac{1}{\sqrt{2\pi}}\frac{e^{i\alpha x} - 1}{x}\).
  Állítás: tetszőleges \(\phi_{\alpha},\phi_{\beta}\) esetén
  \(\left\langle {\mathcal{F}\phi_{\alpha},\mathcal{F}\phi_{\beta}} \right\rangle = \left\langle {\phi_{\alpha},\phi_{\beta}} \right\rangle\),
  \(\left\langle {\mathcal{G}\phi_{\alpha},\mathcal{G}\phi_{\beta}} \right\rangle = \left\langle {\phi_{\alpha},\phi_{\beta}} \right\rangle\)
  és
  \(\left\langle {\mathcal{F}\phi_{\alpha},\phi_{\beta}} \right\rangle = \left\langle {\phi_{\alpha},\mathcal{G}\phi_{\beta}} \right\rangle\)
  is igaz.
\item
  Kiterjesztjük az állítást lépcsős függvényekre, amik láthatóan ilyen
  függvények lineárkombinációi.
\item
  A lépcsős függvények sűrűn vannak
  \(L^{2}\left( {\mathbb{R}} \right)\)-ben. Hasonló állítást kapok ezen
  lépcsős függvényekre. \(\mathcal{F}\) és \(\mathcal{G}\)-t a
  linearitás és korlátosság megtartásával egyértelműen kiterjeszthetjük
  \(L^{2}\left( {\mathbb{R}} \right)\)-re.
\item
  \(\mathcal{F}\) és \(\mathcal{G}\) képlete
  \(L^{2}\left( {\mathbb{R}} \right)\)-en megadandó.
\end{enumerate}

\end{bizonyitas}

\begin{megjegyzes}

Megjegyzés:\\
\(\mathcal{F}\) operátor \({\mathbb{R}}^{n}\)-ben:
\(\left( {\mathcal{F}\phi} \right)\left( x \right) = \frac{1}{\left( \sqrt{2\pi} \right)^{n}}{\int_{{\mathbb{R}}^{n}}{e^{- i{\langle{x,y}\rangle}}\phi\left( y \right)dy}}\),
ha
\(\phi \in L^{2}\left( {\mathbb{R}}^{n} \right) \cap L^{1}\left( {\mathbb{R}} \right)\),
ekkor \(\mathcal{F}\) unitér.

\end{megjegyzes}

\hypertarget{veges-rendu-operatorok}{%
\subsubsection{Véges rendű operátorok}\label{veges-rendu-operatorok}}

\begin{definicio}

Definíció:\\
Legyen \(X\) Hilbert tér! Egy \(\left. A:X\rightarrow X \right.\)
korlátos operátort véges rendűnek nevezünk, ha \(R_{A}\) véges
dimenziós.

\end{definicio}

\begin{pelda}

Példa:\\
Legyenek \(\phi_{1},...,\phi_{m}\) lineárisan függetlenek, akárcsak
\(\psi_{1},\psi_{2},...,\psi_{m}\), mind \(X\)-beli elemek! Az \(A\)
operátort így értelmezzük: \(\left. A:X\rightarrow X \right.\),
\(A\left( f \right): = {\sum\limits_{j = 1}^{m}{\left\langle {f,\psi_{j}} \right\rangle\phi_{j}}}\).
Látható, hogy ez véges rendű. Világos, hogy \(A\) operátor lineáris,
\(R_{A} = \mathcal{L}\left( {\phi_{1},\phi_{2},...,\phi_{m}} \right)\)
véges dimenziós. A korlátos is:
\(\left\| {Af} \right\|_{X} \leq {\sum\limits_{j = 1}^{m}\left\| {\left\langle {f,\psi_{j}} \right\rangle\phi_{j}} \right\|} = {\sum\limits_{j = 1}^{m}{\left| \left\langle {f,\psi_{j}} \right\rangle \right| \cdot \left\| \phi_{j} \right\|}}\),
melyre a Cauchy-Schwarz szerint
\(\leq {\sum\limits_{j = 1}^{m}{\left\| f \right\|_{X} \cdot \left\| \psi_{j} \right\|_{X} \cdot \left\| \phi_{j} \right\|_{X}}} = \left\| f \right\| \cdot {\sum\limits_{j = 1}^{m}{\left\| \psi_{j} \right\|_{X} \cdot \left\| \phi_{j} \right\|_{X}}}\).

\end{pelda}

\begin{allitas}

Állítás:\\
Legyen \(X\) Hilbert tér, \(\left. A:X\rightarrow X \right.\) véges
rendű operátor. Ekkor \(\exists\phi_{1},\phi_{2},...,\phi_{m} \in X\)
lineárisan függetlenek és
\(\exists\psi_{1},\psi_{2},...,\psi_{m} \in X\) lineárisan függetlenek a
fentiek szerint, és \(A\) a fenti alakú.

\end{allitas}

\begin{bizonyitas}

Bizonyítás:\\
\(R_{A}\) véges, \(m\) dimenziós lineáris altér. Legyenek
\(\phi_{1},\phi_{2},...,\phi_{m}\) lineárisan független elemek,
\(\mathcal{L}\left( {\phi_{1},\phi_{2},...,\phi_{m}} \right) = R_{A}\).
Ezek választhatók úgy, hogy ortonormáltak legyenek (a Schmidt
eljárással). Ekkor, ha \(f \in X\),
\(Af = {\sum\limits_{j = 1}^{m}{c_{j}\left( f \right)\phi_{j}}}\).Ebben
a \(c_{j}\) együtthatók egyértelműek,
\(c_{j}\left( f \right) = \left\langle {Af,\phi_{j}} \right\rangle\).
Látjuk, hogy \(c_{j}\) lineáris funkcionál, továbbá korlátos is, és
\(\left| {c_{j}\left( f \right)} \right| = \left| \left\langle {Af,\phi_{j}} \right\rangle \right| \leq \left\| {Af} \right\| \cdot \underset{= 1}{\underset{\}\ }{\left\| \phi_{j} \right\|}} \leq \left\| A \right\| \cdot \left\| f \right\|\).
Riesz-tétel segítségével
\(\left. \exists!\psi_{j} \in X:c_{j}\left( f \right) = \left\langle {f,\psi_{j}} \right\rangle\Rightarrow Af = {\sum\limits_{j = 1}^{m}{c_{j}\left( f \right)\phi_{j}}} = {\sum\limits_{j = 1}^{m}{\left\langle {f,\psi_{j}} \right\rangle\phi_{j}}} \right.\).
Nem nehéz belátni, hogy \(\psi_{1},\psi_{2},...,\psi_{m}\) is lineárisan
függetlenek.

\end{bizonyitas}

\hypertarget{a-masodfaju-egyenlet-veges-rendu-operatorokra}{%
\subsection{A másodfajú egyenlet véges rendű
operátorokra}\label{a-masodfaju-egyenlet-veges-rendu-operatorokra}}

Legyen \(X\) Hilbert tér (véges vagy végtelen dimenziós),
\(\left. A:X\rightarrow X \right.\) véges rendű operátor. Tekintsük az
\(A\) operátornak a másodfajú egyenletét:
\(\left( {\lambda I - A} \right)f = b\), ahol \(b \in X\) adott és
\(f \in X\) keresett. Ezt az előbbiek szerint így írhatjuk:
\(\lambda f - {\sum\limits_{j = 1}^{m}{\left\langle {f,\psi_{j}} \right\rangle\phi_{j}}} = b\).
Belátjuk, hogy \(\lambda \neq 0\) esetén ez az egyenlet ekvivalens egy
lineáris algebrai egyenletrendszerrel.

Az előző egyenletet jobbról \(\psi_{k}\)-val skalárisan szorozva:
\(\lambda\left\langle {f,\psi_{k}} \right\rangle - {\sum\limits_{j = 1}^{m}{\left\langle {f,\psi_{j}} \right\rangle\left\langle {\phi_{j},\psi_{k}} \right\rangle}} = \left\langle {b,\psi_{k}} \right\rangle\),
\(k \in \left\{ {1,2,...,m} \right\}\). Keressük
\(\xi_{j}: = \left\langle {f,\psi_{j}} \right\rangle\)-t, adottak
\(a_{kj}: = \left\langle {\phi_{j},\psi_{k}} \right\rangle\),
\(\beta_{k}: = \left\langle {b,\psi_{k}} \right\rangle\). Ezzel a
jelöléssel:
\(\lambda\xi_{k} - {\sum\limits_{j = 1}^{m}{a_{kj}\xi_{j}}} = \beta_{k}\),
\(k \in \left\{ {1,2,...,m} \right\}\). Ez egy lineáris egyenletrendszer
\(\xi_{k}\) együtthatókra.
\(\xi: = \left( \begin{array}{l} \xi_{1} \\  \vdots \\ \xi_{m} \\ \end{array} \right)\),
\(\beta: = \left( \begin{array}{l} \beta_{1} \\  \vdots \\ \beta_{m} \\ \end{array} \right)\),
\(\mathcal{A}: = \left( \begin{array}{lll} a_{11} & \cdots & a_{1m} \\  \vdots & \ddots & \vdots \\ a_{m1} & \cdots & a_{mm} \\ \end{array} \right)\),
így \(\left( {\lambda\mathcal{I} - \mathcal{A}} \right)\xi = \beta\). Ha
\(f\) kielégíti a másodfajú egyenletet \(\left. \Rightarrow\xi \right.\)
kielégíti a kapott lineáris algebrai egyenletrendszert \(\lambda = 0\)
esetén is!

\begin{allitas}

Állítás:\\
Legyen \(\lambda \neq 0\) és tfh \(\xi\) kielégíti a lineáris algebrai
egyenletrendszert! Ekkor
\(f: = \frac{1}{\lambda}b + \frac{1}{\lambda}{\sum\limits_{j = 1}^{m}{\xi_{j}\phi_{j}}}\)
kielégíti a véges rendű operátorra vonatkozó másodafajú egyenletet.

\end{allitas}

\begin{bizonyitas}

Bizonyítás:\\
Behelyettesítünk a másodfajú egyenletbe, s kihasználjuk, hogy \(\xi\)
kielégíti a lineáris algebrai egyenletrendszert.

\end{bizonyitas}

\begin{tetel}

Tétel:\\
Egy \(f \in X\) elem kielégíti a véges rendű opertárra vonatkozó
másodfajú egyenletet \(\lambda \neq 0\) esetén
\(\left. \Leftrightarrow\xi_{j} = \left\langle {f,\psi_{j}} \right\rangle \right.\)
képlettel értelmezett koordinátákból álló \(\xi\) kielégíti a fenti
lineáris algebrai egyenletrendszert.

\end{tetel}

Ennek alapján a véges rendű operátorokra vonatkozó másodfajú egyenlet
megoldhatóságának elmélete következik a lineáris algebrai
egyenletrendszerek megoldhatóságának elméletéből. Két eset lehetséges:

\begin{enumerate}
\def\labelenumi{\arabic{enumi}.}
\tightlist
\item
  Ha \(\lambda \neq 0\) szám az \(\mathcal{A}\) mátrixnak nem
  sajátértéke
  \(\left. \Leftrightarrow\det\left| {\lambda\mathcal{I} - \mathcal{A}} \right| \neq 0 \right.\),
  ekkor \(\left( {\lambda\mathcal{I} - \mathcal{A}} \right)\xi = \beta\)
  egyenletben \(\forall\beta \in {\mathbb{K}}^{n}\exists!\xi\) megoldás
  \(\left. \Rightarrow\exists!f \right.\) megoldás a
  \(\left( {\lambda I - A} \right)f = b\) egyenletre. Nem nehéz belátni,
  hogy \(f\) folytonosan függ \(b\)-től. Ekkor \(\lambda \neq 0\)
  reguláris érték \(A\)-ra.
\item
  Ha \(\lambda \neq 0\) az \(\mathcal{A}\) mátrixnak sajátértéke
  \(\left. \Rightarrow\lambda \right.\) az \(A\) sajátéréke, s a kétféle
  rang egyenlő. \(\lambda = 0\) végtelen rangú sajátértéke \(A\)-nak (ha
  \(X\) végtelen dimenziós).
\end{enumerate}

\begin{ajanlo}

\begin{ajanlofig}

\href{https://xkcd.com}{\includegraphics[width=5.20833in,height=2.82292in]{wikipedian_protester.png}}

\end{ajanlofig}

Text

\end{ajanlo}

\begin{allitas}

Állítás:\\
Ha \(X\) végtelen dimenziós vektortér, akkor \(\lambda = 0\) végtelen
rangú sajátértéke az operátornak.
\(A\phi: = {\sum\limits_{j = 1}^{m}{\left\langle {\phi,\psi_{j}} \right\rangle\phi_{j}}}\).
\(\lambda = 0\) sajátérték azt jelenti, hogy \(A\phi = 0\phi = 0\)
biztosan teljesül. Mivel \(\phi_{j}\)-k lineárisan függetlenek,
\(\left\langle {\phi,\psi_{j}} \right\rangle = 0\),
\(\left. \forall j \in \left\{ {1,2,...,m} \right\}\Leftrightarrow\phi\bot\mathcal{L}\left( {\psi_{1},\psi_{2},...,\psi_{m}} \right) \right.\).

\end{allitas}

Összefoglalva: legyen \(X\) végtelen dimenziós szeparábilis Hilbert tér!
Ekkor egy \(A\) véges rendű operátor spektruma csak sajátértékekből áll,
mégpedig a 0-tól különböző (véges sok) sajátérték véges rangú (ezek
megegyeznek az \(\mathcal{A}\) mátrix sajátértékeivel, s ranguk is
megegyezik), a 0 pedig végtelen rangú sajátérték. Minden más \(\lambda\)
reguláris érték.

\begin{pelda}

Példa véges rangú operátorokra (elfajult magú integrálegyenletek):

\(X: = L^{2}\left( M \right)\), ahol \(M\) mérhető halmaz.
\(\mathcal{K}\left( {x,y} \right) = {\sum\limits_{j = 1}^{m}{\phi_{j}\left( x \right)\psi_{j}\left( y \right)}}\),
ahol
\(\left. \phi_{j},\psi_{j} \in L^{2}\left( M \right)\Rightarrow\mathcal{K} \in L^{2}\left( {M \times M} \right) \right.\).
\[\begin{aligned}
  \left( {K\phi } \right)\left( x \right) &  = \int_M {{\mathcal{K}}\left( {x,y} \right)\phi \left( y \right)dy}  \\ 
   &  = \int_M {\left[ {\mathop \sum \limits_{j = 1}^m {\phi _j}\left( x \right){\psi _j}\left( y \right)} \right]\phi \left( y \right)dy}  \\ 
   &  = \mathop \sum \limits_{j = 1}^m {\phi _j}\left( x \right)\int_M {{\psi _j}\left( y \right)\phi \left( y \right)dy} . \\ 
\end{aligned} \] Röviden:
\(K\phi = {\sum\limits_{j = 1}^{m}{\phi_{j}\left\langle {\phi,\psi_{j}} \right\rangle}}\).

\end{pelda}

Az előbbiek alapján egy elfajult magú (elsőfajú) integrálegyenlet
megoldása kiszámolható egy lineáirs algebrai egyenletrendszer
megoldásával.

\hypertarget{kompakt-teljesen-folytonos-operatorok}{%
\subsubsection{Kompakt (teljesen folytonos)
operátorok}\label{kompakt-teljesen-folytonos-operatorok}}

\begin{definicio}

Definíció:\\
Egy \(M \subset Y\) halmazt feltételesen (vagy relatíve)
sorozatkompaktnak nevezünk, ha lezárása sorozatkompakt.

\end{definicio}

\begin{megjegyzes}

Megjegyzés:\\
\(M\) feltételesen sorozatkompakt, ha tetszőleges \(M\)-beli sorozatból
kiválasztható konvergens részsorozat. \({\mathbb{R}}^{n}\)-ben a
feltételesen sorozatkompakt halmazok a korlátos halmazok.

\end{megjegyzes}

\begin{definicio}

Definíció:\\
Legyenek \(X,Y\) Banach terek! Egy \(\left. A:X\rightarrow Y \right.\)
lineáris operátort teljesen folytonosnak, avagy kompaktnak nevezünk, ha
\(X\) tetszőleges korlátos halmazát feltételesen (avagy relatíve)
sorozatkompakt halmazba képezi.

\end{definicio}

\begin{megjegyzes}

Megjegyzés:\\
Ekkor \(A\) korlátos is, továbbá két kompakt operátor összege és
számszorosa is kompakt.

\end{megjegyzes}

\begin{allitas}

Állítás:\\
Egy \(\left. A:X\rightarrow Y \right.\) operátor pontosan akkor kompakt,
ha \(\forall\left( x_{k} \right)_{k \in {\mathbb{N}}},x_{k} \in X\)
korlátos sorozatra \(\left( {Ax_{k}} \right)_{k \in {\mathbb{N}}}\)-ból
kiválasztható konvergens részsorozat.

\end{allitas}

\begin{allitas}

Állítás:\\
Legyen \(X\) Hilbert tér, \(\left. A:X\rightarrow X \right.\) véges
rendű operátor. Ekkor \(A\) kompakt.

\end{allitas}

\begin{tetel}

Tétel:\\
Legyenek \(X\), \(Y\) Banach terek, \(A_{j} \in L\left( {X,Y} \right)\)
operátorok kompaktak, és
\(\left. \exists A \in L\left( {X,Y} \right):\lim\limits_{j\rightarrow\infty}A_{j} = A\Rightarrow A \right.\)
is kompakt operátor.

\end{tetel}

\begin{bizonyitas}

Bizonyítás:

Legyen \(\left( x_{k} \right)_{k \in {\mathbb{N}}}\) egy \(X\) -beli
korlátos sorozat. Bizonyítani akarjuk, hogy
\(\left( {Ax_{k}} \right)_{k \in {\mathbb{N}}}\) -nek van konvergens
részsorozata \(Y\) -ban. Tudjuk, hogy \(A \in L\left( {X,Y} \right)\) .
Mivel \(A_{1}\) kompakt, ezért az
\(\left( {A_{1}x_{k}} \right)_{k \in {\mathbb{N}}}\) sorozatból
kiválasztható \(Y\) -ban konvergens részsorozat, legyen ez
\(\left( {A_{1}x_{k1}} \right)_{k \in {\mathbb{N}}}\) !
\(\left( {A_{2}x_{k1}} \right)_{k \in {\mathbb{N}}}\) -ből kiválasztható
konvergens részsorozat, legyen ez
\(\left( {A_{2}x_{k2}} \right)_{k \in {\mathbb{N}}}\) .
\(\left( {A_{3}x_{k2}} \right)_{k \in {\mathbb{N}}}\) -ből megint
kiválasztható... \[\begin{array}{lllllll}
 & x_{1} & x_{2} & \cdots & x_{k} & \cdots & \\
A_{1} & x_{11} & x_{21} & \cdots & x_{k1} & \cdots & {\text{részsorozatra~}\left( {A_{1}x_{k1}} \right)_{k \in {\mathbb{N}}}\text{~konvergens}} \\
A_{2} & x_{12} & x_{22} & \cdots & x_{k2} & \cdots & {\text{részsorozatra~}\left( {A_{2}x_{k2}} \right)_{k \in {\mathbb{N}}}\text{~konvergens}} \\
 \vdots & \vdots & \vdots & \ddots & \vdots & \vdots & \vdots \\
A_{j} & x_{1j} & x_{2j} & \cdots & x_{kj} & \cdots & {\text{részsorozatra~}\left( {A_{j}x_{kj}} \right)_{k \in {\mathbb{N}}}\text{~konvergens}} \\
 \vdots & \vdots & \vdots & & \vdots & \vdots & \vdots \\
\end{array}\]\\
Tekintsük az \(\left( x_{kk} \right)_{k \in {\mathbb{N}}}\) átlós
sorozatot. Belátjuk, hogy
\(\left( {Ax_{kk}} \right)_{k \in {\mathbb{N}}}\) konvergens \(Y\) -ban.
\(\left( x_{kk} \right)_{k \in {\mathbb{N}}}\) az eredeti
\(\left( x_{k} \right)_{k \in {\mathbb{N}}}\) sorozatnak olyan
részsorozata, amely bármelyik sorban levő részsorozatnak a részsorozata,
bizonyos indextől kezdve.

\[\begin{aligned}
  {\left\| {A{x_{kk}} - A{x_{mm}}} \right\|_Y} =  & {\left\| {\left[ {A{x_{kk}} - {A_j}{x_{kk}}} \right] + \left[ {{A_j}{x_{kk}} - {A_j}{x_{mm}}} \right] + \left[ {{A_j}{x_{mm}} - A{x_{mm}}} \right]} \right\|_Y} \\ 
   \leqslant  & {\left\| {\left( {A - {A_j}} \right){x_{kk}}} \right\|_Y} + {\left\| {{A_j}{x_{kk}} - {A_j}{x_{mm}}} \right\|_Y} + {\left\| {\left( {{A_j} - A} \right){x_{mm}}} \right\|_Y} \\ 
   \leqslant  & {\left\| {A - {A_j}} \right\|_{L\left( {X,Y} \right)}} \cdot {\left\| {{x_{kk}}} \right\|_X} +  \\ 
   & \quad  + {\left\| {{A_j}{x_{kk}} - {A_j}{x_{mm}}} \right\|_Y} + {\left\| {{A_j} - A} \right\|_{L\left( {X,Y} \right)}} \cdot {\left\| {{x_{mm}}} \right\|_X} \\ 
\end{aligned} \]

\(\left( x_{kk} \right)_{k \in {\mathbb{N}}}\) korlátos sorozat, ehhez
\(\exists c > 0:\left\| x_{kk} \right\| \leq c\) . Legyen
\(\varepsilon > 0\) tetszőleges. Mivel
\(\lim\limits_{j\rightarrow\infty}\left\| {A_{j} - A} \right\| = 0\) ,
ezért
\(\left. \exists j_{0}:j \geq j_{0}\Rightarrow\left\| {A_{j} - A} \right\| \leq \varepsilon \right.\)
. Válasszuk pl: \(j = j_{0}\) . Mivel
\(\left( {A_{j_{0}}x_{kk}} \right)_{k \in {\mathbb{N}}}\) konvergens,
ezért
\(\left. \exists k_{0}:k,l \geq k_{0}\Rightarrow\left\| {A_{j_{0}}x_{kk} - A_{j_{0}}x_{ll}} \right\| \leq \varepsilon \right.\)
. Tehát \(k,l \geq k_{0}\) esetén
\(\left. \left\| {Ax_{kk} - Ax_{ll}} \right\|_{Y} \leq c\varepsilon + \varepsilon + c\varepsilon = \left( {2c + 1} \right)\varepsilon\Rightarrow\left( {Ax_{kk}} \right) \right.\)
Cauchy sorozat.

\end{bizonyitas}

\textbf{Következmény}: kompakt operátorok alteret képeznek
\(L\left( {X,Y} \right)\)-ban.

\begin{tetel}

Tétel (bizonyítás nélkül):\\
Legyen \(X\) szeparábilis Hilbert tér. Ha
\(\left. A:X\rightarrow X \right.\) kompakt operátor, akkor
\(\left. \exists A_{j}:X\rightarrow X \right.\) véges rendű operátorok,
hogy
\(\lim\limits_{j\rightarrow\infty}\left\| {A_{j} - A} \right\|_{L{({X,X})}} = 0\).

\end{tetel}

Összefoglalva: ha \(X\) szeparábilis Hilbert tér, akkor az
\(\left. A:X\rightarrow X \right.\) korlátos operátor kompakt
\(\Leftrightarrow\) előáll véges rendű operátorok sorozatának norma
szerinti limeszeként.

\begin{pelda}

Példa:\\
Legyen \(X = L^{2}\left( M \right)\) Hilbert tér,
\(\left. K:L^{2}\left( M \right)\rightarrow L^{2}\left( M \right) \right.\)
négyzetesen integrálható magú integráloperátor,
\(\left( {K\phi} \right)\left( x \right): = {\int_{M}{\mathcal{K}\left( {x,y} \right)\phi\left( y \right)dy}}\).
Ez a \(K\) operátor kompakt. Ennek igazolásának alapgondolata: tudjuk,
hogy \(L^{2}\left( M \right)\) szeparábilis Hilbert tér (végtelen
dimenziós). Legyenek ebben teljes ortonormált rendszerek
\(\psi_{1},\psi_{2},...\) illetve \(\phi_{1},\phi_{2},...\). Ekkor
\(\mathcal{K}\left( {x,y} \right) = {\sum\limits_{m = 1}^{\infty}\left( {\sum\limits_{j,k \leq m}{c_{jk}\phi_{j}\left( x \right)\psi_{k}\left( y \right)}} \right)}\),
\(\mathcal{K}_{N}\left( {x,y} \right) = {\sum\limits_{m = 1}^{N}{\sum\limits_{j,k \leq m}{c_{jk}\phi_{j}\left( x \right)\psi_{k}\left( y \right)}}}\),
\(\lim\limits_{N\rightarrow\infty}\left\| {\mathcal{K}_{N} - \mathcal{K}} \right\|_{L^{2}{({M \times M})}} = 0\).
\(\mathcal{K}_{N}\)-nek véges rendű operátorok felelnek meg.
\(\left. \left\| {K_{N} - K} \right\|_{L{({L^{2}{(M)},L^{2}{(M)}})}}\rightarrow 0 \right.\),
ha \(\left. N\rightarrow\infty \right.\).

\end{pelda}

\hypertarget{masodfaju-egyenlet-kompakt-operatorokra}{%
\subsubsection{Másodfajú egyenlet kompakt
operátorokra}\label{masodfaju-egyenlet-kompakt-operatorokra}}

Legyen \(X\) szeparábilis Hilbert tér, és benne egy
\(\left. A:X\rightarrow X \right.\) kompakt operátor. Tekintsük a
\(\left( {\lambda I - A} \right)f = b\) másodfajú egyenletet, melyben
\(\lambda \neq 0\) rögzített. Tudjuk, hogy \(A\) kompakt operátor
tetszőleges előírt pontossággal megközelíthatő egy \(B\) véges rendű
operátorral. \(\left. \exists A_{0}:X\rightarrow X \right.\) véges rendű
operátor, hogy
\(\left\| {A - A_{0}} \right\| < \left| \lambda \right|\).
\(\left. B_{0}: = A - A_{0}\Leftrightarrow A = A_{0} + B_{0} \right.\),
ahol \(A_{0}\) véges rendű, és
\(\left\| B_{0} \right\| < \left| \lambda \right|\) . Tehát a másodfajú
egyenlet így írható:

\[\left[ {\lambda I - \left( {{A_0} + {B_0}} \right)} \right]f = b \Leftrightarrow \left( {\lambda I - {B_0}} \right)f = b + {A_0}f.\]

\(\left. \left| \lambda \right| > \left\| B_{0} \right\|\Rightarrow\left| \lambda \right| > B_{0} \right.\)
korlátos operátor spektrálsugara \(\left. \Rightarrow\lambda \right.\)
reguláris érték \(B_{0}\) operátorra nézve \(\Rightarrow\) a legutóbbi
egyenlet ekvivalens:
\[f = {\left( {\lambda I - {B_0}} \right)^{ - 1}}\left( {b + {A_0}f} \right) = \underbrace {{{\left( {\lambda I - {B_0}} \right)}^{ - 1}}b}_{{\text{adott}}} + {\left( {\lambda I - {B_0}} \right)^{ - 1}}{A_0}f.\]
\(\lambda\)-val beszorozva, átrendezve:
\[\lambda f - \underbrace {\lambda {{\left( {\lambda I - {B_0}} \right)}^{ - 1}}{A_0}}_{: = B: = {B_\lambda }}f = \underbrace {\lambda {{\left( {\lambda I - {B_0}} \right)}^{ - 1}}}_{: = g}.\]
A bevezetett jelöléssel
\(\left( {\lambda I - B_{\lambda}} \right)f = g\) .

Észrevétel: \(B_{\lambda}\) véges rendű operátor, mert \(A_{0}\) véges
rendű operátor. Legyen \(\delta > 0\) rögtített szám, és válasszuk
\(A_{0}\)-t úgy, hogy \(\left\| {A - A_{0}} \right\| < \delta\) legyen.
Ekkor az előbbi gondolatmenet érvényes \(\forall\lambda\)-ra, \(A_{0}\)
nem függ \(\lambda\)-tól, ha \(\lambda \geq \delta\) (de \(\delta\)-tól
igen). \(A_{0}\) véges rendű operátor \(\lambda \geq \delta\) esetén, és
\(A_{0}f = {\sum\limits_{j = 1}^{m}{\left\langle {f,\psi_{j}} \right\rangle\phi_{j}}}\)
alakban írható.
\(Bf = B_{\lambda}f = \lambda\left( {\lambda I - B_{0}} \right)^{- 1}{\sum\limits_{j = 1}^{m}{\left\langle {f,\psi_{j}} \right\rangle\phi_{j}}} = {\sum\limits_{j = 1}^{m}{\lambda\left\langle {f,\psi_{j}} \right\rangle\left( {\lambda I - B_{0}} \right)^{- 1}\phi_{j}}}\).
A másodfajú egyenlet:
\(\lambda f - {\sum\limits_{j = 1}^{m}{\lambda\left\langle {f,\psi_{j}} \right\rangle\left( {\lambda I - B_{0}} \right)^{- 1}\phi_{j}}} = g = g_{\lambda}\).

\begin{ajanlo}

\begin{ajanlofig}

\href{https://xkcd.com}{\includegraphics[width=5.20833in,height=2.82292in]{wikipedian_protester.png}}

\end{ajanlofig}

Text

\end{ajanlo}

Tehát kaptuk, hogy
\(\lambda f - {\sum\limits_{j = 1}^{m}{\lambda\left\langle {f,\psi_{j}} \right\rangle\left( {\lambda I - B_{0}} \right)^{- 1}\phi_{j}}} = g = g_{\lambda}\).
Ez megfelel egy lineáris algebrai egyenletrendszernek:
\(\lambda\mathcal{I}\xi - \mathcal{B}_{\lambda}\xi = \beta_{\lambda}\).
Ekkor
\(\det\left( {\lambda\mathcal{I} - \mathcal{B}_{\lambda}} \right) = 0\)
egyenlet gyökei a sajátértékek. A mátrix (\(\mathcal{B}_{\lambda}\)) és
az operátor (\(B_{\lambda}\)) sajátértékei azonosak az eredeti operátor
(\(A\)) sajátértékeivel, és rangjuk is azonos. Belátható, hogy a mátrix
elemei a \(\lambda\) változónak holomorf függvényei! Így a determináns
is holomorf függvénye \(\lambda\)-nak. Tudjuk, hogy egy holomorf
függvény gyökei nem torlódhatnak egy véges pontban, hacsak nem az
azonosan 0 függvény. Mivel \(\lambda < \left\| A \right\|\), ezért csak
véges sok gyök van. Tehát tetszőleges rögzített \(\delta\) esetén \(A\)
operátornak véges sok \(\delta\)-nál nagyobb abszolút értékű sajátértéke
van, s ezek véges rangúak.

\begin{tetel}

Tétel:\\
Ha \(A\) kompakt operátor, akkor \(A\)-nak legfeljebb megszámlálhatóan
végtelen sok sajátértéke van, a 0-tól különböző sajátértékek véges
rangúak, s a sajátértékek csak a 0-ban torlódhatnak. (Gondoljunk csak a
\(\delta: = 1/k\), \(k \in {\mathbb{N}}\) esetre!)

\end{tetel}

\begin{tetel}

Tétel (biz. nélkül):\\
minden \(\lambda \neq 0\), ami nem sajátérték, az reguláris érték \(A\)
(kompakt operátorra) nézve.

\end{tetel}

Következmény: ha \(\lambda \neq 0\) nem sajátérték,
\(\left( {\lambda I - A} \right)f = b\) másodfajú egyenletnek
\(\forall b\)-re létezik egyetlen \(f\) megoldás, és ez folytonosan függ
\(b\)-től.

Mi a helyzet, ha \(\lambda\) sajátérték?

Emlékeztető: tetszőleges korlátos lineáris operátor esetén
\[\left. {\overline{R_{\lambda I-A}}}^{\bot} = S_{\overline{\lambda}}\left( A^{*} \right)\Leftrightarrow\overline{R_{\lambda I-A}} = S_{\overline{\lambda}}\left( A^{*} \right)^{\bot}. \right.\]Ha
\(R_{\lambda I - A}\) zárt altér, akkor
\(R_{\lambda I - A} = \overline{R_{\lambda I-A}} = S_{\overline{\lambda}}\left( A^{*} \right)^{\bot}\).

\begin{tetel}

Tétel:\\
Ha \(A\) kompakt operátor, akkor \(\lambda \neq 0\) esetén
\(R_{\lambda I - A}\) zárt altér.

\end{tetel}

\begin{bizonyitas}

Bizonyítás:\\
Látható, hogy \(R_{\lambda I - A}\) lineáris altér. Azt kell
bizonyítani, hogy \(R_{\lambda I - A}\) zárt halmaz. Legyen tetszőleges
\(\psi_{j} \in R_{\lambda I - A}\) és \(\exists\lim\psi_{j} = \psi\),
ekkor \(\psi \in R_{\lambda I - A}\)? Mivel
\(\left. \psi_{j} \in R_{\lambda I - A}\Rightarrow\exists\phi_{j} \in X:\left( {\lambda I - A} \right)\phi_{j} = \psi_{j} \right.\).
Jelöljük:
\[S_{\lambda}\left( A \right): = \left\{ {\phi \in X:\left( {\lambda I - A} \right) = 0} \right\}.\]
Ekkor \(S_{\lambda}\left( A \right)\) zárt lineáris altér (\(A\)
folytonos). A Riesz-tétel következtében
\[\left. X: = S_{\lambda}\left( A \right) \oplus S_{\lambda}\left( A \right)^{\bot}\Leftrightarrow\forall x \in X\exists!x_{1},x_{2}:x_{1} \in S_{\lambda}\left( A \right),x_{2} \in S_{\lambda}\left( A \right)^{\bot}, \right.\]
ahol \(x = x_{1} + x_{2}\). Ennek megfelelően
\(X \ni \phi_{j} = f_{j} + g_{j}\), ahol
\(f_{j} \in S_{\lambda}\left( A \right)\).
\(g_{j} \in S_{\lambda}\left( A \right)^{\bot}\),
\[{\psi _j} = \left( {\lambda I - A} \right){\phi _j} = \underbrace {\left( {\lambda I - A} \right){f_j}}_0 + \left( {\lambda I - A} \right){g_j} \Rightarrow \left( {\lambda I - A} \right){g_j} = {\psi _j}.\]

\begin{allitas}

Kis állítás (a bizonyításon belül)::\\
\(\left( g_{j} \right)_{j \in {\mathbb{N}}}\) korlátos sorozat
\(X\)-ben.

\end{allitas}

\begin{bizonyitas}

Bizonyítás (a kis állításé):\\
Indirekt feltesszük, hogy
\(\exists\left( g_{j_{k}} \right)_{k \in {\mathbb{N}}}\) részsorozat,
hogy
\(\lim\limits_{k\rightarrow\infty}\left\| g_{j_{k}} \right\|_{X} = \infty\).
Legyen \(h_{j_{k}} = \frac{g_{j_{k}}}{\left\| g_{j_{k}} \right\|_{X}}\),
ekkor \(\left\| h_{j_{k}} \right\|_{X} = 1\).
\(\left( {\lambda I - A} \right)g_{j_{k}} = \psi_{j_{k}}\) egyenletet
osztva \(\left\| g_{j_{k}} \right\|\)-val:
\(\left. \left( {\lambda I - A} \right)h_{j_{k}} = \frac{\psi_{j_{k}}}{\left\| g_{j_{k}} \right\|_{X}}\rightarrow 0_{X} \right.\),
ugyanis \(\psi_{j}\) konvergens \(\Rightarrow\) korlátos.
\(\lim\limits_{k\rightarrow\infty}\left( {\lambda h_{j_{k}} - Ah_{j_{k}}} \right) = 0_{X}\).
\(\left( h_{j_{k}} \right)\) korlátos sorozat (mert
\(\left\| h_{j_{k}} \right\| = 1\)), \(A\) kompakt operátor, ezért
\(\exists\left( {\widetilde{h}}_{j_{k}} \right)\) részsorozat, amelyre
\(\left( {A{\widetilde{h}}_{j_{k}}} \right)\) konvergens
\(\left. \Leftrightarrow\left( {\lambda{\widetilde{h}}_{j_{k}}} \right)_{k \in {\mathbb{N}}} \right.\)
is konvergens.
\(\left. \lambda \neq 0\Rightarrow\left( {\widetilde{h}}_{j_{k}} \right) \right.\)
konvergens,
\(\left. \left( {\widetilde{h}}_{j_{k}} \right)_{k \in {\mathbb{N}}}\rightarrow h_{0}\Rightarrow\left( {\lambda I - A} \right){\widetilde{h}}_{j_{k}}\rightarrow 0\Rightarrow\left( {\lambda I - A} \right)h_{0} = 0 \right.\).
Ebből következik, hogy \(h_{0} \in S_{\lambda}\left( A \right)\).
Másrészt \(h_{j_{k}} = \frac{g_{jk}}{\left\| g_{jk} \right\|}\),
\(\left. g_{j_{k}} \in S_{\lambda}\left( A \right)^{\bot}\Rightarrow h_{j_{k}} \in S_{\lambda}\left( A \right)^{\bot}\Rightarrow \right.\)
limeszben \(h_{0} \in S_{\lambda}\left( A \right)^{\bot}\). Másrészt
\(h_{0} \in S_{\lambda}\left( A \right)\), így \(h_{0} = 0\), de ez meg
nem lehet, mert
\(\left. \left\| {\widetilde{h}}_{j_{k}} \right\| = 1\Rightarrow\left\| h_{0} \right\| = 1 \right.\)
kéne lennie.

\end{bizonyitas}

Tehát \(\left( {\lambda I - A} \right)g_{j} = \psi_{j}\),
\(\lim\left( \psi_{j} \right) = \psi\), \(\left\| g_{j} \right\|_{X}\)
korlátos. Mivel \(A\) kompakt és \(g_{j}\) korlátos
\(\left. \Rightarrow\exists{\widetilde{g}}_{j_{k}} \right.\)
részsorozat, hogy \(A{\widetilde{g}}_{j_{k}}\) konvergens.
\(\psi_{j_{k}}\) is konvergens
\(\left. \Rightarrow\lambda g_{j_{k}} \right.\) is konvergens,
\(\left. \lambda \neq 0\Rightarrow\left( g_{j_{k}} \right) \right.\)
konvergens. \(\left. g_{j_{k}}\rightarrow g_{0} \right.\) \(X\)-ben,
\(g_{0} \in X\).
\(\left. \left( {\lambda I - A} \right)g_{0} = \psi\Rightarrow\psi \in R_{\lambda I - A} \right.\).

\end{bizonyitas}

\begin{tetel}

Tétel\\
(bizonyítás nélkül): legyen \(\left. A:X\rightarrow X \right.\) kompakt
operátor. Ekkor \(A^{*}\) is kompakt. Továbbá \(\lambda \neq 0\) az
\(A\)-nak sajátértéke
\(\left. \Leftrightarrow\overline{\lambda} \right.\) sajátértéke
\(A^{*}\)-nak, és ekkor a rangok egyenlők.

\end{tetel}

Összefoglalás (Fredholm alternatíva): legyen
\(\left. A:X\rightarrow X \right.\) kompakt operátor, \(\lambda \neq 0\)
tetszőleges szám s \(\left( {\lambda I - A} \right)f = b\) másodfajú
egyenlet. Ekkor két eset lehetséges:

\begin{enumerate}
\def\labelenumi{\arabic{enumi}.}
\tightlist
\item
  Ha \(\lambda \neq 0\) az \(A\)-nak nem sajátértéke (legfeljebb
  megszámlálhatóan végtelen sok, véges rangú, 0-ban torlódó
  sajátértékek), akkor a másodfajú egyenletnek \(\forall b \in X\)
  esetén \(\exists!f\) megoldása és ez folytonosan függ \(b\)-től
  (\(\left( {\lambda I - A} \right)^{- 1}\) folytonos)
\item
  Ha \(\lambda \neq 0\) sajátérték, akkor a másodfajú egyenletnek a
  megoldása nem egyértelmű, a homogén egyenletnek véges sok lineárisan
  független megoldása van. A megoldás pontosan létezik, ha
  \(b\bot S_{\overline{\lambda}}\left( A^{*} \right)\) minden elemére.
  Ez annyi db ortogonalitási feltétel, amennyi a \(\lambda\) sajátérték
  rangja.
\end{enumerate}

\hypertarget{onadjungalt-kompakt-operatorok}{%
\subsubsection{Önadjungált kompakt
operátorok}\label{onadjungalt-kompakt-operatorok}}

\begin{tetel}

Tétel:\\
Legyen \(X\) szeparábilis Hilbert tér,
\(\left. A:X\rightarrow X \right.\) kompakt és önadjungált operátor,
\(A \neq 0\). Ekkor \(\exists\lambda_{1}\) sajátérték:
\(\left| \lambda_{1} \right| = \left\| A \right\| = \sup\left\{ {\left| \left\langle {Ax,x} \right\rangle \right|:\left\| x \right\|_{X} = 1} \right\}\).

\end{tetel}

\begin{megjegyzes}

Megjegyzés:\\
Ha \(\lambda_{1}\) az \(A\) operátor olyan sajátértéke, amelyre
\(\left| \lambda_{1} \right| = \left\| A \right\|\) és \(x_{1}\) olyan
sajátelem, hogy \(\left\| x_{1} \right\| = 1\), azaz
\(Ax_{1} = \lambda_{1}x_{1}\), \(\left\| x_{1} \right\| = 1\), akkor
\(\left| \left\langle {Ax_{1},x_{1}} \right\rangle \right| = \left| \left\langle {\lambda_{1}x_{1},x_{1}} \right\rangle \right| = \left| {\lambda_{1}\left\langle {x_{1},x_{1}} \right\rangle} \right| = \left| \lambda_{1} \right| = \left\| A \right\| = \sup\left\{ \left| {\left\langle {Ax,x} \right\rangle:\left\| x \right\|_{X} = 1} \right| \right\}\).
Más szóval, az
\(\left. x\mapsto\left| \left\langle {Ax,x} \right\rangle \right| \right.\),
ahol \(\left\| x \right\| = 1\), ez a függvény felveszi a suprémumot az
\(x = x_{1}\) sajátelemen, a maximum (ami most a suprémum is) értéke
\(= \left| \lambda_{1} \right|\). Fordítva: ha \(x^{*}\) olyan, hogy
\(\left\| x^{*} \right\| = 1\), és arra
\(\left| \left\langle {Ax,x} \right\rangle \right|\) maximális, akkor ez
sajátelem és a maximum egyenlő a sajátérték abszolút értékével. Ugyanis
\(\left| \left\langle {Ax^{*},x^{*}} \right\rangle \right| \leq \left\| {Ax^{*}} \right\| \cdot \left\| x^{*} \right\| \leq \left\| A \right\| \cdot \left\| x^{*} \right\|^{2} = \left\| A \right\|\),
a Cauchy-Schwarz egyenlőtlenségben egyenlőség pontosan akkor áll fenn,
amikor \(Ax^{*} \parallel x^{*}\), azaz \(Ax^{*} = const \cdot x^{*}\).

\end{megjegyzes}

További sajátértékek, sajátelemek keresése.

Legyen \(X_{1}: = \left\{ {x \in X:x\bot x_{1}} \right\}\), ahol
\(A_{1}: = \left. A \right|_{X_{1}}\), a leszűkítés, és
\(Ax_{1} = \lambda_{1}x_{1}\).

\begin{allitas}

Állítás:\\
\(X_{1}\) invariáns altér, azaz
\(\left. x \in X_{1}\Rightarrow Ax \in X_{1} \right.\).

\end{allitas}

\begin{bizonyitas}

Bizonyítás:\\
Tfh \(x \in X_{1}\)!
\(\left\langle {Ax,x_{1}} \right\rangle = \left\langle {x,Ax_{1}} \right\rangle = \left\langle {x,\lambda_{1}x_{1}} \right\rangle = \lambda_{1}\left\langle {x,x_{1}} \right\rangle = 0\),
tehát \(Ax \in X_{1}\). Az előbbi tételt alkalmazhatjuk az \(A_{1}\)
operátorra \(X_{1}\) Hilbert térben. Ekkor \(\exists\lambda_{2}\)
sajátérték, hogy
\(\left| \lambda_{2} \right| = \left\| A_{1} \right\| = \sup\left\{ {\left\langle {A_{1}x,x} \right\rangle:\left\| x \right\|_{X} = 1,x \in X_{1}} \right\}\).
A maximum helye \(x_{2}\) sajátelem helyén van,
\(\lambda_{2}x_{2} = Ax_{2}\), \(x_{2}\bot x_{1}\). Így egymás után
megkaphatjuk az \(A\) operátor sajátértékeit és sajátelemeit,
\(\left| \lambda_{1} \right| \geq \left| \lambda_{2} \right| \geq ...\).
Ha \(A\) véges rendű, akkor az eljárás véges sok lépés után befejeződik.

\end{bizonyitas}

\begin{tetel}

Tétel:\\
Legyenek az \(A\) önadjungált operátor sajátértékei \(\lambda_{1}\),
\(\lambda_{2}\),\ldots{} és sajátelemei \(x_{1}\), \(x_{2}\),\ldots{} A
sajátelemekről feltehető, hogy ortonormált rendszert alkotnak. Ekkor
\(\forall x \in X\) elemre
\(Ax = {\sum\limits_{k}{\lambda_{k}\left\langle {x,x_{k}} \right\rangle x_{k}}}\).
Az \(\left( x_{k} \right)\) ortonormált rendszert kibővítve a
\(\lambda = 0\)-hoz tartozó sajátelemek ortonormált rendszerével, akkor
ezek egy teljes ortonormált rendszert alkotnak.

\end{tetel}

\end{document}
